% Generated mechanically from frozen input p07/ja/resolve.tex.
% Do not edit this generated file; edit the bound locale source instead.

% OK, start here.
%

\setcounter{chapter}{53}
\chapter{曲面の特異点解消}



\phantomsection
\label{resolve-section-phantom}


\section{序論}
\label{resolve-section-introduction}

\noindent
本章では，Lipman \cite{Lipman} に従い，また主として
\cite{Artin-Lipman} における Artin の解説に従って，曲面の特異点解消を論じる．
主要結果（定理 \ref{resolve-theorem-resolve}）によれば，$2$ 次元 Noether 概形 $Y$ は，
有限な正規化 $Y^\nu \to Y$ をもち，その特異点 $y_i \in Y^\nu$ が有限個で，
各 $i$ に対して完備化 $\mathcal{O}_{Y^\nu, y_i}^\wedge$ が正規であれば，
特異点解消をもつ．

\medskip\noindent
実際，$Y$ が $2$ 次元概形で，準優秀基礎環 $R$ 上有限型であるとする
（例えば，体上，または分数体の標数が $0$ である Dedekind 整域，
例えば $\mathbf{Z}$ 上であるとする）．このとき $Y$ の正規化は有限であり，
その特異点は有限個で，局所環の完備化は正規である．
More on Algebra の各節
\ref{more-algebra-section-singular-locus},
\ref{more-algebra-section-G-ring}，および
\ref{more-algebra-section-excellent}
における議論，および More on Algebra，補題
\ref{more-algebra-lemma-normal-goes-up} を参照されたい．
したがって，このような $Y$ は特異点解消をもつ．

\medskip\noindent
証明の概略は次の通りである．$A$ を次元 $2$ の Noether 局所整域とする．
証明は以下の段階からなる．
\begin{enumerate}
\item[N] $A$ をその正規化で置き換える；
\item[V] Grauert--Riemenschneider の定理を証明する；
\item[B] 最大値 $g$ が存在し，それが $H^1(X, \mathcal{O}_X)$ の長さを
すべての正規改変 $X \to \Spec(A)$ にわたってとった最大値であることを示し，
$g = 0$ の場合に帰着する；
\item[R] $A$ は $g = 0$ のとき有理特異点を定めるという．この場合，
有限回のブローアップの後で $A$ が Gorenstein かつ $g = 0$ であると仮定できる；
\item[D] $A$ は，$g = 0$ かつ $A$ が Gorenstein であるとき，有理二重点を
定めるという．この場合には特異点を明示的に解消する．
\end{enumerate}
これらの各段階では，環 $A$ に仮定が必要となる．以下，順に論じる．

\medskip\noindent
N について：ここでは $A$ の正規化が有限であると仮定する必要がある
（これは自動的には成り立たない）．本章の大部分では，ある正規化が有限であることを
必要とするとき，対象の概形が Nagata であると仮定する．しかし Nagata であるという条件は，
定理の仮定よりもわずかに強い．この小さな問題を避ける一つの方法は，
% P07-ERR-1094
環 $A$ が形式的不分岐である（すなわち，その完備化が被約である）ことと，
補題 \ref{resolve-lemma-formally-unramified} を用いることであった．しかし本章の証明では，
補題 \ref{resolve-lemma-normalization-completion} を用い，完備化上での正規化の有限性から
$A$ 上での正規化の有限性を導く方が容易である．

\medskip\noindent
V について：これは命題 \ref{resolve-proposition-Grauert-Riemenschneider} であり，
大まかには，正規改変 $f : X \to \Spec(A)$ に対して
$R^1f_*\omega_X = 0$ が成り立つと主張する．ここで $\omega_X$ は $X/A$ の
双対化加群である（注意 \ref{resolve-remark-dualizing-setup}）．実際，双対性により，
この結果は対象 $Rf_*\mathcal{O}_X$ に関する導来圏 $D(A)$ での主張
（補題 \ref{resolve-lemma-R1-injective}）と同値である．その一方で，証明には，
特殊ファイバーの各成分が正の余法層をもつという標準的事実
（補題 \ref{resolve-lemma-nontrivial-normal-bundle}）を用いる．

\medskip\noindent
B について：これは，ある意味で証明中もっとも微妙な部分である．記述自体はやや一般的だが，
最終的にこの段階の結論を必要とするのは，$A$ が完備 Noether 局所環である場合だけである．
用語は定義 \ref{resolve-definition-reduce-to-rational} で定める．上で定義した $g$ が有界ならば，
直接的な議論により，正規改変 $X \to \Spec(A)$ であって，$X$ のすべての
特異点が有理特異点となるものを見いだせる（補題 \ref{resolve-lemma-reduce-to-rational}）．
有限拡大 $A \subset B$ が与えられたとき，次の二つの場合には，$g$ が $B$ について
有界であることが，$A$ についての有界性から従うことを示す：
% P07-ERR-1095
(1) 分数体の拡大が分離的な場合（補題 \ref{resolve-lemma-reduce-to-rational} を参照）；
(2) 分数体の拡大次数が $p$，標数が $p$ で，$A$ が正則かつ完備である場合
（補題 \ref{resolve-lemma-go-up-degree-p} を参照）．

\medskip\noindent
R について：ここでは $g = 0$ の場合を Gorenstein の場合に帰着する．
議論全体を可能にする素晴らしい事実は，有理曲面特異点のブローアップが正規であることである
（補題 \ref{resolve-lemma-blow-up-normal-rational} を参照）．

\medskip\noindent
D について：有理二重点の解消は，ほぼ直接的な計算によって進める．
節 \ref{resolve-section-rational-double-points} を参照されたい．有理二重点は超曲面特異点である
（これは正しいが，ここでは必要ないため証明しない）．その局所方程式は次の形をしている．
$$
a_{11} x_1^2 + a_{12} x_1x_2 + a_{13}x_1x_3 + a_{22} x_2^2 +
a_{23} x_2x_3 + a_{33} x_3^2 =
\sum a_{ijk} x_ix_jx_k
$$
重複度が $2$ であり，任意のブローアップの後にも $2$ のままであるため二次部分は
零ではありえないこと，および各ブローアップが正規であることを用いると，速やかに
解消が得られる．ただし，ブローアップのたびに減少する不変量は存在しないので，
この過程が停止することを示すために，節 \ref{resolve-section-arcs} の形式弧に関する議論を用いる．

\medskip\noindent
以上を組み合わせるには，さらにいくらかの作業が必要である．主な難点は，完備化
$A^\wedge$ の解消を $\Spec(A)$ の解消へ持ち上げることにある．そのためまず，
解消が存在すれば，正規化ブローアップによる解消が存在することを示す
（補題 \ref{resolve-lemma-existence-implies-existence-by-normalized-blowing-ups}）．
補題 \ref{resolve-lemma-normalized-blowup-completion} により，正規化ブローアップの列は
完備化から持ち上げることができる．この事実は，完備局所環 $A$ の解消を証明する際にも用いる．
というのも，本章の方法は，有限包含 $A_0 \subset A$ の次数に関する帰納法によるからであり，
ここで $A_0$ は正則である（補題 \ref{resolve-lemma-resolve-complete} を参照）．
B の段階でより強い結果（例えば Lipman の論文で証明されている結果）を用いれば，
この段階は不要となる．




\section{正標数におけるトレース写像}
\label{resolve-section-trace}

\noindent
本節の結果の一部は，de Rham Cohomology，節 \ref{derham-section-trace} における
微分形式上のトレースに関する，はるかに一般的な議論から導ける．
この点については注意 \ref{resolve-remark-compare-Garel} を参照されたい．

\medskip\noindent
素数 $p$ を固定する．$R$ を $\mathbf{F}_p$-代数とする．
$a \in R$ が与えられたとき，$S = R[x]/(x^p - a)$ とおく．$R$-線形写像
$$
\text{Tr}_x : \Omega_{S/R} \longrightarrow \Omega_R
$$
を次の規則で定める．
$$
x^i\text{d}x \longmapsto
\left\{
\begin{matrix}
0 & \text{if} & 0 \leq i \leq p - 2, \\
\text{d}a & \text{if} & i = p - 1
\end{matrix}
\right.
$$
これは意味をもつ．実際，$\Omega_{S/R}$ は $R$-加群として自由であり，
$x^i\text{d}x$（$0 \leq i \leq p - 1$）を基底にもつ．
次の補題より，このトレース写像は良定義，すなわち座標 $x$ の選び方に依存しない．

\begin{lemma}
\label{resolve-lemma-trace-well-defined}
$\varphi : R[x]/(x^p - a) \to R[y]/(y^p - b)$ を $R$-代数準同型とする．
このとき $\text{Tr}_x = \text{Tr}_y \circ \varphi$ である．
\end{lemma}

\begin{proof}
$\varphi(x) = \lambda_0 + \lambda_1 y + \ldots + \lambda_{p - 1}y^{p - 1}$
と書き，$\lambda_i \in R$ とする．$x$ を
$\lambda_0 + \lambda_1 y + \ldots + \lambda_{p - 1}y^{p - 1}$
へ送る写像が $R$-代数準同型 $R[x]/(x^p - a) \to R[y]/(y^p - b)$ を
誘導するための条件は，
$$
a = \lambda_0^p + \lambda_1^p b + \ldots + \lambda_{p - 1}^pb^{p - 1}
$$
が環 $R$ で成り立つことと同値である．多項式環
$$
R_{univ} = \mathbf{F}_p[b, \lambda_0, \ldots, \lambda_{p - 1}]
$$
と，その元
$a = \lambda_0^p + \lambda_1^p b + \ldots + \lambda_{p - 1}^pb^{p - 1}$
% P07-ERR-1096
を考える．さらに，普遍代数写像
$\varphi_{univ} : R_{univ}[x]/(x^p - a) \to R_{univ}[y]/(y^p - b)$
であって，$x$ を
$\lambda_0 + \lambda_1 y + \ldots + \lambda_{p - 1}y^{p - 1}$
へ送る．標準写像
$$
R_{univ} \longrightarrow R
$$
で $b, \lambda_i$ をそれぞれ $b, \lambda_i$ へ送るものを得る．構成により，
交換図式
$$
\xymatrix{
R_{univ}[x]/(x^p - a) \ar[r] \ar[d]_{\varphi_{univ}} &
R[x]/(x^p - a) \ar[d]^\varphi \\
R_{univ}[y]/(y^p - b) \ar[r] & R[y]/(y^p - b)
}
$$
を得て，水平射はトレース写像と両立する．したがって，写像
$\varphi_{univ}$ に対して補題を証明すれば十分である．よって
$R = \mathbf{F}_p[b, \lambda_0, \ldots, \lambda_{p - 1}]$ が多項式環であると
仮定してよい．この場合，$\text{Tr}_y(\varphi(x)^i\text{d}\varphi(x))$ を
$i = 0 , \ldots, p - 1$ について計算して補題を確認する．

\medskip\noindent
$0 \leq i \leq p - 2$ の場合を考える．
$$
(\lambda_0 + \lambda_1 y + \ldots + \lambda_{p - 1}y^{p - 1})^i
(\lambda_1 + 2 \lambda_2 y + \ldots + (p - 1)\lambda_{p - 1}y^{p - 2})
$$
を環 $R[y]/(y^p - b)$ で展開する．$y^{p - 1}$ の係数が零であることを
示さなければならない．そのためには，上の式を $y$ の多項式とみたとき，
$y^{pk - 1}$ の係数がすべて零であることを示せば十分である．
この多項式を
$$
\frac{\text{d}}{(i + 1)\text{d}y}
(\lambda_0 + \lambda_1 y + \ldots + \lambda_{p - 1}y^{p - 1})^{i + 1}
$$
と書けば，実際 $y^{kp - 1}$ の係数はすべて零である．

\medskip\noindent
$i = p - 1$ の場合を考える．
$$
(\lambda_0 + \lambda_1 y + \ldots + \lambda_{p - 1}y^{p - 1})^{p - 1}
(\lambda_1 + 2 \lambda_2 y + \ldots + (p - 1)\lambda_{p - 1}y^{p - 2})
$$
を環 $R[y]/(y^p - b)$ で展開する．証明を終えるには，$y^{p - 1}$ の係数と
$\text{d}b$ の積が $\text{d}a$ に等しいことを示せばよい．ここで $R$ は
$S/pS$ であり，
$S = \mathbf{Z}[b, \lambda_0, \ldots, \lambda_{p - 1}]$．
上の式を $y$ の多項式とみると，これは
$$
\frac{\text{d}}{p\text{d}y}
(\lambda_0 + \lambda_1 y + \ldots + \lambda_{p - 1}y^{p - 1})^p
$$
に等しい．$\frac{\text{d}}{\text{d}y}(y^{pk}) = pk y^{pk - 1}$ なので，
$y^{pk}$ の，次の多項式における係数を
$(\lambda_0 + \lambda_1 y + \ldots + \lambda_{p - 1}y^{p - 1})^p$
について，$p$ を法として調べれば十分である．これらの項の和は
$$
\lambda_0^p + \lambda_1^py^p + \ldots + \lambda_{p - 1}^py^{p(p - 1)}
\bmod p
$$
である．したがって，作用素 $\frac{\text{d}}{p\text{d}y}$ を適用し，
$y^p - b$ を法として簡約すると，
% P07-ERR-0381
$$
\lambda_1^p + 2\lambda_2^pb + \ldots + (p - 1)\lambda_{p - 1}b^{p - 2}
$$
を得る．これが求める $y^{p - 1}$ の係数である．ここで
$a = \lambda_0^p + \lambda_1^p b + \ldots + \lambda_{p - 1}^pb^{p - 1}$
が $R$ で成り立つことから，
$$
\text{d}a = (\lambda_1^p  + 2 \lambda_2^p b + \ldots +
(p - 1) \lambda_{p - 1}^p b^{p - 2}) \text{d}b
$$
を $R$ で得る．これで $y^{p - 1}$ の係数が所望のものであることが分かった．
\end{proof}

\begin{lemma}
\label{resolve-lemma-trace-higher}
$\mathbf{F}_p \subset \Lambda \subset R \subset S$ を環の拡大とし，
$S$ は $R[x]/(x^p - a)$ と同型であると仮定する．ここで $a \in R$ である．
このとき，標準的な $R$-線形写像
$$
\text{Tr} :
\Omega^{t + 1}_{S/\Lambda}
\longrightarrow
\Omega_{R/\Lambda}^{t + 1}
$$
であって，$t \geq 0$ に対し
$$
\eta_1 \wedge \ldots \wedge \eta_t \wedge x^i\text{d}x
\longmapsto
\left\{
\begin{matrix}
0 & \text{if} & 0 \leq i \leq p - 2, \\
\eta_1 \wedge \ldots \wedge \eta_t \wedge \text{d}a & \text{if} & i = p - 1
\end{matrix}
\right.
$$
と作用するものが存在する．ここで $\eta_i \in \Omega_{R/\Lambda}$ である．さらに
$\text{Tr}$ は，写像 $S \otimes_R \Omega_{R/\Lambda}^{t + 1} \to
\Omega_{S/\Lambda}^{t + 1}$ の像を零にする．
\end{lemma}

\begin{proof}
$t = 0$ のときは，合成
$$
\Omega_{S/\Lambda} \to \Omega_{S/R} \to \Omega_R \to \Omega_{R/\Lambda}
$$
を用いる．ここで第2の写像は補題 \ref{resolve-lemma-trace-well-defined} の写像である．
完全列
$$
H_1(L_{S/R}) \xrightarrow{\delta} \Omega_{R/\Lambda} \otimes_R S \to
\Omega_{S/\Lambda} \to \Omega_{S/R} \to 0
$$
がある（Algebra，補題 \ref{algebra-lemma-exact-sequence-NL}）．
加群 $\Omega_{S/R}$ は $S$ 上自由で基底 $\text{d}x$ をもち，
加群 $H_1(L_{S/R})$ は $S$ 上自由で基底 $x^p - a$ をもつ．後者を $\delta$ は
$-\text{d}a \otimes 1$ へ送る．これは $\Omega_{R/\Lambda} \otimes_R S$ の元である．特に，
$$
M = \Coker(R \to \Omega_{R/\Lambda}, 1 \mapsto -\text{d}a)
$$
とおくと，$\Coker(\delta) = M \otimes_R S$ であることが分かる．標準写像
$$
\Omega^{t + 1}_{S/\Lambda} \to
\wedge_S^t(\Coker(\delta)) \otimes_S \Omega_{S/R} =
\wedge^t_R(M) \otimes_R \Omega_{S/R}
$$
を得る．さて，写像 $\text{Tr} : \Omega_{S/R} \to \Omega_{R/\Lambda}$
（補題 \ref{resolve-lemma-trace-well-defined}）の像は $R\text{d}a$ に
含まれるので，像の元との外積をとると $\text{d}a$ が零になることが分かる．
したがって標準写像
$$
\wedge^t_R(M) \otimes_R \Omega_{S/R} \to \Omega_{R/\Lambda}^{t + 1}
$$
で
$\overline{\eta}_1 \wedge \ldots \wedge \overline{\eta}_t \wedge \omega$
を $\eta_1 \wedge \ldots \wedge \eta_t \wedge \text{Tr}(\omega)$ へ送るものがある．
\end{proof}

\begin{remark}
\label{resolve-remark-compare-Garel}
$\mathbf{F}_p \subset \Lambda \subset R \subset S$ および $\text{Tr}$ を
補題 \ref{resolve-lemma-trace-higher} の通りとする．de Rham Cohomology，命題
\ref{derham-proposition-Garel} により，複体の標準写像
$$
\Theta_{S/R} :
\Omega_{S/\Lambda}^\bullet
\longrightarrow
\Omega_{R/\Lambda}^\bullet
$$
がある．de Rham Cohomology，例 \ref{derham-example-Garel} の計算により，
$\Theta_{S/R}(x^i \text{d}x) = \text{Tr}_x(x^i\text{d}x)$ がすべての $i$ に
対して成り立つ．$\text{Trace}_{S/R} = \Theta^0_{S/R}$ は恒等的に零であり，また
$$
\Theta_{S/R}(a \wedge b) = a \wedge \Theta_{S/R}(b)
$$
が $a \in \Omega^i_{R/\Lambda}$ および $b \in \Omega^j_{S/\Lambda}$ に対して
成り立つので，$\text{Tr} = \Theta_{S/R}$ を得る．$\text{Tr}$ を用いる利点は，
その構成がはるかに初等的なことである．
\end{remark}

\begin{lemma}
\label{resolve-lemma-trace-extends}
$S$ を $\mathbf{F}_p$ 上の概形とする．$f : Y \to X$ を $S$ 上の Noether
正規整概形の間の有限射とする．次を仮定する．
\begin{enumerate}
\item 関数体の拡大は次数 $p$ の純非分離拡大である；
\item $\Omega_{X/S}$ は連接 $\mathcal{O}_X$-加群である（例えば，
$X$ が $S$ 上有限型である場合）．
\end{enumerate}
$i \geq 1$ に対して標準写像
$$
\text{Tr} : f_*\Omega^i_{Y/S} \longrightarrow (\Omega_{X/S}^i)^{**}
$$
% P07-ERR-1097
があり，$X$ の一般点におけるその茎は補題 \ref{resolve-lemma-trace-higher} の
トレース写像を復元する．
\end{lemma}

\begin{proof}
完全列 $f^*\Omega_{X/S} \to \Omega_{Y/S} \to \Omega_{Y/X} \to 0$ から，
$\Omega_{Y/S}$，したがって $f_*\Omega_{Y/S}$ も連接加群であることが分かる．
よって，一般点におけるトレース写像が，$x \in X$ かつ
$\dim(\mathcal{O}_{X, x}) = 1$ を満たす点での茎へ延長することを示せば十分である
（Divisors，補題 \ref{divisors-lemma-describe-reflexive-hull} を参照）．
したがって次の段落で扱う場合に帰着する．

\medskip\noindent
$X = \Spec(A)$ および $Y = \Spec(B)$ とし，$A$ は離散付値環，
$B$ は $A$ 上有限であると仮定する．誘導される分数体の拡大 $L/K$ は
純非分離なので，$B$ も局所環である．したがって $B$ も離散付値環である．
このとき，次のいずれかが成り立つ．
\begin{enumerate}
\item $B/A$ の分岐指数は $p$ であり，したがって
$B = A[x]/(x^p - a)$ である．ここで $a \in A$ は素元である；
\item $\mathfrak m_B = \mathfrak m_A B$ であり，剰余体
$B/\mathfrak m_A B$ は次数 $p$ の純非分離拡大であり，その基礎体は
$\kappa_A = A/\mathfrak m_A$ である．任意の $x \in B$ であって，その剰余類が
$\kappa_A$ に属さないものを選べば，
$B = A[x]/(x^p - a)$ となる．ここで $a \in A$ は単元である．
\end{enumerate}
$\Spec(\Lambda) \subset S$ を，$X$ が $\Spec(\Lambda)$ の中へ写るような
アフィン開部分概形とする．
補題 \ref{resolve-lemma-trace-higher} を適用すると，トレース写像が
$\Omega^i_{B/\Lambda} \to \Omega^i_{A/\Lambda}$
へ延長されることが，すべての $i \geq 1$ に対して分かる．
\end{proof}


















\section{二次変換}
\label{resolve-section-quadratic}

\noindent
本節では，曲面上の非特異点をブローアップしたときに何が起こるかを調べる．
このような射を正式に {\it 二次変換} と定義することにはためらいがある．
というのも，一方では別の名称がしばしば用いられ，他方では
「二次変換」という語が異なる意味で用いられることもあるからである．
% P07-ERR-0382: frozen source grammar defect; intended sense translated.

\begin{lemma}
\label{resolve-lemma-blowup}
$(A, \mathfrak m, \kappa)$ を次元 $2$ の正則局所環とする．
$f : X \to S = \Spec(A)$ を $A$ の $\mathfrak m$ に沿うブローアップとし，
その例外因子を $E$ とする．閉埋め込み
% P07-ERR-0383: frozen source typo "wotj"; intended sense translated.
$$
r : X \longrightarrow \mathbf{P}^1_S
$$
この閉埋め込みは，$S$ 上
\begin{enumerate}
\item $r|_E : E \to \mathbf{P}^1_\kappa$ は同型であり，
\item $\mathcal{O}_X(E) = \mathcal{O}_X(-1) =
r^*\mathcal{O}_{\mathbf{P}^1}(-1)$ であり，
\item $\mathcal{C}_{E/X} = (r|_E)^*\mathcal{O}_{\mathbf{P}^1}(1)$ かつ
$\mathcal{N}_{E/X} = (r|_E)^*\mathcal{O}_{\mathbf{P}^1}(-1)$
\end{enumerate}
を満たすものが存在する．
\end{lemma}

\begin{proof}
$A$ は次元 $2$ の正則局所環なので，$\mathfrak m = (x, y)$ と書ける．
このとき，$x$ と $y$ を次数 $1$ に置くと，Rees 代数
$\bigoplus_{n \geq 0} \mathfrak m^n$ を $A$ 上で生成する．
$X = \text{Proj}(\bigoplus_{n \geq 0} \mathfrak m^n)$ であることを思い出そう
（Divisors, 補題 \ref{divisors-lemma-blowing-up-affine}）．
したがって，全射
$$
A[T_0, T_1] \longrightarrow \bigoplus\nolimits_{n \geq 0} \mathfrak m^n,
\quad
T_0 \mapsto x,\ T_1 \mapsto y
$$
は次数付き $A$-代数の全射であり，閉埋め込み
$r : X \to \mathbf{P}^1_S = \text{Proj}(A[T_0, T_1])$
を誘導し，$\mathcal{O}_X(1) = r^*\mathcal{O}_{\mathbf{P}^1_S}(1)$ を満たす
（Constructions, 補題
\ref{constructions-lemma-surjective-graded-rings-generated-degree-1-map-proj}）．
一方，Divisors, 補題
\ref{divisors-lemma-blowing-up-gives-effective-Cartier-divisor} により
$\mathcal{O}_X(E) = \mathcal{O}_X(-1)$ だから，これで (2) が従う．

\medskip\noindent
(1) を示すには，
$$
\left(\bigoplus\nolimits_{n \geq 0} \mathfrak m^n\right) \otimes_A \kappa =
\bigoplus\nolimits_{n \geq 0} \mathfrak m^n/\mathfrak m^{n + 1} \cong
\kappa[\overline{x}, \overline{y}]
$$
が多項式代数であることに注意する
（Algebra, 補題 \ref{algebra-lemma-regular-graded}）．
これにより，$X \to S$ の $\Spec(\kappa)$ 上のファイバーは
$\text{Proj}(\kappa[\overline{x}, \overline{y}]) = \mathbf{P}^1_\kappa$ に等しい
（Constructions, 補題 \ref{constructions-lemma-base-change-map-proj}）．
$E$ は $X$ の閉部分概形であり，$\mathfrak m\mathcal{O}_X$ によって定まる．
すなわち $E = X_\kappa$ であった．射 $r$ の選び方から，$r|_E$ は実際に
$E = X_\kappa$ と $\mathbf{P}^1_S \to S$ の特殊ファイバーとの同一視を与える．

\medskip\noindent
(3) は (1), (2) および Divisors, 補題
\ref{divisors-lemma-conormal-effective-Cartier-divisor} から従う．
\end{proof}

\begin{lemma}
\label{resolve-lemma-blowup-regular}
$(A, \mathfrak m, \kappa)$ を次元 $2$ の正則局所環とする．
$f : X \to S = \Spec(A)$ を $A$ の $\mathfrak m$ に沿うブローアップとする．
このとき $X$ は既約正則スキームである．
\end{lemma}

\begin{proof}
Divisors, 補題 \ref{divisors-lemma-blow-up-integral-scheme} および
Algebra, 補題 \ref{algebra-lemma-regular-domain} により，$X$ は整である．
$X$ が正則であることを示すには，$\mathcal{O}_{X, x}$ が
閉点 $x \in X$ に対して正則であることを確かめれば十分である
（Properties, 補題 \ref{properties-lemma-characterize-regular}）．
$x \in X$ を閉点とする．$f$ は固有だから，$x$ は $\mathfrak m$ に写る．
すなわち $x$ は例外因子 $E$ の点である．
ここで $E$ は有効 Cartier 因子であり，$E \cong \mathbf{P}^1_\kappa$ である．
したがって，$g \in \mathfrak m_x \subset \mathcal{O}_{X, x}$ が
$E$ の局所方程式ならば，
$\mathcal{O}_{X, x}/(g) \cong \mathcal{O}_{\mathbf{P}^1_\kappa, x}$．
$\mathbf{P}^1_\kappa$ は，$\kappa$ 上の多項式環のスペクトルである
二つのアフィン開集合で被覆されるので，Algebra, 補題
\ref{algebra-lemma-dim-affine-space} により
$\mathcal{O}_{\mathbf{P}^1_\kappa, x}$ は正則である．
結論は Algebra, 補題 \ref{algebra-lemma-regular-mod-x} から従う．
\end{proof}

\begin{lemma}
\label{resolve-lemma-blowup-pic}
$(A, \mathfrak m, \kappa)$ を次元 $2$ の正則局所環とする．
$f : X \to S = \Spec(A)$ を $A$ の $\mathfrak m$ に沿うブローアップとする．
このとき $\Pic(X) = \mathbf{Z}$ であり，$\mathcal{O}_X(E)$ により生成される．
\end{lemma}

\begin{proof}
$E = \mathbf{P}^1_\kappa$ の Picard 群は $\mathbf{Z}$ で，
$\mathcal{O}(1)$ により生成されることを思い出そう
（Divisors, 補題 \ref{divisors-lemma-Pic-projective-space-UFD}）．
補題 \ref{resolve-lemma-blowup} により，可逆 $\mathcal{O}_X$-加群
$\mathcal{O}_X(E)$ の例外因子への制限は $\mathcal{O}(-1)$ である．
したがって $\mathcal{O}_X(E)$ は $\Pic(X)$ の無限巡回部分群を生成する．
$A$ は正則なので UFD である
（More on Algebra, 補題 \ref{more-algebra-lemma-regular-local-UFD}）．
したがって穴あきスペクトル
$U = S \setminus \{\mathfrak m\} = X \setminus E$ の Picard 群は自明である
（Divisors, 補題 \ref{divisors-lemma-open-subscheme-UFD}）．
ゆえに任意の可逆 $\mathcal{O}_X$-加群 $\mathcal{L}$ に対し，
同型 $s : \mathcal{O}_U \to \mathcal{L}|_U$ が存在する．
このとき $s$ は $\mathcal{L}$ の正則有理型切断であり，
$\text{div}_\mathcal{L}(s) = nE$ がある $n \in \mathbf{Z}$ に対して成り立つ
（Divisors, 定義 \ref{divisors-definition-divisor-invertible-sheaf}）．
Divisors, 補題 \ref{divisors-lemma-normal-c1-injective} と
補題 \ref{resolve-lemma-blowup-regular} による $X$ の正規性から，
$\mathcal{L} = \mathcal{O}_X(nE)$ と結論する．
\end{proof}

\begin{lemma}
\label{resolve-lemma-cohomology-of-blowup}
$(A, \mathfrak m, \kappa)$ を次元 $2$ の正則局所環とする．
$f : X \to S = \Spec(A)$ を $A$ の $\mathfrak m$ に沿うブローアップとする．
$\mathcal{F}$ を準連接 $\mathcal{O}_X$-加群とする．
\begin{enumerate}
\item $H^p(X, \mathcal{F}) = 0$ は $p \not \in \{0, 1\}$ に対して成り立つ．
\item $H^1(X, \mathcal{O}_X(n)) = 0$ は $n \geq -1$ に対して成り立つ．
\item $H^1(X, \mathcal{F}) = 0$ は，$\mathcal{F}$ または $\mathcal{F}(1)$
が大域生成されるときに成り立つ．
\item $H^0(X, \mathcal{O}_X(n)) = \mathfrak m^{\max(0, n)}$．
\item $\text{length}_A H^1(X, \mathcal{O}_X(n)) = -n(-n - 1)/2$
は $n < 0$ のとき成り立つ．
\end{enumerate}
\end{lemma}

\begin{proof}
$\mathfrak m = (x, y)$ とする．$X$ はアフィン・ブローアップ代数
$A[\frac{\mathfrak m}{x}]$ と $A[\frac{\mathfrak m}{y}]$ の
スペクトルで被覆される．これは，$x$ と $y$ を次数 $1$ に置くと
Rees 代数 $\bigoplus \mathfrak m^n$ を $A$ 上で生成するからである．
Divisors, 補題 \ref{divisors-lemma-blowing-up-affine} および
Constructions, 補題 \ref{constructions-lemma-proj-quasi-compact} を参照せよ．
Constructions, 補題 \ref{constructions-lemma-proj-separated} により
$X$ は分離的だから，Cohomology of Schemes, 補題
\ref{coherent-lemma-vanishing-nr-affines} により，準連接層のコホモロジーは
次数 $\geq 2$ で消える．

\medskip\noindent
$i : E \to X$ を例外因子の埋め込みとする
（Divisors, 定義 \ref{divisors-definition-blow-up}）．
$\mathcal{O}_X(-E) = \mathcal{O}_X(1)$ は $f$-相対的に豊富であった
（Divisors, 補題
\ref{divisors-lemma-blowing-up-gives-effective-Cartier-divisor}）．
したがって，$H^1(X, \mathcal{O}_X(-nE)) = 0$ となる
$n > 0$ が存在する
（Cohomology of Schemes, 補題 \ref{coherent-lemma-kill-by-twisting}）．
フィルトレーション
$$
\mathcal{O}_X(-nE) \subset \mathcal{O}_X(-(n - 1)E) \subset
\ldots \subset \mathcal{O}_X(-E) \subset \mathcal{O}_X \subset \mathcal{O}_X(E)
$$
を考える．その逐次商は層
$$
\mathcal{O}_X(-t E)/\mathcal{O}_X(-(t + 1)E) =
\mathcal{O}_X(t)/\mathcal{I}(t) =
i_*\mathcal{O}_E(t)
$$
である．ここで $\mathcal{I} = \mathcal{O}_X(-E)$ は $E$ のイデアル層である．
補題 \ref{resolve-lemma-blowup} により $E = \mathbf{P}^1_\kappa$ であり，
$\mathcal{O}_E(1)$ は実際に $\mathbf{P}^1$ 上の構造層の通常の
Serre 捻りに対応する．したがって $\mathcal{O}_E(t)$ の次数 $1$ の
コホモロジーは $t \geq -1$ のとき消える
（Cohomology of Schemes, 補題
\ref{coherent-lemma-cohomology-projective-space-over-ring}）．
これは $H^1(X, i_*\mathcal{O}_E(t))$ に等しい
（Cohomology of Schemes, 補題
\ref{coherent-lemma-relative-affine-cohomology}）から，
$H^1(X, \mathcal{O}_X(-(t + 1)E)) \to H^1(X, \mathcal{O}_X(-tE))$
は $t \geq -1$ に対して全射である．ゆえに
$$
0 = H^1(X, \mathcal{O}_X(-nE))
\longrightarrow
H^1(X, \mathcal{O}_X(-tE)) = H^1(X, \mathcal{O}_X(t))
$$
は $t \geq -1$ に対して全射であり，(2) が示された．

\medskip\noindent
$\mathcal{F}$ が大域生成されるとする．これは短完全列
$$
0 \to \mathcal{G} \to \bigoplus\nolimits_{i \in I} \mathcal{O}_X
\to \mathcal{F} \to 0
$$
が存在することを意味する．
$H^1(X, \bigoplus_{i \in I} \mathcal{O}_X) =
\bigoplus_{i \in I} H^1(X, \mathcal{O}_X)$ であることは
Cohomology, 補題 \ref{cohomology-lemma-quasi-separated-cohomology-colimit}
による．(2) により
$H^1(X, \mathcal{O}_X) = 0$ である．
$\mathcal{F}(1)$ が大域生成される場合には，全射
$\bigoplus_{i \in I} \mathcal{O}_X(-1) \to \mathcal{F}$ を取り，
同様に論じればよい．すなわち (3) は (2) から従う．

\medskip\noindent
(4) について，十分大きいすべての $n$ に対して
$\Gamma(X, \mathcal{O}_X(n)) = \mathfrak m^n$ であることに注意する
（Cohomology of Schemes, 補題
\ref{coherent-lemma-recover-tail-graded-module}）．
$n \geq 0$ ならば，短完全列
$$
0 \to \mathcal{O}_X(n) \to \mathcal{O}_X(n - 1) \to
i_*\mathcal{O}_E(n - 1) \to 0
$$
と左端の層の $H^1$ の消滅を用いて，可換図式
% P07-ERR-1098: frozen source quotient order is reversed; formula retained.
$$
\xymatrix{
0 \ar[r] &
\mathfrak m^{\max(0, n)} \ar[r] \ar[d] &
\mathfrak m^{\max(0, n - 1)} \ar[r] \ar[d] &
\mathfrak m^{\max(0, n)}/\mathfrak m^{\max(0, n - 1)} \ar[r] \ar[d] & 0\\
0 \ar[r] &
\Gamma(X, \mathcal{O}_X(n)) \ar[r] &
\Gamma(X, \mathcal{O}_X(n - 1)) \ar[r] &
\Gamma(E, \mathcal{O}_E(n - 1)) \ar[r] & 0
}
$$
を得る．その各行は完全である．実際，$n < 0$ の場合も，右側の群が
零なので各行は完全である．補題 \ref{resolve-lemma-blowup} の証明で，
右の縦射が同型であることを見た（詳細は省く）．したがって，
左の縦射が同型ならば中央の縦射も同型である．このようにして，
$n$ に関する降下帰納法により (4) が成り立つことが分かる．

\medskip\noindent
最後に，$n$ に関する降下帰納法と列
$$
0 \to \mathcal{O}_X(n) \to \mathcal{O}_X(n - 1) \to
i_*\mathcal{O}_E(n - 1) \to 0
$$
を用いて (5) を示す．$n \geq -1$ ならば
$H^1(X, \mathcal{O}_X(n)) = 0$ であることは既に分かっている．また，
$$
H^1(X, i_*\mathcal{O}_E(-2)) =
H^1(E, \mathcal{O}_E(-2)) =
H^1(\mathbf{P}^1_\kappa, \mathcal{O}(-2)) \cong \kappa
$$
は Cohomology of Schemes, 補題
\ref{coherent-lemma-cohomology-projective-space-over-ring} により成り立ち，
これは長さ $1$ の $A$-加群である．したがってコホモロジー長完全列から，
$n = -2$ に対して (5) が成り立つ．以下も同様である．
\end{proof}

\begin{lemma}
\label{resolve-lemma-blowup-improve}
$(A, \mathfrak m)$ を次元 $2$ の正則局所環とする．
$f : X \to S = \Spec(A)$ を $A$ の $\mathfrak m$ に沿うブローアップとする．
$\mathfrak m^n \subset I \subset \mathfrak m$ をイデアルとする．
$d \geq 0$ は，
$$
I \mathcal{O}_X \subset \mathcal{O}_X(-dE)
$$
を満たす最大の整数であるとする．ここで $E$ は例外因子である．
$\mathcal{I}' = I\mathcal{O}_X(dE) \subset \mathcal{O}_X$ とおく．
このとき $d > 0$ であり，層 $\mathcal{O}_X/\mathcal{I}'$ は
$x_1, \ldots, x_r$ という $X$ の有限個の閉点に台をもち，さらに
\begin{align*}
\text{length}_A(A/I)
& >
\text{length}_A \Gamma(X, \mathcal{O}_X/\mathcal{I}') \\
& \geq
\sum\nolimits_{i = 1, \ldots, r}
\text{length}_{\mathcal{O}_{X, x_i}}
(\mathcal{O}_{X, x_i}/\mathcal{I}'_{x_i})
\end{align*}
\end{lemma}

\begin{proof}
$I \subset \mathfrak m$ なので，$I$ のすべての元は $E$ 上で消える．
したがって $d \geq 1$ である．一方，
$\mathfrak m^n \subset I$ なので $d \leq n$ である．短完全列
$$
0 \to I\mathcal{O}_X \to \mathcal{O}_X \to \mathcal{O}_X/I\mathcal{O}_X \to 0
$$
を考える．$I\mathcal{O}_X$ は大域生成されるので，補題
\ref{resolve-lemma-cohomology-of-blowup} により
$H^1(X, I\mathcal{O}_X) = 0$ である．ゆえに全射
$A/I \to \Gamma(X, \mathcal{O}_X/I\mathcal{O}_X)$ を得る．さらに短完全列
$$
0 \to
\mathcal{O}_X(-dE)/I\mathcal{O}_X \to
\mathcal{O}_X/I\mathcal{O}_X \to
\mathcal{O}_X/\mathcal{O}_X(-dE) \to 0
$$
を考える．Divisors, 補題 \ref{divisors-lemma-codim-1-part} により，
$\mathcal{O}_X(-dE)/I\mathcal{O}_X$ は $X$ の有限個の閉点に台をもつ．
特に，この連接層の高次コホモロジー群は消える（詳細は省く）．
したがって，次の図式において
$$
\xymatrix{
& & A/I \ar[d] \\
0 \ar[r] &
\Gamma(X, \mathcal{O}_X(-dE)/I\mathcal{O}_X) \ar[r] &
\Gamma(X, \mathcal{O}_X/I\mathcal{O}_X) \ar[r] &
\Gamma(X, \mathcal{O}_X/\mathcal{O}_X(-dE)) \ar[r] & 0
}
$$
下の行は完全であり，縦射は全射である．
$$
\text{length}_A \Gamma(X, \mathcal{O}_X(-dE)/I\mathcal{O}_X) <
\text{length}_A(A/I)
$$
が成り立つ．実際，
$\Gamma(X, \mathcal{O}_X/\mathcal{O}_X(-dE))$ は零でない．
$1 \in \Gamma(X, \mathcal{O}_X)$ の像は，$d > 0$ だから零でないのである．

\medskip\noindent
証明を終えるため，上の結果を補題の主張に読み替える．
$\mathcal{O}_X(dE)$ は可逆なので，
$$
\mathcal{O}_X/\mathcal{I}' =
\mathcal{O}_X(-dE)/I\mathcal{O}_X \otimes_{\mathcal{O}_X} \mathcal{O}_X(dE).
$$
したがって $\mathcal{O}_X/\mathcal{I}'$ と
$\mathcal{O}_X(-dE)/I\mathcal{O}_X$ は，同じ有限個の閉点，
たとえば $x_1, \ldots, x_r \in E \subset X$ に台をもつ．さらに
$$
\Gamma(X, \mathcal{O}_X(-dE)/I\mathcal{O}_X) =
\bigoplus \mathcal{O}_X(-dE)_{x_i}/I\mathcal{O}_{X, x_i}
\cong
\bigoplus \mathcal{O}_{X, x_i}/\mathcal{I}'_{x_i} =
\Gamma(X, \mathcal{O}_X/\mathcal{I}')
$$
を得る．これは局所環上の可逆加群が自明だからである．
これで狭義不等式を得る．また，
$$
\text{length}_A(\mathcal{O}_{X, x_i}/\mathcal{I}'_{x_i}) \geq
\text{length}_{\mathcal{O}_{X, x_i}}(\mathcal{O}_{X, x_i}/\mathcal{I}'_{x_i})
$$
は長さの定義から直ちに分かるので，二番目の不等式も得られる．
\end{proof}

\begin{lemma}
\label{resolve-lemma-differentials-of-blowup}
$(A, \mathfrak m, \kappa)$ を次元 $2$ の正則局所環とする．
$f : X \to S = \Spec(A)$ を $A$ の $\mathfrak m$ に沿うブローアップとする．
このとき $\Omega_{X/S} = i_*\Omega_{E/\kappa}$ である．ここで
$i : E \to X$ は例外因子の埋め込みである．
\end{lemma}

\begin{proof}
$\mathbf{P}^1 = \mathbf{P}^1_S$ と書き，$r : X \to \mathbf{P}^1$ を
補題 \ref{resolve-lemma-blowup} のものとする．このとき完全列
$$
\mathcal{C}_{X/\mathbf{P}^1} \to r^*\Omega_{\mathbf{P}^1/S} \to
\Omega_{X/S} \to 0
$$
がある
（Morphisms, 補題 \ref{morphisms-lemma-differentials-relative-immersion}）．
Morphisms, 補題 \ref{morphisms-lemma-base-change-differentials} により
$\Omega_{\mathbf{P}^1/S}|_E = \Omega_{E/\kappa}$ だから，最初の射が
標準写像 $r^*\Omega_{\mathbf{P}^1/S} \to i_*\Omega_{E/\kappa}$ の
核への全射を定めることを示せば十分である．これは局所的に示せる．
補題 \ref{resolve-lemma-blowup} の証明と同じ記号を用いると，$X$ の
あるアフィン開集合上で，射 $f$ は環準同型
$$
A \to A[t]/(xt - y)
$$
に対応する．ここで $x, y \in \mathfrak m$ は生成元である．したがって
$\text{d}(xt - y) = x\text{d}t$ かつ
$y\text{d}t = t \cdot x \text{d}t$ であり，求める主張が従う．
\end{proof}



\section{二次変換による支配}
\label{resolve-section-dominating-by-quadratic}

\noindent
上の結果を用いると，点に沿うブローアップが，正則な $2$ 次元スキームの
任意の改変を支配することを示せる．

\medskip\noindent
$X$ をスキームとし，$x \in X$ を閉点とする．通常どおり，
$i : x = \Spec(\kappa(x)) \to X$ を閉部分概形とみなす．
{\it ブローアップ $X' \to X$，すなわち $X$ の $x$ におけるブローアップ}
とは，$X$ の閉部分概形 $x \subset X$ に沿うブローアップである．
$X$ が局所 Noether スキームならば，Divisors, 補題
\ref{divisors-lemma-blowing-up-projective} により
$X' \to X$ は射影的（特に固有）であることに注意せよ．

\begin{lemma}
\label{resolve-lemma-make-ideal-principal}
$X$ を Noether スキームとする．$T \subset X$ を閉点 $x$ の有限集合で，
$\mathcal{O}_{X, x}$ が次元 $2$ の正則環となることが
$x \in T$ に対して成り立つものとする．
$\mathcal{I} \subset \mathcal{O}_X$ を準連接イデアル層で，
$\mathcal{O}_X/\mathcal{I}$ が $T$ に台をもつものとする．
このとき列
$$
X_n \to X_{n - 1} \to \ldots \to X_1 \to X_0 = X
$$
で，各 $X_{i + 1} \to X_i$ が $X_i$ の閉点で $T$ の点の上にあるものに沿う
ブローアップであり，$\mathcal{I}\mathcal{O}_{X_n}$ が可逆イデアル層となる
ものが存在する．
\end{lemma}

\begin{proof}
$T = \{x_1, \ldots, x_r\}$ と書く．$I_i$ で $\mathcal{I}$ の
$x_i$ における茎を表す．
$$
n_i = \text{length}_{\mathcal{O}_{X, x_i}}(\mathcal{O}_{X, x_i}/I_i)
$$
とおく．$\mathcal{O}_X/\mathcal{I}$ は $T$ に台をもつので，これは有限である．
したがって $\mathcal{O}_{X, x_i}/I_i$ の台は
$\{\mathfrak m_{x_i}\}$ に等しい
（Algebra, 補題 \ref{algebra-lemma-support-point}）．
$\sum n_i$ に関する帰納法を用いる．$n_i = 0$ がすべての
$i$ に対して成り立つならば，$\mathcal{I} = \mathcal{O}_X$ であり，証明は終わる．

\medskip\noindent
$n_i > 0$ とする．$X' \to X$ を $X$ の $x_i$ に沿うブローアップとする
（補題の前の議論を参照）．
$\Spec(\mathcal{O}_{X, x_i}) \to X$ は平坦なので，
$X' \times_X \Spec(\mathcal{O}_{X, x_i})$ は環
$\mathcal{O}_{X, x_i}$ の極大イデアルに沿うブローアップである
（Divisors, 補題
\ref{divisors-lemma-flat-base-change-blowing-up}）．
したがって，可換図式
$$
\xymatrix{
\text{Proj}(\bigoplus\nolimits_{d \geq 0} \mathfrak m_{x_i}^d) \ar[r] \ar[d] &
X' \ar[d] \\
\Spec(\mathcal{O}_{X, x_i}) \ar[r] & X
}
$$
の正方形は引き戻し正方形である．$E \subset X'$ および
$E' \subset \text{Proj}(\bigoplus\nolimits_{d \geq 0} \mathfrak m_{x_i}^d)$
を例外因子とする．$d \geq 1$ を，イデアル
% P07-ERR-1099: frozen source switches from I_i to undefined \mathcal{I}_i.
$\mathcal{I}_i \subset \mathcal{O}_{X, x_i}$ に対して
補題 \ref{resolve-lemma-blowup-improve} で得られる整数とする．
図式の横射は平坦であり，$E' \to E$ は全射であり，かつ $E'$ は $E$ の
引き戻しなので，
$$
\mathcal{I}\mathcal{O}_{X'} \subset \mathcal{O}_{X'}(-dE)
$$
を得る（詳細の一部は省く）．
$\mathcal{I}' = \mathcal{I}\mathcal{O}_{X'}(dE) \subset \mathcal{O}_{X'}$
とおく．このとき $\mathcal{O}_{X'}/\mathcal{I}'$ は有限個の閉点
$T' \subset |X'|$ に台をもつ．実際，これは
$X \setminus \{x_i\}$ 上でも，
$\text{Proj}(\bigoplus\nolimits_{d \geq 0} \mathfrak m_{x_i}^d)$
への引き戻しについても成り立つ．補題
\ref{resolve-lemma-blowup-improve} の最後の主張により，茎
$\mathcal{O}_{X', x'}/\mathcal{I}'\mathcal{O}_{X', x'}$ の長さを
$x'$ が $x_i$ の上にあるすべての場合について足した総和は
$< n_i$ である．したがって長さの総和は減少した．

\medskip\noindent
帰納法の仮定により，列
$$
X'_n \to \ldots \to X'_1 \to X'
$$
で，$T'$ の上にある閉点でのブローアップからなり，
$\mathcal{I}'\mathcal{O}_{X'_n}$ が可逆となるものが存在する．
$\mathcal{I}'\mathcal{O}_{X'}(-dE) = \mathcal{I}\mathcal{O}_{X'}$ だから，
$\mathcal{I}\mathcal{O}_{X'_n} =
\mathcal{I}'\mathcal{O}_{X'_n}(-d(f')^{-1}E)$
である．ここで $f' : X'_n \to X'$ は合成射である．
Divisors, 補題
\ref{divisors-lemma-blow-up-pullback-effective-Cartier} により
$(f')^{-1}E$ は有効 Cartier 因子であることに注意せよ．
したがって Divisors, 補題
\ref{divisors-lemma-sum-effective-Cartier-divisors} により証明が終わる．
\end{proof}

\begin{lemma}
\label{resolve-lemma-dominate-by-blowing-up-in-points}
$X$ を Noether スキームとする．$T \subset X$ を閉点 $x$ の有限集合で，
$\mathcal{O}_{X, x}$ が次元 $2$ の正則局所環となるものとする．
$f : Y \to X$ を，$U = X \setminus T$ 上で同型となるスキームの固有射とする．
このとき列
$$
X_n \to X_{n - 1} \to \ldots \to X_1 \to X_0 = X
$$
で，各 $X_{i + 1} \to X_i$ が，$X_i$ の閉点 $x_i$ で
$T$ の点の上にあるものに沿うブローアップであり，かつその合成が
$X_n \to Y \to X$ と分解するものが存在する．
\end{lemma}

\begin{proof}
More on Flatness, 補題 \ref{flat-lemma-dominate-modification-by-blowup}
により，$U$-許容ブローアップ $X' \to X$ で $Y \to X$ を支配するものが存在する．
したがって，イデアル層 $\mathcal{I} \subset \mathcal{O}_X$ で，
$\mathcal{O}_X/\mathcal{I}$ が $T$ に台をもち，$Y$ が $X$ の
$\mathcal{I}$ に沿うブローアップとなるものが存在すると仮定してよい．
補題 \ref{resolve-lemma-make-ideal-principal} により，列
$$
X_n \to X_{n - 1} \to \ldots \to X_1 \to X_0 = X
$$
で，各 $X_{i + 1} \to X_i$ が，$X_i$ の閉点 $x_i$ で
$T$ の点の上にあるものに沿うブローアップであり，
$\mathcal{I}\mathcal{O}_{X_n}$ が可逆イデアル層となるものが存在する．
ブローアップの普遍性
（Divisors, 補題
\ref{divisors-lemma-universal-property-blowing-up}）により，
求める分解を得る．
\end{proof}

\begin{lemma}
\label{resolve-lemma-extend-rational-map-blowing-up}
% P07-ERR-1100: global finiteness step lacks a Noetherian/quasi-compact hypothesis.
$S$ をスキームとする．$X$ を $S$ 上の正則な次元 $2$ のスキームとし，
$Y$ を $S$ 上の固有スキームとする．$S$-有理写像
$f : U \to Y$ が $X$ から $Y$ へ与えられたとする．このとき列
$$
X_n \to X_{n - 1} \to \ldots \to X_1 \to X_0 = X
$$
と $S$-射 $f_n : X_n \to Y$ で，$X_{i + 1} \to X_i$ が
$X_i$ の閉点で $U$ の上にないものに沿うブローアップであり，
$f_n$ と $f$ が一致するものが存在する．
\end{lemma}

\begin{proof}
$U$ は余次元 $1$ のすべての点を含むと仮定してよい
（Morphisms, 補題 \ref{morphisms-lemma-extend-across}）．
したがって $T \subset X$ は $U$ の補集合であり，有限個の閉点からなり，
その局所環は次元 $2$ の正則環である．
Divisors, 補題
\ref{divisors-lemma-extend-rational-map-after-modification} を適用すると，
固有射 $p : X' \to X$ で $U$ 上で同型となるものと，射
$f' : X' \to Y$ で $f$ と $U$ 上で一致するものを得る．
射 $p : X' \to X$ に補題
\ref{resolve-lemma-dominate-by-blowing-up-in-points} を適用する．合成
$X_n \to X' \to Y$ が求める射である．
\end{proof}






\section{正規化ブローアップによる支配}
\label{resolve-section-normalized-blowups}

\noindent
本節では，曲面の改変を点における正規化ブローアップの列によって
支配できることを証明する．

\begin{definition}
\label{resolve-definition-normalized-blowup}
$X$ を，任意の準コンパクト開部分が有限個の既約成分をもつスキームとする．
$x \in X$ を閉点とする．{\it $X$ の $x$ における正規化ブローアップ}とは，
合成 $X'' \to X' \to X$ のことである．ここで $X' \to X$ は
$X$ の $x$ に沿うブローアップであり，$X'' \to X'$ は $X'$ の正規化である．
\end{definition}

\noindent
ここで正規化 $X'' \to X'$ は定義される．実際，
Divisors, 補題 \ref{divisors-lemma-blow-up-and-irreducible-components}
により，スキーム $X'$ は有限個の既約成分をもつ開部分による開被覆をもつ．
正規化の定義については Morphisms, 定義
\ref{morphisms-definition-normalization} を参照せよ．

\medskip\noindent
一般に，$X$ が Noether スキームであっても，正規化ブローアップは
固有であるとは限らない．スキームが Nagata であるとは，Nagata 環の
スペクトルであるアフィン開集合による開被覆をもつことをいう
（Properties, 定義 \ref{properties-definition-nagata}）．

\begin{lemma}
\label{resolve-lemma-Nagata-normalized-blowup}
定義 \ref{resolve-definition-normalized-blowup} において $X$ が Nagata ならば，
$X$ の $x$ における正規化ブローアップは正規かつ Nagata であり，
$X$ 上固有である．
\end{lemma}

\begin{proof}
ブローアップ射 $X' \to X$ は固有である
（$X$ は局所 Noether なので Divisors, 補題
\ref{divisors-lemma-blowing-up-projective} を適用できる）．
したがって $X'$ は Nagata である
（Morphisms, 補題 \ref{morphisms-lemma-finite-type-nagata}）．
ゆえに正規化 $X'' \to X'$ は有限であり
（Morphisms, 補題 \ref{morphisms-lemma-nagata-normalization}），
$X'' \to X$ も固有であると結論できる
（Morphisms, 補題 \ref{morphisms-lemma-finite-proper} および
\ref{morphisms-lemma-composition-proper}）．したがって正規化ブローアップは，
正規な（Morphisms, 補題
\ref{morphisms-lemma-normalization-normal}）Nagata 代数空間である．
\end{proof}

\noindent
次の補題では，$X$ が有限個の既約成分をもつことを保証するため，
Noether であると仮定する必要がある．すると $f : Y \to X$ の固有性により，
$Y$ も有限個の既約成分をもつことが保証され，$f$ が双有理であると
要求することに意味がある
（Morphisms, 定義 \ref{morphisms-definition-birational}）．

\begin{lemma}
\label{resolve-lemma-dominate-by-normalized-blowing-up}
$X$ を Noether かつ Nagata で次元 $2$ のスキームとする．
$f : Y \to X$ を固有双有理射とする．このとき可換図式
$$
\xymatrix{
X_n \ar[r] \ar[d] &
X_{n - 1} \ar[r] &
\ldots \ar[r] &
X_1 \ar[r] &
X_0 \ar[d] \\
Y \ar[rrrr]  & & & & X
}
$$
で，$X_0 \to X$ が正規化であり，$X_{i + 1} \to X_i$ が
$X_i$ のある閉点における正規化ブローアップであるものが存在する．
\end{lemma}

\begin{proof}
以下では，Morphisms, 節
\ref{morphisms-section-nagata},
\ref{morphisms-section-dimension-formula}, および
\ref{morphisms-section-normalization} の結果を断りなく用いる．
$Y$ をその正規化で置き換えてよい．$X_0 \to X$ を正規化とする．
射 $Y \to X$ は $X_0$ を経由して分解する．したがって，
$X$ と $Y$ はともに正規であると仮定してよい．

\medskip\noindent
% P07-ERR-1101: frozen source assigns the codimension/fibre loci to Y, but the cited lemma places them in X.
$X$ と $Y$ は正規であると仮定する．射 $f : Y \to X$ は，余次元
$0$ および $1$ である $Y$ のすべての点と，ファイバーが有限となる $Y$ のすべての点を
含む開部分上で同型である．Varieties, 補題
\ref{varieties-lemma-modification-normal-iso-over-codimension-1} を参照せよ．
したがって，閉点の有限集合 $T \subset X$ で，$f$ が
$X \setminus T$ 上で同型となるものが存在する．各 $x \in T$ に対し，
ファイバー $Y_x$ は次元 $1$ で，$\kappa(x)$ 上固有な幾何学的連結スキームである．
More on Morphisms, 補題
\ref{more-morphisms-lemma-geometrically-connected-fibres-towards-normal} を参照せよ．
したがって
$$
BadCurves(f) = \{C \subset Y\text{ は閉} \mid
\dim(C) = 1, f(C) = \text{一点}\}
$$
は有限集合である．$BadCurves(f)$ の元の個数に関する帰納法で補題を証明する．
基底の場合は $BadCurves(f)$ が空の場合であり，このとき $f$ は同型である．

\medskip\noindent
$x \in T$ を固定する．$X' \to X$ を $X$ の $x$ における正規化ブローアップとし，
$Y'$ を $Y \times_X X'$ の正規化とする．図式は
$$
\xymatrix{
Y' \ar[r]_{f'} \ar[d] & X' \ar[d] \\
Y \ar[r]^f & X
}
$$
である．$x' \in X'$ を $x$ の上にある閉点で，ファイバー $Y'_{x'}$ の次元が
$\geq 1$ となるものとする．$C' \subset Y'$ を $Y'_{x'}$ の既約成分，
すなわち $C' \in BadCurves(f')$ とする．$Y' \to Y \times_X X'$ は有限なので，
$C'$ はある既約成分 $C \subset Y_x$ に写らなければならない．
$C \in BadCurves(f)$ であることは明らかである．$Y' \to Y$ は双有理であり，
したがって余次元 $1$ である $Y$ の点上で同型なので，単射
$$
BadCurves(f') \longrightarrow BadCurves(f)
$$
を得る．したがって，これらの正規化ブローアップを有限回行った後に，
少なくとも一つの悪い曲線を除去できること，すなわち表示した写像が
全射でないことを示せば十分である．

\medskip\noindent
Zariski による議論を用いて悪い曲線を除去する．
$C \in BadCurves(f)$ をとり，これは先の $x$ の上にあるとする．
$\mathcal{O}_{Y, C}$ で，$Y$ の $C$ の一般点における局所環を表す．
% P07-ERR-1102: frozen source uses the undefined local ring O_{X,C}; C is a curve in Y.
$u \in \mathcal{O}_{X, C}$ を，剰余体 $R(C)$ における像が $\kappa(x)$ 上
超越的となる元として選ぶ（これは Varieties, 補題
\ref{varieties-lemma-dimension-locally-algebraic} により，$R(C)$ の
超越次数が $1$ の $\kappa(x)$-拡大だから可能である）．
$u = a/b$ を $a, b \in \mathcal{O}_{X, x}$ として書ける．これは
$\mathcal{O}_{Y, C}$ と $\mathcal{O}_{X, x}$ が同じ分数体をもつからである．
$u$ の選び方により，$a, b \in \mathfrak m_x$ でなければならない．したがって
$$
N_{u, a, b} = \min
\{\text{ord}_{\mathcal{O}_{Y, C}}(a), \text{ord}_{\mathcal{O}_{Y, C}}(b)\} > 0
$$
よって，この整数に関する降下帰納法を行える．$X' \to X$ を $x$ の
正規化ブローアップとし，$Y'$ を上と同様に $X' \times_X Y$ の正規化とする．
$C$ が，ある悪い曲線 $C' \subset Y'$ で $x' \in X'$ の上にあるものの像ならば，
% P07-ERR-0386: frozen source omits \in between a', b' and the local ring.
$a', b' \mathcal{O}_{X', x'}$ の選択で
$N_{u, a', b'} < N_{u, a, b}$ を満たすものが存在することを示す．
これで証明が終わる．実際，$X' \to X$ はブローアップを経由して分解するので，
非零元 $d \in \mathfrak m_{x'}$ で，$a = a' d$ および $b = b' d$ を
満たすものが存在する（すなわち，$d$ としてブローアップの例外因子の
局所方程式をとる）．$Y' \to Y$ は $C$ の一般点を含む開部分上で
同型なので（上で見た），$\mathcal{O}_{Y', C'} = \mathcal{O}_{Y, C}$ である．
したがって
$$
\text{ord}_{\mathcal{O}_{Y, C}}(a) =
\text{ord}_{\mathcal{O}_{Y', C'}}(a' d) =
\text{ord}_{\mathcal{O}_{Y', C'}}(a') +
\text{ord}_{\mathcal{O}_{Y', C'}}(d) >
\text{ord}_{\mathcal{O}_{Y', C'}}(a')
$$
$b$ についても同様であり，証明は完了する．
\end{proof}

\begin{lemma}
\label{resolve-lemma-extend-rational-map-normalized-blowing-up}
$S$ をスキームとする．$X$ を $S$ 上の Noether かつ Nagata で
次元 $2$ のスキームとする．$Y$ を $S$ 上の固有スキームとする．
$S$-有理写像 $f : U \to Y$ が $X$ から $Y$ へ与えられたとき，列
$$
X_n \to X_{n - 1} \to \ldots \to X_1 \to X_0 \to X
$$
と $S$-射 $f_n : X_n \to Y$ で，$X_0 \to X$ が正規化，
$X_{i + 1} \to X_i$ が $X_i$ のある閉点における正規化ブローアップであり，
$f_n$ と $f$ が一致するものが存在する．
\end{lemma}

\begin{proof}
Divisors, 補題 \ref{divisors-lemma-extend-rational-map-after-modification}
を適用すると，固有射 $p : X' \to X$ で $U$ 上で同型となるものと，
射 $f' : X' \to Y$ で $f$ と $U$ 上で一致するものを得る．
射 $p : X' \to X$ に補題 \ref{resolve-lemma-dominate-by-normalized-blowing-up}
を適用する．合成 $X_n \to X' \to Y$ が求める射である．
\end{proof}




\section{局所環上での改変}
\label{resolve-section-modifications}

\noindent
$S$ をスキームとする．$s_1, \ldots, s_n \in S$ を相異なる閉点とする．
開埋め込み
$$
U = S \setminus \{s_1, \ldots, s_n\} \longrightarrow S
$$
が準コンパクトであると仮定する．$FP_{S, \{s_1, \ldots, s_n\}}$ で，
有限表示射 $f : X \to S$ のうち，同型 $f^{-1}(U) \to U$ を誘導するものの
圏を表す．射は $S$ 上のスキームの射である．各 $i$ に対して
$S_i = \Spec(\mathcal{O}_{S, s_i})$ とおき，$V_i = S_i \setminus \{s_i\}$
とする．$FP_{S_i, s_i}$ で，有限表示射 $g_i : Y_i \to S_i$ のうち，
同型 $g_i^{-1}(V_i) \to V_i$ を誘導するものの圏を表す．
射は $S_i$ 上の射である．基底変換により関手
\begin{equation}
\label{resolve-equation-equivalence}
F :
FP_{S, \{s_1, \ldots, s_n\}}
\longrightarrow
FP_{S_1, s_1} \times \ldots \times FP_{S_n, s_n}
\end{equation}
が定まる．本章の問題の少なくとも一部を局所環の場合へ帰着するため，
次の補題を用いる．

\begin{lemma}
\label{resolve-lemma-equivalence}
関手 $F$（\ref{resolve-equation-equivalence}）は圏同値である．
\end{lemma}

\begin{proof}
$n = 1$ の場合，これは Limits, 補題 \ref{limits-lemma-modifications} である．
$n > 1$ の場合もまったく同様に証明でき，またはこの場合から導ける．
例えば，$g_i : Y_i \to S_i$ を $FP_{S_i, s_i}$ の対象とする．
すると $n = 1$ の場合により，有限表示射 $f'_i : X'_i \to S$ で，
$S \setminus \{s_i\}$ 上同型であり，$S_i$ への基底変換が $g_i$ となるものを
見つけられる．そこで
$$
f : X = X'_1 \times_S \ldots \times_S X'_n \to S
$$
とおく．これは $FP_{S, \{s_1, \ldots, s_n\}}$ の対象であり，
$S_i \to S$ による基底変換で $g_i$ が復元される．したがって関手は
本質的全射である．充満忠実性の証明は省略する．
\end{proof}

\begin{lemma}
\label{resolve-lemma-equivalence-properties}
$S, s_i, S_i$ は（\ref{resolve-equation-equivalence}）のものとする．
$f : X \to S$ が $g_i : Y_i \to S_i$ に $F$ のもとで対応するならば，
$f$ が分離的，固有，有限であることは，それぞれ $g_i$ が
$i = 1, \ldots, n$ のすべてについてそうであることと同値である．
\end{lemma}

\begin{proof}
Limits, 補題 \ref{limits-lemma-modifications-properties} から従う．
\end{proof}

\begin{lemma}
\label{resolve-lemma-equivalence-fibre}
$S, s_i, S_i$ は（\ref{resolve-equation-equivalence}）のものとする．
$f : X \to S$ が $g_i : Y_i \to S_i$ に $F$ のもとで対応するならば，
$X_{s_i} \cong (Y_i)_{s_i}$ は $\kappa(s_i)$ 上のスキームの同型である．
\end{lemma}

\begin{proof}
これは明らかである．
\end{proof}

\begin{lemma}
\label{resolve-lemma-equivalence-sequence-blowups}
$S, s_i, S_i$ は（\ref{resolve-equation-equivalence}）のものとし，
$f : X \to S$ が $g_i : Y_i \to S_i$ に $F$ のもとで対応すると仮定する．
このとき分解
$$
X = Z_m \to Z_{m - 1} \to \ldots \to Z_1 \to Z_0 = S
$$
で，$f$ を分解し，$Z_{j + 1} \to Z_j$ が $Z_j$ の閉点 $z_j$ における
ブローアップであって，この点が $\{s_1, \ldots, s_n\}$ の上にあるものが
存在するための必要十分条件は，各 $i$ に対して分解
$$
Y_i = Z_{i, m_i} \to Z_{i, m_i - 1} \to \ldots \to Z_{i, 1} \to Z_{i, 0} = S_i
$$
で，$g_i$ を分解し，$Z_{i, j + 1} \to Z_{i, j}$ が $Z_{i, j}$ の閉点
$z_{i, j}$ におけるブローアップであって，この点が $s_i$ の上にあるものが
存在することである．
\end{lemma}

\begin{proof}
まずブローアップの列
$Z_m \to Z_{m - 1} \to \ldots \to Z_1 \to Z_0 = S$
から始める．最初の射 $Z_1 \to S$ は，$s_i$ の一つ，例えば $s_1$ を
ブローアップして得られる．$F$ を $Z_1 \to S$ に適用すると，
% P07-ERR-1103: frozen source has a malformed/redundant clause describing the first base-changed blowup.
$Z_{1, 1} \to S_1$ が $s_1$ におけるブローアップで，その中心が $s_1$ であり，
それ以外では $Z_{i, 0} = S_i$（$i > 1$）であることがわかる．
次の段階では，$s_i$（$i \geq 2$）の一つを $Z_1$ 上でブローアップするか，
または閉点 $z_1$ で，$Z_1 \to S$ の $s_1$ 上のファイバーに属するものを選ぶ．
前者の場合に何をすべきかは明らかであり，後者の場合には
$(Z_1)_{s_1} \cong (Z_{1, 1})_{s_1}$
（補題 \ref{resolve-lemma-equivalence-fibre}）を用いて，閉点
$z_{1, 1} \in Z_{1, 1}$ で $z_1$ に対応するものを得る．つぎに $Z_{1, 2} \to Z_{1, 1}$ を
$z_{1, 1}$ におけるブローアップとする．この操作を続けることで，
各 $g_i$ の分解を構成する．

\medskip\noindent
逆に，ブローアップの列
$Z_{i, m_i} \to Z_{i, m_i - 1} \to \ldots \to Z_{i, 1} \to Z_{i, 0} = S_i$
が与えられたとき，まったく同じ方法で $S$ のブローアップの列を構成する．
\end{proof}

\noindent
補題 \ref{resolve-lemma-equivalence-sequence-blowups} の
正規化ブローアップ版は次のとおりである．

\begin{lemma}
\label{resolve-lemma-equivalence-sequence-normalized-blowups}
$S, s_i, S_i$ は（\ref{resolve-equation-equivalence}）のものとし，
$f : X \to S$ が $g_i : Y_i \to S_i$ に $F$ のもとで対応すると仮定する．
$S$ の任意の準コンパクト開部分が有限個の既約成分をもつと仮定する．
このとき分解
$$
X = Z_m \to Z_{m - 1} \to \ldots \to Z_1 \to Z_0 = S
$$
で，$f$ を分解し，$Z_{j + 1} \to Z_j$ が $Z_j$ の閉点 $z_j$ における
% P07-ERR-0389: frozen source uses undefined x_i instead of the distinguished s_i.
正規化ブローアップであって，この点が $\{x_1, \ldots, x_n\}$ の上にあるものが
存在するための必要十分条件は，各 $i$ に対して分解
$$
Y_i = Z_{i, m_i} \to Z_{i, m_i - 1} \to \ldots \to Z_{i, 1} \to Z_{i, 0} = S_i
$$
で，$g_i$ を分解し，$Z_{i, j + 1} \to Z_{i, j}$ が $Z_{i, j}$ の閉点
$z_{i, j}$ における正規化ブローアップであって，この点が $s_i$ の上にあるものが
存在することである．
\end{lemma}

\begin{proof}
$S$ に関する仮定は，正規化するスキームが局所的に有限個の既約成分を
もつことを逐次保証し，主張に意味をもたせるために用いられる．
この点を踏まえれば，補題は補題
\ref{resolve-lemma-equivalence-sequence-blowups} の証明とまったく同じ議論から従う．
\end{proof}








\section{消滅}
\label{resolve-section-vanishing}

\noindent
本節では，しばしば次の設定のもとで議論する．改変とは整スキームの間の
固有双有理射であることを思い出そう
（Morphisms, 定義 \ref{morphisms-definition-modification}）．

\begin{situation}
\label{resolve-situation-vanishing}
% P07-ERR-0390: frozen source omits the finite verb after the displayed triple.
ここで $(A, \mathfrak m, \kappa)$ を次元 $2$ の局所 Noether 正規整域とする．
$s$ を $S = \Spec(A)$ の閉点とし，$U = S \setminus \{s\}$ とおく．
$f : X \to S$ を改変とする．$C_1, \ldots, C_r$ で，特殊ファイバー
$X_s$，すなわち $f$ の特殊ファイバーの既約成分を表す．
\end{situation}

\noindent
Varieties, 補題
\ref{varieties-lemma-modification-normal-iso-over-codimension-1}
により，射 $f$ は同型 $f^{-1}(U) \to U$ を定める．特殊ファイバー
$X_s$ は $\Spec(\kappa)$ 上固有であり，Varieties, 補題
\ref{varieties-lemma-dimension-fibre-in-higher-codimension} により
次元は高々 $1$ である．Stein 分解（More on Morphisms, 補題
\ref{more-morphisms-lemma-geometrically-connected-fibres-towards-normal}）
により $f_*\mathcal{O}_X = \mathcal{O}_S$ であり，特殊ファイバー
$X_s$ は $\kappa$ 上幾何学的連結である．$X_s$ の次元が $0$ ならば，
$f$ は有限であり（More on Morphisms, 補題
\ref{more-morphisms-lemma-proper-finite-fibre-finite-in-neighbourhood}），
したがって同型である
（Morphisms, 補題 \ref{morphisms-lemma-finite-birational-over-normal}）．
この自明な場合は除外することにし，$\dim(C_i) = 1$ が
$i = 1, \ldots, r$ に対して成り立つと結論する．

\begin{lemma}
\label{resolve-lemma-dominate-by-scheme-modification}
状況 \ref{resolve-situation-vanishing} において，$U$-許容ブローアップ
$X' \to S$ で $X$ を支配するものが存在する．
\end{lemma}

\begin{proof}
これは More on Flatness, 補題
\ref{flat-lemma-dominate-modification-by-blowup} の特別な場合である．
\end{proof}

\begin{lemma}
\label{resolve-lemma-nice-meromorphic-function}
状況 \ref{resolve-situation-vanishing} において，非零元 $f \in \mathfrak m$ で，
すべての $i = 1, \ldots, r$ に対して次が存在するものがある：
\begin{enumerate}
\item 閉点 $x_i \in C_i$ で，$x_i \not \in C_j$ が $j \not = i$ のとき成り立つもの；
\item $f = g_i f_i$ という $f$ の $\mathcal{O}_{X, x_i}$ における分解で，
$g_i \in \mathfrak m_{x_i}$ が $\mathcal{O}_{C_i, x_i}$ の非零元に写るもの．
\end{enumerate}
\end{lemma}

\begin{proof}
状況 \ref{resolve-situation-vanishing} の後の観察を断りなく用いる．
$x_i \in C_i$ を，$C_j$ のうち $j \not = i$ を満たすどのものにも
属さない閉点としてとる．
$g_i \in \mathfrak m_{x_i}$ を，$\mathcal{O}_{C_i, x_i}$ の非零元に写るものとしてとる．
$A$ の分数体は
% P07-ERR-0392: frozen source uses undefined X_i rather than X in this local ring.
$\mathcal{O}_{X_i, x_i}$ の分数体なので，$g_i = a_i/b_i$ と書けるような
$a_i, b_i \in A$ が存在する．$f = \prod a_i$ ととる．
\end{proof}

\begin{lemma}
\label{resolve-lemma-nontrivial-normal-bundle}
状況 \ref{resolve-situation-vanishing} において $X$ は正規であると仮定する．
$Z \subset X$ を空でない有効 Cartier 因子で，集合論的に
$Z \subset X_s$ となるものとする．このとき $Z$ の余法層は自明でない．
より正確には，ある $i$ が存在して $C_i \subset Z$ かつ
$\deg(\mathcal{C}_{Z/X}|_{C_i}) > 0$ となる．
\end{lemma}

\begin{proof}
状況 \ref{resolve-situation-vanishing} の後の観察を断りなく用いる．
$f$ を補題 \ref{resolve-lemma-nice-meromorphic-function} の関数とする．
$\xi_i \in C_i$ を一般点とする．$\mathcal{O}_i$ を $X$ の $\xi_i$ における
局所環とする．このとき $\mathcal{O}_i$ は離散付値環である．
$e_i$ を $f$ の $\mathcal{O}_i$ における付値とすると，$e_i > 0$ である．
$h_i \in \mathcal{O}_i$ を $Z$ の局所方程式とし，$d_i$ をその付値とする．
このとき $d_i \geq 0$ である．$i$ を $d_i/e_i$ が最大となるものとして一つ選んで
固定する（このとき $d_i > 0$ である．これは $Z$ が空でないことから従う）．
$f$ を $f^{d_i}$ で，
$Z$ を $e_iZ$ で置き換える．これは，関係式
$\mathcal{O}_X(e_i Z) = \mathcal{O}_X(Z)^{\otimes e_i}$，
余法層と $\mathcal{O}_X(Z)$ の関係（Divisors, 補題
\ref{divisors-lemma-invertible-sheaf-sum-effective-Cartier-divisors}
および \ref{divisors-lemma-conormal-effective-Cartier-divisor}），ならびに
次数が $e_i$ 倍されること（Varieties, 補題
\ref{varieties-lemma-degree-tensor-product}）により許される．
% P07-ERR-0393: Japanese closes the parenthetical citation missing in the frozen source.
$\mathcal{I}$ を $Z$ のイデアル層とすると，
$\mathcal{C}_{Z/X} = \mathcal{I}|_Z$ である．$\overline{f}$ を $f$ の
$\Gamma(Z, \mathcal{O}_Z)$ における像として考える．
上の選択により，$\overline{f}$ は $Z$ の各既約成分の一般点で消える
（これらはすべて $C_j$ の一般点である．実際，$Z$ は特殊ファイバーに含まれる）．
一方，Divisors, 補題
\ref{divisors-lemma-normal-effective-Cartier-divisor-S1} により
$Z$ は $(S_1)$ である．したがってスキーム $Z$ は埋め込まれた随伴点をもたず，
$\overline{f} = 0$ と結論する（Divisors, 補題
\ref{divisors-lemma-S1-no-embedded} および
\ref{divisors-lemma-restriction-injective-open-contains-weakly-ass}）．
よって $f$ は $\mathcal{I}$ の大域切断であり，構成により
$\mathcal{I}_{\xi_i}$ を生成する．したがって，$s_i$ を $f$ の
$\Gamma(C_i, \mathcal{I}|_{C_i})$ における像とすると，これは非零である．
しかし $f$ の選び方により，$s_i$ は $x_i$ で零点をもつ．
ゆえに Varieties, 補題 \ref{varieties-lemma-check-invertible-sheaf-trivial}
により $\mathcal{I}|_{C_i}$ の次数は $> 0$ である．
\end{proof}

\begin{lemma}
\label{resolve-lemma-H1-injective}
状況 \ref{resolve-situation-vanishing} において，$X$ は正規で $A$ は Nagata と仮定する．
写像
$$
H^1(X, \mathcal{O}_X) \longrightarrow H^1(f^{-1}(U), \mathcal{O}_X)
$$
は単射である．
\end{lemma}

\begin{proof}
$0 \to \mathcal{O}_X \to \mathcal{E} \to \mathcal{O}_X \to 0$ を，
非自明な元 $\xi$，すなわち $H^1(X, \mathcal{O}_X)$ の元に対応する拡大とする
（Cohomology, 補題 \ref{cohomology-lemma-h1-extensions}）．
$\pi : P = \mathbf{P}(\mathcal{E}) \to X$ を $\mathcal{E}$ に付随する射影束とする．
全射 $\mathcal{E} \to \mathcal{O}_X$ は切断 $\sigma : X \to P$ を定め，
その余法層は $\mathcal{O}_X$ と同型である（Divisors, 補題
\ref{divisors-lemma-conormal-sheaf-section-projective-bundle}）．
$\xi$ の $f^{-1}(U)$ への制限が自明ならば，写像
$\mathcal{E}|_{f^{-1}(U)} \to \mathcal{O}_{f^{-1}(U)}$ で，単射
$\mathcal{O}_X \to \mathcal{E}$ を分裂させるものを得る．
これは第二の切断 $\sigma' : f^{-1}(U) \to P$ を定め，$\sigma$ と交わらない．
$\xi$ は非自明なので，
$\sigma'$ は $X$ 全体へ延長されて $\sigma$ と交わらないままでいることはできない．
$X' \subset P$ を $\sigma'$ のスキーム論的像とする
（Morphisms, 定義 \ref{morphisms-definition-scheme-theoretic-image}）．
図式は
$$
\xymatrix{
& X' \ar[rd]_g \ar[r] & P \ar[d]_\pi \\
f^{-1}(U) \ar[ru]_{\sigma'} \ar[rr] & & X \ar@/_/[u]_\sigma
}
$$
である．射 $P \setminus \sigma(X) \to X$ はアフィンである．
$X' \cap \sigma(X) = \emptyset$ ならば，$X' \to X$ はアフィンかつ固有であり，
したがって有限（Morphisms, 補題 \ref{morphisms-lemma-finite-proper}），
さらに同型である（$X$ は正規なので，Morphisms, 補題
\ref{morphisms-lemma-finite-birational-over-normal} を参照）．
これは上で述べたことに反する．

\medskip\noindent
$X^\nu$ を $X'$ の正規化とする．$A$ は Nagata なので，
$X^\nu \to X'$ は有限である（Morphisms, 補題
\ref{morphisms-lemma-nagata-normalization} および
\ref{morphisms-lemma-ubiquity-nagata}）．$Z \subset X^\nu$ を
有効 Cartier 因子 $\sigma(X) \subset P$ の引き戻しとする．
上より $Z$ は空でなく，$X^\nu \to S$ の閉ファイバーに含まれる．
$P \to X$ は滑らかなので，$\sigma(X)$ は有効 Cartier 因子である
（Divisors, 補題 \ref{divisors-lemma-section-smooth-regular-immersion}）．
したがって $Z \subset X^\nu$ も有効 Cartier 因子である．
$\sigma(X)$ の $P$ における余法層は $\mathcal{O}_X$ なので，
$Z$ の $X^\nu$ における余法層（これは先験的に可逆である）は，
Morphisms, 補題 \ref{morphisms-lemma-conormal-functorial-flat} により
$\mathcal{O}_Z$ である．これは補題
\ref{resolve-lemma-nontrivial-normal-bundle} に反し，証明は完了する．
\end{proof}

\begin{lemma}
\label{resolve-lemma-R1-injective}
状況 \ref{resolve-situation-vanishing} において，$X$ は正規で $A$ は Nagata と仮定する．
このとき
$$
\Hom_{D(A)}(\kappa[-1], Rf_*\mathcal{O}_X)
$$
は零である．ここでは $D(A) = D_\QCoh(\mathcal{O}_S)$ により
$Rf_*\mathcal{O}_X$ を $D(A)$ の対象とみなしている．
\end{lemma}

\begin{proof}
$Rf_*$ と $Lf^*$ の随伴性により，このような写像は写像
$\alpha : Lf^*\kappa[-1] \to \mathcal{O}_X$ と同じものである．
$$
H^i(Lf^*\kappa[-1]) =
\left\{
\begin{matrix}
0 & \text{if} & i > 1 \\
\mathcal{O}_{X_s} & \text{if} & i = 1 \\
\text{ある }\mathcal{O}_{X_s}\text{-加群} & \text{if} & i \leq 0
\end{matrix}
\right.
$$
であることに注意する．$\Hom(H^0(Lf^*\kappa[-1]), \mathcal{O}_X) = 0$ である．
実際，$\mathcal{O}_X$ は捩れなしであり，
$\Ext$ に対するスペクトル系列
（Cohomology on Sites, 例
\ref{sites-cohomology-example-hom-complex-into-sheaf}）から
$\Hom_{D(\mathcal{O}_X)}(Lf^*\kappa[-1], \mathcal{O}_X)$ は
$\Ext^1_{\mathcal{O}_X}(\mathcal{O}_{X_s}, \mathcal{O}_X)$ に等しい．
したがって $\alpha : Lf^*\kappa[-1] \to \mathcal{O}_X$ は拡大
$$
0 \to \mathcal{O}_X \to \mathcal{E} \to \mathcal{O}_{X_s} \to 0
$$
によって与えられる．補題 \ref{resolve-lemma-H1-injective} により，この拡大を
全射 $\mathcal{O}_X \to \mathcal{O}_{X_s}$ で引き戻したものは零である
（この引き戻しは $f^{-1}(U)$ 上で明らかに分裂するからである）．
したがって $1 \in \mathcal{O}_{X_s}$ は大域切断 $s$ に持ち上がり，
これは $\mathcal{E}$ の切断である．$s$ にイデアル層 $\mathcal{I}$ を掛ける．
ここでこれは $X_s$ のイデアル層であり，その結果
$\mathcal{O}_X$-加群写像 $c_s : \mathcal{I} \to \mathcal{O}_X$ を得る．
$f_*$ を適用すると，$A$-線形写像 $f_*c_s : \mathfrak m \to A$ を得る．
$A$ は Noether 正規局所整域なので，この写像はある元 $a \in A$ による
乗法で与えられる．$s$ を $s -  a$ に取り替えると，$s$ は
$\mathcal{I}$ で零化され，望むとおり拡大は自明となる．
\end{proof}

\begin{remark}
\label{resolve-remark-dualizing-setup}
$X$ を次元 $2$ の整 Noether 正規スキームとする．この場合，次は同値である：
\begin{enumerate}
\item $X$ は双対化複体 $\omega_X^\bullet$ をもつ；
\item ある連接 $\mathcal{O}_X$-加群 $\omega_X$ が存在して，
$\omega_X[n]$ が双対化複体となる．ここで $n$ は任意の整数でよい．
\end{enumerate}
これは $X$ が Cohen--Macaulay であること
（Properties, 補題 \ref{properties-lemma-normal-dimension-2-Cohen-Macaulay}）と
Duality for Schemes, 補題 \ref{duality-lemma-dualizing-module-CM-scheme} から従う．
この状況では，Duality for Schemes, 節
\ref{duality-section-dualizing-module} に従い，
$\omega_X$ を{\it 双対化加群}と呼ぶ．特に $A$ が次元 $2$ の
Noether 正規局所整域であるとき，上の条件が成り立つならば
{\it $A$ は双対化加群 $\omega_A$ をもつ}という．この場合，
$X \to \Spec(A)$ が正規改変ならば，$X$ も双対化加群をもつ．
Duality for Schemes, 例 \ref{duality-example-proper-over-local} を参照せよ．
この状況では常に $\omega_X$ で，$\omega_A$ に関して正規化された
双対化加群，すなわち $\omega_X[2]$ が $\omega_A[2]$ に関して
相対的に正規化された双対化複体となるものを表す．
Duality for Schemes, 節 \ref{duality-section-glue} を参照せよ．
\end{remark}

\noindent
次の命題の Grauert--Riemenschneider 消滅は，補題
\ref{resolve-lemma-R1-injective} と双対性の一般論から形式的に従う．

\begin{proposition}[Grauert-Riemenschneider]
\label{resolve-proposition-Grauert-Riemenschneider}
状況 \ref{resolve-situation-vanishing} において，次を仮定する：
\begin{enumerate}
\item $X$ は正規スキームである；
\item $A$ は Nagata であり，双対化複体 $\omega_A^\bullet$ をもつ．
\end{enumerate}
$\omega_X$ を $X$ の双対化加群とする（注 \ref{resolve-remark-dualizing-setup}）．
このとき $R^1f_*\omega_X = 0$ である．
\end{proposition}

\begin{proof}
この証明では同一視 $D(A) = D_\QCoh(\mathcal{O}_S)$ を用いて，
準連接 $\mathcal{O}_S$-加群を $A$-加群と同一視する．さらに
$\omega_A^\bullet$ は正規化されていると仮定してよい．
Dualizing Complexes, 節 \ref{dualizing-section-dualizing-local} を参照せよ．
$X$ は次元 $2$ の Noether 正規スキームなので Cohen--Macaulay である
（Properties, 補題
\ref{properties-lemma-normal-dimension-2-Cohen-Macaulay}）．
したがって $\omega_X^\bullet = \omega_X[2]$ である
（Duality for Schemes, 補題
\ref{duality-lemma-dualizing-module-CM-scheme} および Duality for Schemes, 例
\ref{duality-example-proper-over-local} における正規化）．
命題が偽ならば，非零写像 $R^1f_*\omega_X \to \kappa$ をとれる．
言い換えると，非零写像
$\alpha : Rf_*\omega_X^\bullet \to \kappa[1]$ を得る．
$R\Hom_A(-, \omega_A^\bullet)$ を適用すると，非零写像
$$
\beta : \kappa[-1] \longrightarrow Rf_*\mathcal{O}_X
$$
を得るが，これは補題 \ref{resolve-lemma-R1-injective} に反する．
$R\Hom_A(-, \omega_A^\bullet)$ が述べたとおりに作用することを見るには，
まず
$$
R\Hom_A(\kappa[1], \omega_A^\bullet) =
R\Hom_A(\kappa, \omega_A^\bullet)[-1] =
\kappa[-1]
$$
に注意する．これは $\omega_A^\bullet$ が正規化されていることによる．また
$$
R\Hom_A(Rf_*\omega_X^\bullet, \omega_A^\bullet) =
Rf_*R\SheafHom_{\mathcal{O}_X}(\omega_X^\bullet, \omega_X^\bullet) =
Rf_*\mathcal{O}_X
$$
が成り立つ．
% P07-ERR-1104: frozen source's “The first equality by ...” clause lacks a finite predicate.
第一の等号は Duality for Schemes, 例
\ref{duality-example-iso-on-RSheafHom-noetherian} と，構成により
$\omega_X^\bullet = f^!\omega_A^\bullet$ であることから従い，
第二の等号は $\omega_X^\bullet$ が $X$ の双対化複体であることから従う
（これは Duality for Schemes, 補題
\ref{duality-lemma-shriek-dualizing} に遡る）．
\end{proof}





\section{有界性}
\label{resolve-section-bounded}

\noindent
本節では，$2$ 次元スキームについて有理特異点の場合への帰着に至る議論を始める．

\begin{lemma}
\label{resolve-lemma-exact-sequence}
$(A, \mathfrak m, \kappa)$ を次元 $2$ の Noether 正規局所整域とする．
可換図式
$$
\xymatrix{
X' \ar[rd]_{f'} \ar[rr]_g & & X \ar[ld]^f \\
& \Spec(A)
}
$$
を考える．ここで $f$ と $f'$ は状況 \ref{resolve-situation-vanishing} における改変であり，
$X$ は正規である．このとき短完全列
$$
0 \to H^1(X, \mathcal{O}_X) \to H^1(X', \mathcal{O}_{X'}) \to
H^0(X, R^1g_*\mathcal{O}_{X'}) \to 0
$$
を得る．さらに $\dim(\text{Supp}(R^1g_*\mathcal{O}_{X'})) = 0$ であり，
$R^1g_*\mathcal{O}_{X'}$ は大域切断で生成される．
\end{lemma}

\begin{proof}
状況 \ref{resolve-situation-vanishing} の後の観察を断りなく用いる．$X$ は正規で，
$g$ は支配的かつ双有理なので，$g_*\mathcal{O}_{X'} = \mathcal{O}_X$ である．
例えば More on Morphisms, 補題
\ref{more-morphisms-lemma-geometrically-connected-fibres-towards-normal} を参照せよ．
$g$ のファイバーの次元は $\leq 1$ なので，$R^pg_*\mathcal{O}_{X'} = 0$ が
$p > 1$ に対して成り立つ．例えば Cohomology of Schemes, 補題
\ref{coherent-lemma-higher-direct-images-zero-above-dimension-fibre} を参照せよ．
% P07-ERR-0394: frozen source refers to undefined g' rather than the diagram's g.
$R^1g_*\mathcal{O}_{X'}$ の台は，$X$ の点のうち，$g'$ のファイバーの次元が
$\geq 1$ となるものの集合に含まれる．したがってこれは，既約成分
$C' \subset X'_s$ のうち $X_s$ の点に写るものの像の集合に含まれ，
その集合は閉点の有限集合である（$X'_s \to X_s$ は固有な $1$ 次元スキームの
$\kappa$ 上の射であることを思い出せ）．
よって Cohomology of Schemes, 補題
\ref{coherent-lemma-coherent-support-dimension-0} により，
$R^1g_*\mathcal{O}_{X'}$ は大域的に生成される．射 $f : X \to S$ と
上の参照結果を用いると，$H^p(X, \mathcal{F}) = 0$ が $p > 1$ に対し，
任意の連接 $\mathcal{O}_X$-加群 $\mathcal{F}$ について成り立つことが分かる．
したがって補題の短完全列は，$g$ と $\mathcal{O}_{X'}$ に対する Leray
スペクトル系列から従う．Cohomology, 補題 \ref{cohomology-lemma-Leray} を参照せよ．
\end{proof}

\begin{lemma}
\label{resolve-lemma-bound-a-torsion}
$(A, \mathfrak m, \kappa)$ を次元 $2$ の Nagata 正規局所整域とする．
$a \in A$ を非零とする．ある整数 $N$ が存在して，任意の改変
$f : X \to \Spec(A)$ で $X$ が正規ならば，$A$-加群
$$
M_{X, a} = \Coker(A \longrightarrow H^0(Z, \mathcal{O}_Z))
$$
を考える．ここで $Z \subset X$ は $a$ により切り出される閉部分スキームであり，
この加群の長さは $N$ で抑えられる．
\end{lemma}

\begin{proof}
短完全列
$
0 \to \mathcal{O}_X \xrightarrow{a} \mathcal{O}_X \to \mathcal{O}_Z \to 0
$
より
\begin{equation}
\label{resolve-equation-a-torsion}
M_{X, a} = H^1(X, \mathcal{O}_X)[a]
\end{equation}
を得る．ここで $N[a] = \{n \in N \mid an = 0\}$ とし，$A$-加群を $N$ で表した．
したがって $a$ が $b$ を割れば，
$M_{X, a} \subset M_{X, b}$ である．ある $c \in A$ について加群
$M_{X, c}$ の長さが有界であると仮定する．このとき各 $X$ に対して完全列
$$
0 \to M_{X, c} \to M_{X, c^2} \to M_{X, c}
$$
があり，第二の矢印は $c$ による乗法で与えられる．したがって
$M_{X, c^2}$ の長さも有界である．ゆえに，$c \in A$ について補題が成り立ち，
$a$ が $c^n$ をある $n > 0$ に対して割るようなものを見いだせば十分である．
More on Algebra, 補題 \ref{more-algebra-lemma-divides-radical} により，
$A/(a)$ は被約環であると仮定してよい．

\medskip\noindent
$A/(a)$ は被約であると仮定する．$A/(a) \subset B$ を，その全商環における
$A/(a)$ の正規化とする．$A$ は Nagata なので，$\Coker(A \to B)$ は有限である．
この有限加群の長さが求める上界であると主張する．実際，補題における
$f : X \to \Spec(A)$ を考え，$Z' \subset Z$ を $Z \cap f^{-1}(U)$ の
スキーム論的閉包とする．すると例えば Varieties, 補題
\ref{varieties-lemma-finite-in-codim-1} により，$Z' \to \Spec(A/(a))$ は有限である．
したがって $Z' = \Spec(B')$ であり，$A/(a) \subset B' \subset B$ である．一方，写像
$$
H^0(Z, \mathcal{O}_Z) \to H^0(Z', \mathcal{O}_{Z'})
$$
は単射であると主張する．実際，$s \in H^0(Z, \mathcal{O}_Z)$ が核に属すれば，
$s$ の $f^{-1}(U) \cap Z$ への制限は零である．したがって $s$ の
$H^1(X, \mathcal{O}_X)$ における像は
$H^1(f^{-1}(U), \mathcal{O}_X)$ において消える．補題
\ref{resolve-lemma-H1-injective} により，$s$ はある元 $\tilde s$（$A$ のもの）から来る．
しかし仮定により $\tilde s$ は $B'$ で零に写るので，$s = 0$ である．
以上を合わせると，$M_{X, a}$ は $B'/A$ の部分商であることが分かる．
すなわち $B'$ のすべての元が $\mathcal{O}_Z$ の大域切断へ延長されるとは
限らないが，いずれにせよ $M_{X, a}$ の長さは $B/A$ の長さで抑えられる．
\end{proof}

\noindent
場合によっては，特異点解消は有理特異点の場合へ帰着される．

\begin{definition}
\label{resolve-definition-reduce-to-rational}
$(A, \mathfrak m, \kappa)$ を次元 $2$ の Nagata 正規局所整域とする．
\begin{enumerate}
\item $A$ が {\it 有理特異点を定める}というのは，任意の正規改変
$X \to \Spec(A)$ に対して $H^1(X, \mathcal{O}_X) = 0$ となるときである．
\item {\it $A$ に対して有理特異点への帰着が可能である}というのは，$A$-加群
$$
H^1(X, \mathcal{O}_X)
$$
の長さが，すべての改変 $X \to \Spec(A)$ で $X$ が正規であるものにわたって
有界となるときである．
\end{enumerate}
\end{definition}

\noindent
(2) の用語の意味は補題 \ref{resolve-lemma-reduce-to-rational} で説明される．
次の補題は大まかにいえば，$\Spec(A)$ の改変の局所環は，$A$ が
有理特異点を定めるならば，それ自身も有理特異点を定めることを述べている．

\begin{lemma}
\label{resolve-lemma-rational-propagates}
$(A, \mathfrak m, \kappa)$ を，有理特異点を定める次元 $2$ の Nagata
正規局所整域とする．$A \subset B$ を，同じ分数体をもつ整域の
本質的有限型な局所拡大で，$\dim(B) = 2$ かつ $B$ が正規であるものとする．
このとき $B$ は有理特異点を定める．
\end{lemma}

\begin{proof}
有限型な $A$-代数 $C$ を，$B = C_\mathfrak q$ がある素イデアル
$\mathfrak q \subset C$ に対して成り立つように選ぶ．$C$ を $C$ の $B$ における像で置き換えることで，
$C$ は $A$ の分数体と同じ分数体をもつ整域であると仮定してよい．
閉埋め込み $\Spec(C) \to \mathbf{A}^n_A$ を選び，$\mathbf{P}^n_A$ における
閉包をとると，$B$ は $\mathcal{O}_{X, x}$ と同型であることが分かる．ここで
$x \in X$ は，ある射影的改変 $X \to \Spec(A)$ の閉点である．
（Morphisms, 補題 \ref{morphisms-lemma-dimension-formula} により
$\kappa(x)$ は $\kappa$ 上有限であり，続いて Morphisms, 補題
\ref{morphisms-lemma-algebraic-residue-field-extension-closed-point-fibre}
により $x$ は閉点である．）
$\nu : X^\nu \to X$ を正規化とする．$A$ は Nagata なので，射 $\nu$ は有限である
（Morphisms, 補題 \ref{morphisms-lemma-nagata-normalization}）．
したがって More on Morphisms, 補題
\ref{more-morphisms-lemma-category-projective} により，$X^\nu$ は $A$ 上射影的である．
$B = \mathcal{O}_{X, x}$ は正規なので，
$\mathcal{O}_{X, x} = (\nu_*\mathcal{O}_{X^\nu})_x$ である．
ゆえに点 $x^\nu \in X^\nu$ で $x$ の上にあるものは一意であり，
$\mathcal{O}_{X^\nu, x^\nu} = \mathcal{O}_{X, x}$ となる．したがって
$X$ は正規かつ $A$ 上射影的であると仮定してよい．
$Y \to \Spec(\mathcal{O}_{X, x}) = \Spec(B)$ を改変とし，$Y$ は正規であるとする．
$H^1(Y, \mathcal{O}_Y) = 0$ を示さなければならない．Limits, 補題
\ref{limits-lemma-modifications} により，スキームの射 $g : X' \to X$ で，
$X \setminus \{x\}$ 上で同型であり，かつ
$X' \times_X \Spec(\mathcal{O}_{X, x})$ が $Y$ と同型となるものを見いだせる．Limits, 補題
\ref{limits-lemma-modifications-properties} により $g$ は固有なので，改変である．
% P07-ERR-0395: frozen source twice says “at a point of x'”; Japanese uses the intended construction.
$X'$ の点 $x'$ における局所環は，$X$ の $g(x')$ における局所環と，
$g(x') \not = x$ の場合には同型である．一方 $g(x') = x$ の場合には，
$X'$ の $x'$ における局所環は，対応する点における $Y$ の局所環と同型である．
したがって $X'$ は正規である．実際，$X$ と $Y$ はともに正規である．
$H^1(X', \mathcal{O}_{X'}) = 0$ は $A$ に関する仮定から従う．
補題 \ref{resolve-lemma-exact-sequence} により $R^1g_*\mathcal{O}_{X'} = 0$ であり，
これは明らかに，望むとおり $H^1(Y, \mathcal{O}_Y) = 0$ を意味する．
\end{proof}

\begin{lemma}
\label{resolve-lemma-reduce-to-rational}
$(A, \mathfrak m, \kappa)$ を次元 $2$ の Nagata 正規局所整域とする．
$A$ に対して有理特異点への帰着が可能ならば，閉点における正規化ブローアップの有限列
$$
X = X_n \to X_{n - 1} \to \ldots \to X_1 \to X_0 = \Spec(A)
$$
で，任意の閉点 $x \in X$ に対して局所環 $\mathcal{O}_{X, x}$ が
有理特異点を定めるものが存在する．特に $X \to \Spec(A)$ は改変であり，
$X$ は $A$ 上射影的な正規スキームである．
\end{lemma}

\begin{proof}
$X \to \Spec(A)$ を改変で，$X$ が正規かつ $H^1(X, \mathcal{O}_X)$ の
長さを最大にするものとして選ぶ．補題 \ref{resolve-lemma-exact-sequence} により，
任意のさらなる改変 $g : X' \to X$ で $X'$ が正規であるものに対して
$R^1g_*\mathcal{O}_{X'} = 0$ かつ
$H^1(X, \mathcal{O}_X) = H^1(X', \mathcal{O}_{X'})$ である．

\medskip\noindent
$x \in X$ を閉点とする．$\mathcal{O}_{X, x}$ が有理特異点を定めることを示す．
$Y \to \Spec(\mathcal{O}_{X, x})$ を改変とし，$Y$ は正規であるとする．
$H^1(Y, \mathcal{O}_Y) = 0$ を示さなければならない．Limits, 補題
\ref{limits-lemma-modifications} により，スキームの射 $g : X' \to X$ で，
$X \setminus \{x\}$ 上で同型であり，かつ
$X' \times_X \Spec(\mathcal{O}_{X, x})$ が $Y$ と同型となるものを見いだせる．Limits, 補題
\ref{limits-lemma-modifications-properties} により $g$ は固有なので，改変である．
% P07-ERR-0395: second copied “point of x'” occurrence; frozen authority remains unchanged.
$X'$ の点 $x'$ における局所環は，$X$ の $g(x')$ における局所環と，
$g(x') \not = x$ の場合には同型である．一方 $g(x') = x$ の場合には，
$X'$ の $x'$ における局所環は，対応する点における $Y$ の局所環と同型である．
したがって $X'$ は正規である．実際，$X$ と $Y$ はともに正規である．
最大性により $R^1g_*\mathcal{O}_{X'} = 0$ である（第一段落を参照）．
これは明らかに，望むとおり $H^1(Y, \mathcal{O}_Y) = 0$ を意味する．

\medskip\noindent
以上により，正規改変 $X$（$\Spec(A)$ のもの）で，閉点における $X$ の局所環が
すべて有理特異点を定めるものを一つ得た．次に，正規化ブローアップの列
$X_n \to \ldots \to X_1 \to \Spec(A)$ を，$X_n$ が $X$ を支配するように選ぶ．
補題 \ref{resolve-lemma-dominate-by-normalized-blowing-up} を参照せよ．
$x' \in X_n$ を $x \in X$ に写る閉点とし，環写像
$\mathcal{O}_{X, x} \to \mathcal{O}_{X_n, x'}$ に補題
\ref{resolve-lemma-rational-propagates} を適用すると，$\mathcal{O}_{X_n, x'}$ が
有理特異点を定めることが分かる．
\end{proof}

\begin{lemma}
\label{resolve-lemma-go-up-separable}
$A \to B$ を，次元 $2$ の Nagata 正規局所整域の間の有限単射な局所環写像とする．
誘導される分数体の拡大は分離的であると仮定する．$A$ に対して有理特異点への
帰着が可能ならば，$B$ に対しても可能である．
\end{lemma}

\begin{proof}
$n$ を分数体拡大 $L/K$ の次数とする．
$\text{Trace}_{L/K} : L \to K$ をトレースとする．拡大は有限分離的なので，
% P07-ERR-0396: frozen trace-pairing formula uses undefined f instead of input h.
トレース対 $(h, g) \mapsto \text{Trace}_{L/K}(fg)$ は $L$ を台とする $K$ 上の
非退化双線形形式である．Fields, 補題
\ref{fields-lemma-separable-trace-pairing} を参照せよ．
$b_1, \ldots, b_n \in B$ を選び，これらが $L$ の $K$ 上の基底をなすようにする．
上より $d = \det(\text{Trace}_{L/K}(b_ib_j)) \in A$ は非零である．

\medskip\noindent
$Y \to \Spec(B)$ を改変で，$Y$ が正規であるものとする．
% P07-ERR-0397: frozen lemma invokes U without defining the punctured open.
$U$-許容ブローアップ $X'$ を $\Spec(A)$ 上で，狭義変換 $Y'$（$Y$ から得られるもの）が
$X'$ 上有限となるように見いだせる．More on Flatness, 補題
\ref{flat-lemma-finite-after-blowing-up} を参照せよ．図式は
$$
\xymatrix{
Y' \ar[d] \ar[r] & Y \ar[r] & \Spec(B) \ar[d] \\
X' \ar[rr] & & \Spec(A)
}
$$
$X'$ と $Y'$ をそれぞれの正規化で置き換えることで，$X'$ と $Y'$ は
$\Spec(A)$ と $\Spec(B)$ の正規改変であると仮定してよい．
このようにして，可換図式
$$
\xymatrix{
Y \ar[d]_\pi \ar[r]_-g & \Spec(B) \ar[d] \\
X \ar[r]^-f & \Spec(A)
}
$$
で，$X$ と $Y$ が $\Spec(A)$ と $\Spec(B)$ の正規改変であり，
$\pi$ が有限であるものが存在する場合へ帰着される．

\medskip\noindent
$L$ の $K$ 上のトレース写像は，$\mathcal{O}_X$-加群の写像
$\text{Trace} : \pi_*\mathcal{O}_Y \to \mathcal{O}_X$ へ延長される．写像
$$
\Phi : \pi_*\mathcal{O}_Y \longrightarrow \mathcal{O}_X^{\oplus n},\quad
s \longmapsto (\text{Trace}(b_1s), \ldots, \text{Trace}(b_ns))
$$
を考える．この写像は単射であり（一般点で単射だからである），また写像
$$
\mathcal{O}_X^{\oplus n} \longrightarrow \pi_*\mathcal{O}_Y,\quad
(s_1, \ldots, s_n) \longmapsto \sum b_i s_i
$$
があって，これと $\Phi$ との合成の行列は $\text{Trace}(b_ib_j)$ である．
したがって $\Phi$ の余核は $d$ で零化される．ゆえに完全列
$$
H^0(X, \Coker(\Phi)) \to H^1(Y, \mathcal{O}_Y) \to
H^1(X, \mathcal{O}_X)^{\oplus n}
$$
を得る．仮定により右辺は有界なので，$d$-捩れが
$H^1(Y, \mathcal{O}_Y)$ において有界であることを示せば十分である．これは補題
\ref{resolve-lemma-bound-a-torsion} と式 (\ref{resolve-equation-a-torsion}) の内容である．
\end{proof}

\begin{lemma}
\label{resolve-lemma-regular-rational}
$A$ を次元 $2$ の Nagata 正則局所環とする．このとき $A$ は有理特異点を定める．
\end{lemma}

\begin{proof}
（$A$ が Nagata であるという仮定はこの証明には不要であるが，有理特異点の概念は
Nagata な $2$ 次元正規局所整域の場合にしか定義していない．）
$X \to \Spec(A)$ を改変で，$X$ が正規であるものとする．補題
\ref{resolve-lemma-dominate-by-blowing-up-in-points} により，$X$ をスキーム $X_n$ により
支配でき，このスキームは列
$$
X_n \to X_{n - 1} \to \ldots \to X_1 \to X_0 = \Spec(A)
$$
の最後の項である．ここで各射は閉点におけるブローアップである．
補題 \ref{resolve-lemma-blowup-regular} により
スキーム $X_i$ は正則であり，特に正規である（Algebra, 補題
\ref{algebra-lemma-regular-normal}）．補題 \ref{resolve-lemma-exact-sequence} により
$H^1(X, \mathcal{O}_X) \subset H^1(X_n, \mathcal{O}_{X_n})$ である．
したがって $H^1(X_n, \mathcal{O}_{X_n}) = 0$ を示せば十分である．
補題 \ref{resolve-lemma-exact-sequence} を再び用いると，
$R^1(X_i \to X_{i - 1})_*\mathcal{O}_{X_i} = 0$ を $i = 1, \ldots, n$ に対して
示せば十分であることが分かる．
これは補題 \ref{resolve-lemma-cohomology-of-blowup} から従う．
\end{proof}

\begin{lemma}
\label{resolve-lemma-bound-dualizing-implies-bound}
$A$ を次元 $2$ の Nagata 正規局所整域で，双対化複体 $\omega_A^\bullet$ を
もつものとする．零でない $d \in A$ が存在して，すべての正規改変
$X \to \Spec(A)$ に対してトレース写像
$$
\Gamma(X, \omega_X) \to \omega_A
$$
の余核が $d$ により零化されるならば，$A$ に対して
有理特異点への帰着が可能である．
\end{lemma}

\begin{proof}
主張にある $X \to \Spec(A)$ に対して $H^1(X, \mathcal{O}_X)$ を抑えなければならない．
$\omega_X$ を，Grauert--Riemenschneider の主張（命題
\ref{resolve-proposition-Grauert-Riemenschneider}）における $X$ の双対化加群とする．
トレース写像は Duality for Schemes, 節 \ref{duality-section-trace} で記述された写像
$Rf_*\omega_X \to \omega_A$ である．Grauert--Riemenschneider により
% P07-ERR-0398: Japanese separates the two clauses run together in the frozen source.
$Rf_*\omega_X = f_*\omega_X$ である．したがってトレース写像は確かに写像
$\Gamma(X, \omega_X) \to \omega_A$ を与える．双対性により
$Rf_*\omega_X = R\Hom_A(Rf_*\mathcal{O}_X, \omega_A)$ である
（ここでは，$\omega_X[2]$ が $X$ 上の双対化複体であり，$\omega_A[2]$ に関して
正規化されていることを用いた．Duality for Schemes, 補題
\ref{duality-lemma-duality-bootstrap}，より直接には節
\ref{duality-section-duality}，さらに直接には例
\ref{duality-example-iso-on-RSheafHom-noetherian} を参照せよ）．
識別三角形
$$
A \to Rf_*\mathcal{O}_X \to R^1f_*\mathcal{O}_X[-1] \to A[1]
$$
は $R\Hom_A(-, \omega_A)$ により短完全列
$$
0 \to f_*\omega_X \to \omega_A \to
\Ext_A^2(R^1f_*\mathcal{O}_X, \omega_A) \to 0
$$
へ移される（また $\Ext_A^i(R^1f_*\mathcal{O}_X, \omega_A) = 0$ が
$i \not = 2$ に対して成り立つ．これは以下の議論からも従う）．
$R^1f_*\mathcal{O}_X$ の台は $\{\mathfrak m\}$ に含まれるので，
局所双対定理により
$$
\Ext_A^2(R^1f_*\mathcal{O}_X, \omega_A) =
\Ext_A^0(R^1f_*\mathcal{O}_X, \omega_A[2]) =
\Hom_A(R^1f_*\mathcal{O}_X, E)
$$
は $R^1f_*\mathcal{O}_X$ の Matlis 双対である（また他の Ext 群は零である）．
Dualizing Complexes, 補題
\ref{dualizing-lemma-special-case-local-duality} を参照せよ．Matlis 双対に内在する
圏同値（Dualizing Complexes, 命題 \ref{dualizing-proposition-matlis}）により，
$R^1f_*\mathcal{O}_X$ が $d$ で零化されなければ，上の $\Ext^2$ も零化されない．
したがって $H^1(X, \mathcal{O}_X)$ は $d$ で零化される．ゆえに求める有界性は
補題 \ref{resolve-lemma-bound-a-torsion} と式 (\ref{resolve-equation-a-torsion}) から従う．
\end{proof}

\begin{lemma}
\label{resolve-lemma-compare-differentials-dualizing}
$p$ を素数とする．$A$ を次元 $2$，標数 $p$ の正則局所環とする．
$A_0 \subset A$ を，$\Omega_{A/A_0}$ が有限階数 $r < \infty$ の自由加群となる
部分環とする．$\omega_A = \Omega^r_{A/A_0}$ とおく．$X \to \Spec(A)$ が
閉点におけるブローアップの列の結果ならば，写像
$$
\varphi_X : (\Omega^r_{X/\Spec(A_0)})^{**} \longrightarrow \omega_X
$$
で，一般点における与えられた同一視を延長するものが存在する．
\end{lemma}

\begin{proof}
$A$ は Gorenstein であることに注意する（Dualizing Complexes, 補題
\ref{dualizing-lemma-regular-gorenstein}）．したがって可逆加群 $\omega_A$ は
実際に双対化加群として働く．さらに，補題における任意の $X$ は可逆な
双対化加群 $\omega_X$ をもつ．実際 $X$ は正則であり（したがって Gorenstein），
$A$ 上固有である．注 \ref{resolve-remark-dualizing-setup} と補題
\ref{resolve-lemma-blowup-regular} を参照せよ．写像
$\varphi_X : (\Omega^r_{X/A_0})^{**} \to \omega_X$ が構成されていると仮定し，
$b : X' \to X$ を閉点におけるブローアップとする．
$\Omega^r_X = (\Omega^r_{X/A_0})^{**}$ および
$\Omega^r_{X'} = (\Omega^r_{X'/A_0})^{**}$ とおく．$\omega_{X'} = b^!(\omega_X)$
なので，写像 $\Omega^r_{X'} \to \omega_{X'}$ は写像
$Rb_*(\Omega^r_{X'}) \to \omega_X$ と同じものである．注
\ref{resolve-remark-dualizing-setup} と Duality for Schemes, 節
\ref{duality-section-duality} の議論を参照せよ．したがって写像
$$
Rb_*(\Omega^r_{X'}) \longrightarrow \Omega^r_X
$$
を構成すれば十分である．層 $\Omega^r_{X'}$ と $\Omega^r_X$ は可逆である．
Divisors, 補題 \ref{divisors-lemma-reflexive-over-regular-dim-2} を参照せよ．
完全列
$$
b^*\Omega_{X/A_0} \to \Omega_{X'/A_0} \to \Omega_{X'/X} \to 0
$$
を考える．局所計算により，$\Omega_{X'/X}$ は例外因子 $E$ 上の可逆加群と
同型であることが分かる．補題 \ref{resolve-lemma-differentials-of-blowup} を参照せよ．
したがって
$$
\Omega^r_{X'} \cong (b^*\Omega^r_X)(E)
\quad\text{または}\quad
\Omega^r_{X'} \cong b^*\Omega^r_X
$$
のいずれかが成り立つ．Divisors, 補題
\ref{divisors-lemma-wedge-product-ses} を参照せよ．
（第二の可能性は標数零では決して生じないが，標数 $p$ では生じうる．）
いずれの場合も補題 \ref{resolve-lemma-cohomology-of-blowup} により，
$R^1b_*(\Omega^r_{X'}) = 0$ かつ $b_*(\Omega^r_{X'}) = \Omega^r_X$ である．
\end{proof}

\begin{lemma}
\label{resolve-lemma-go-up-degree-p}
$p$ を素数とする．$A$ を次元 $2$，標数 $p$ の完備正則局所環とする．
$L/K$ を次数 $p$ の非分離拡大とする．ここで $K$ は $A$ の分数体である．
$B \subset L$ を $A$ の整閉包とする．このとき $B$ に対して
有理特異点への帰着が可能である．
\end{lemma}

\begin{proof}
$A = k[[x, y]]$ である．等式 $L = K[x]/(x^p - f)$ がある $f \in A$ に対して成り立つ．
% P07-ERR-0399: Japanese uses the idiomatic notation declaration missing “by” upstream.
$g \in B$ で $x$ の剰余類を表す．すなわち $g^p = f$ を満たす元である．
Algebra, 補題 \ref{algebra-lemma-derivative-zero-pth-power} により，
$\text{d}f$ は $\Omega_{K/\mathbf{F}_p}$ において非零である．More on Algebra, 補題
\ref{more-algebra-lemma-power-series-ring-subfields} により，部分体
$k^p \subset k' \subset k$ で，$p^e = [k : k'] < \infty$ かつ
$\text{d}f$ が $\Omega_{K/K_0}$ において非零となるものが存在する．ここで
$K_0$ は $A_0 = k'[[x^p, y^p]] \subset A$ の分数体である．このとき
$$
\Omega_{A/A_0} =
A \otimes_k \Omega_{k/k'} \oplus A \text{d}x \oplus A \text{d}y
$$
は階数 $e + 2$ の有限自由加群である．$\omega_A = \Omega^{e + 2}_{A/A_0}$ とおく．
補題 \ref{resolve-lemma-trace-extends} の標準写像
$$
\text{Tr} :
\Omega^{e + 2}_{B/A_0}
\longrightarrow
\Omega^{e + 2}_{A/A_0} = \omega_A
$$
を考える．双対性により，これは写像
$$
c : \Omega^{e + 2}_{B/A_0} \to \omega_B = \Hom_A(B, \omega_A)
$$
を定める．主張：$c$ の余核は $B$ のある非零元で零化される．

\medskip\noindent
$\text{d}f$ は $\Omega_{A/A_0}$ において非零なので，
$\eta_1, \ldots, \eta_{e + 1} \in \Omega_{A/A_0}$ を，
$\theta = \eta_1 \wedge \ldots \wedge \eta_{e + 1} \wedge \text{d}f$ が
$\omega_A = \Omega^{e + 2}_{A/A_0}$ において非零となるように選べる．
主張を証明するため，元 $\omega_i$ を $\Omega^{e + 2}_{B/A_0}$ から
$i = 0, \ldots, p - 1$ の各々について選び，それらが
$\varphi_i \in \omega_B = \Hom_A(B, \omega_A)$ に写り，
$\varphi_i(g^j) = \delta_{ij}\theta$ が $j = 0, \ldots, p - 1$ に対して成り立つものを構成する．
$\{1, g, \ldots, g^{p - 1}\}$ は $L/K$ の基底なので，これで主張が従う．
$\eta = \eta_1 \wedge \ldots \wedge \eta_{e + 1}$ とおき，
$\theta = \eta \wedge \text{d}f$ とする．
$\omega_i = \eta \wedge g^{p - 1 - i}\text{d}g$ とおく．すると構成により
$$
\varphi_i(g^j) = \text{Tr}(g^j \eta \wedge g^{p - 1 - i}\text{d}g) =
\text{Tr}(\eta \wedge g^{p - 1 - i + j}\text{d}g) = \delta_{ij} \theta
$$
である．これは補題 \ref{resolve-lemma-trace-higher} におけるトレース写像の明示的記述による．

\medskip\noindent
$Y \to \Spec(B)$ を正規改変とする．補題 \ref{resolve-lemma-go-up-separable} の証明と
まったく同様にして，$Y$ が改変 $X$（$\Spec(A)$ のもの）上有限である場合へ帰着できる．
補題 \ref{resolve-lemma-dominate-by-blowing-up-in-points} により，さらに
$X \to \Spec(A)$ は閉点におけるブローアップの列の結果であると仮定してよい．図式は
$$
\xymatrix{
Y \ar[d]_\pi \ar[r]_-g & \Spec(B) \ar[d] \\
X \ar[r]^-f & \Spec(A)
}
$$
補題 \ref{resolve-lemma-trace-extends} を $\pi$ に適用でき，
$$
\pi_*(\Omega^{e + 2}_{Y/A_0})
\xrightarrow{\text{Tr}}
(\Omega^{e + 2}_{X/A_0})^{**}
\xrightarrow{\varphi_X}
\omega_X
$$
の第一の矢印を得る．第二の矢印は補題
\ref{resolve-lemma-compare-differentials-dualizing} から来る
（$f$ は閉点におけるブローアップの列だからである）．
有限射 $\pi$ に対する双対性により，これは写像
$$
c_Y : \Omega^{e + 2}_{Y/A_0} \longrightarrow \omega_Y
$$
で，上の写像 $c$ を延長するものに対応する．したがって
$\Gamma(Y, \omega_Y) \to \omega_B$ の像は $c$ の像を含む．
主張により，その余核は $B$ の固定された非零元で零化される．
補題 \ref{resolve-lemma-bound-dualizing-implies-bound} により結論を得る．
\end{proof}






\section{有理特異点}
\label{resolve-section-rational-singularities}

\noindent
本節では，有理特異点から Gorenstein 有理特異点へ帰着する．
\cite{Lipman-rational} および \cite{Mattuck} を参照せよ．

\begin{situation}
\label{resolve-situation-rational}
% P07-ERR-0400: Japanese supplies the finite verb omitted in the frozen situation.
ここで $(A, \mathfrak m, \kappa)$ は，有理特異点を定める次元 $2$ の Nagata
正規局所整域である．$s$ を $S = \Spec(A)$ の閉点とし，
$U = S \setminus \{s\}$ とおく．$f : X \to S$ を改変であって，$X$ が正規であるものとする．
% P07-ERR-0401: Japanese uses an idiomatic notation declaration missing “by” upstream.
$C_1, \ldots, C_r$ で，特殊ファイバー $X_s$（射 $f$ のファイバー）の既約成分を表すことにする．
\end{situation}

\begin{lemma}
\label{resolve-lemma-globally-generated}
% P07-ERR-0402: frozen introductory situation fragment is joined to its declaration.
状況 \ref{resolve-situation-rational} において，$\mathcal{F}$ を準連接
$\mathcal{O}_X$-加群とする．このとき
\begin{enumerate}
\item $H^p(X, \mathcal{F}) = 0$ である（$p \not \in \{0, 1\}$）．
\item $H^1(X, \mathcal{F}) = 0$ である（$\mathcal{F}$ が大域的に生成される場合）．
\end{enumerate}
\end{lemma}

\begin{proof}
(1) は Cohomology of Schemes, 補題
\ref{coherent-lemma-higher-direct-images-zero-above-dimension-fibre} から従う．
$\mathcal{F}$ が大域的に生成されるならば，全射
$\bigoplus_{i \in I} \mathcal{O}_X \to \mathcal{F}$ が存在する．(1) と
コホモロジーの長完全列により，これは $H^1$ 上の全射を誘導する．
$H^1(X, \mathcal{O}_X) = 0$ である（$S$ は有理特異点をもつ）．また
$H^1(X, -)$ は直和と可換である（Cohomology, 補題
\ref{cohomology-lemma-quasi-separated-cohomology-colimit}）．したがって結論を得る．
\end{proof}

\begin{lemma}
\label{resolve-lemma-sections-powers-I-rational}
状況 \ref{resolve-situation-rational} において，$E = X_s$ は有効 Cartier 因子であると仮定する．
$\mathcal{I}$ を $E$ のイデアル層とする．このとき
$H^0(X, \mathcal{I}^n) = \mathfrak m^n$ かつ
$H^1(X, \mathcal{I}^n) = 0$ である．
\end{lemma}

\begin{proof}
$H^0(X, \mathcal{O}_X) = A$ である．状況 \ref{resolve-situation-vanishing} の後の議論を
参照せよ．すると
$\mathfrak m \subset H^0(X, \mathcal{I}) \subset H^0(X, \mathcal{O}_X)$ である．
$X_s \not = \emptyset$ なので，第二の包含は等号ではない．したがって
$H^0(X, \mathcal{I}) = \mathfrak m$ である．
$\mathcal{I}^n = \mathfrak m^n\mathcal{O}_X$ なので，補題
\ref{resolve-lemma-globally-generated} により $H^1(X, \mathcal{I}^n) = 0$ である．

\medskip\noindent
$x_1, \ldots, x_{\mu + 1}$ を $\mathfrak m$ の生成元として選ぶ．これらは
$\mathcal{I}$ の大域切断を定め，それを生成する．したがって短完全列
$$
0 \to \mathcal{F} \to \mathcal{O}_X^{\oplus \mu + 1} \to \mathcal{I} \to 0
$$
を得る．このとき $\mathcal{F}$ は有限局所自由 $\mathcal{O}_X$-加群であり，
その階数は $\mu$ である．また，Constructions, 補題
\ref{constructions-lemma-globally-generated-omega-twist-1} により
$\mathcal{F} \otimes \mathcal{I}$ は大域的に生成される．したがって
$\mathcal{F} \otimes \mathcal{I}^n$ はすべての $n \geq 1$ に対して大域的に
生成される．ゆえに $n \geq 2$ に対して完全列
$$
0 \to \mathcal{F} \otimes \mathcal{I}^{n - 1} \to
(\mathcal{I}^{n - 1})^{\oplus \mu + 1} \to
\mathcal{I}^n \to 0
$$
を考えられる．補題 \ref{resolve-lemma-globally-generated} により
$H^1(X, \mathcal{F} \otimes \mathcal{I}^{n - 1}) = 0$ であることを用いて
コホモロジーの長完全列を適用すると，$H^0(X, \mathcal{I}^n)$ の各元は
$\sum x_i a_i$ の形であることが分かる．ここで
$a_i \in H^0(X, \mathcal{I}^{n - 1})$ である．帰納法により
$H^0(X, \mathcal{I}^n) = \mathfrak m^n$ が従う．
\end{proof}

\begin{lemma}
\label{resolve-lemma-blow-up-normal-rational}
状況 \ref{resolve-situation-rational} において，$\Spec(A)$ の $\mathfrak m$ に沿う
ブローアップは正規である．
\end{lemma}

\begin{proof}
$X' \to \Spec(A)$ をそのブローアップとする．言い換えると
$$
X' = \text{Proj}(A \oplus \mathfrak m \oplus \mathfrak m^2 \oplus \ldots).
$$
である．これは Rees 代数の Proj である．
% P07-ERR-0403: Japanese supplies the demonstrative subject omitted upstream.
特に，$X'$ は整であり，$X' \to \Spec(A)$ は射影的改変であることが分かる．
$X$ を $X'$ の正規化とする．$A$ は Nagata なので，
$\nu : X \to X'$ は有限である（Morphisms, 補題
\ref{morphisms-lemma-nagata-normalization}）．$E' \subset X'$ を例外因子とし，
$E \subset X$ をその逆像とする．$\mathcal{I}' \subset \mathcal{O}_{X'}$ と
$\mathcal{I} \subset \mathcal{O}_X$ をそれぞれのイデアル層とする．
$\mathcal{I}' = \mathcal{O}_{X'}(1)$ であることを思い出そう
（Divisors, 補題 \ref{divisors-lemma-blowing-up-projective}）．
$\mathcal{I} = \nu^*\mathcal{I}'$ であり，$E$ は有効 Cartier 因子である
（Divisors, 補題 \ref{divisors-lemma-pullback-effective-Cartier-defined}）．
$\nu$ が同型であることを示したい．$\nu$ は有限なので，
$\mathcal{O}_{X'} \to \nu_*\mathcal{O}_X$ が同型であることを示せば十分である．
そうでなければ，ある $n \geq 0$ が存在して
$$
H^0(X', (\mathcal{I}')^n) \not =
H^0(X', (\nu_*\mathcal{O}_X) \otimes (\mathcal{I}')^n)
$$
となる．例えば，準連接 $\mathcal{O}_{X'}$-加群はその随伴次数付き加群から
復元できるからである．Properties, 補題
\ref{properties-lemma-ample-quasi-coherent} を参照せよ．射影公式により
$$
H^0(X', (\nu_*\mathcal{O}_X) \otimes (\mathcal{I}')^n) =
H^0(X, \nu^*(\mathcal{I}')^n) =
H^0(X, \mathcal{I}^n) = \mathfrak m^n
$$
である．最後の等号は補題 \ref{resolve-lemma-sections-powers-I-rational} による．
一方，明らかに単射 $\mathfrak m^n \to H^0(X', (\mathcal{I}')^n)$ がある．
$H^0(X', (\mathcal{I}')^n)$ は捩れなしなので，すべての $n$ に対して等号が
成り立つと結論でき，したがって $X = X'$ である．
\end{proof}

\begin{lemma}
\label{resolve-lemma-cohomology-blow-up-rational}
% P07-ERR-0402: second frozen introductory fragment is joined to its declaration.
状況 \ref{resolve-situation-rational} において，$X$ を $\Spec(A)$ の $\mathfrak m$ に沿う
ブローアップとする．$E \subset X$ を例外因子とする．通常どおり
$\mathcal{O}_X(1) = \mathcal{I}$ および
$\mathcal{O}_E(1) = \mathcal{O}_X(1)|_E$ とおくと，次が成り立つ．
\begin{enumerate}
\item $E$ は $\kappa$ 上の固有 Cohen--Macaulay 曲線である．
% P07-ERR-0406: Japanese supplies the list punctuation missing upstream.
\item $\mathcal{O}_E(1)$ は非常に豊富である．
\item $\deg(\mathcal{O}_E(1)) \geq 1$ であり，等号が成り立つのは
$A$ が正則局所環である場合に限る．
\item $H^1(E, \mathcal{O}_E(n)) = 0$ である（$n \geq 0$）．
\item $H^0(E, \mathcal{O}_E(n)) = \mathfrak m^n/\mathfrak m^{n + 1}$
である（$n \geq 0$）．
\end{enumerate}
\end{lemma}

\begin{proof}
$\mathcal{O}_X(1)$ は構成により非常に豊富なので，その特殊ファイバー $E$ への
制限も非常に豊富である．補題 \ref{resolve-lemma-blow-up-normal-rational} により
スキーム $X$ は正規である．すると Divisors, 補題
\ref{divisors-lemma-normal-effective-Cartier-divisor-S1} により，$E$ は
Cohen--Macaulay である．補題 \ref{resolve-lemma-sections-powers-I-rational} を適用でき，
完全列
$$
0 \to \mathcal{I}^{n + 1} \to \mathcal{I}^n \to i_*\mathcal{O}_E(n) \to 0
$$
とコホモロジーの長完全列から (4) と (5) を得る．特に
$$
\deg(\mathcal{O}_E(1)) = \chi(E, \mathcal{O}_E(1)) - \chi(E, \mathcal{O}_E) =
\dim(\mathfrak m/\mathfrak m^2) - 1
$$
である．これは Varieties, 定義
\ref{varieties-definition-degree-invertible-sheaf} による．したがって (3) も従う．
\end{proof}

\begin{lemma}
\label{resolve-lemma-dualizing-rational}
状況 \ref{resolve-situation-rational} において，$A$ は双対化複体
$\omega_A^\bullet$ をもつと仮定する．$\omega_X$ を $X$ の双対化加群とすると，
トレース写像 $H^0(X, \omega_X) \to \omega_A$ は同型であり，したがって標準写像
$f^*\omega_A \to \omega_X$ が存在する．
\end{lemma}

\begin{proof}
Grauert--Riemenschneider（命題
\ref{resolve-proposition-Grauert-Riemenschneider}）により
$Rf_*\omega_X = f_*\omega_X$ である．双対性により短完全列
$$
0 \to f_*\omega_X \to \omega_A \to
\Ext^2_A(R^1f_*\mathcal{O}_X, \omega_A) \to 0
$$
を得る（例えば補題 \ref{resolve-lemma-bound-dualizing-implies-bound} の証明を参照せよ）．
$A$ は有理特異点を定めるので $f_*\omega_X = \omega_A$ である．
\end{proof}

\begin{lemma}
\label{resolve-lemma-dualizing-blow-up-rational}
状況 \ref{resolve-situation-rational} において，$A$ は双対化複体
$\omega_A^\bullet$ をもち，正則でないと仮定する．$X$ を $\Spec(A)$ の
$\mathfrak m$ に沿うブローアップとし，例外因子を $E \subset X$ とする．
$\omega_X$ を $X$ の双対化加群とする．このとき
\begin{enumerate}
\item $\omega_E = \omega_X|_E \otimes \mathcal{O}_E(-1)$ であり，
\item $H^1(X, \omega_X(n)) = 0$ である（$n \geq 0$）．
\item 補題 \ref{resolve-lemma-dualizing-rational} の写像
$f^*\omega_A \to \omega_X$ は全射である．
\end{enumerate}
\end{lemma}

\begin{proof}
以下では補題 \ref{resolve-lemma-cohomology-blow-up-rational} の結果を暗黙に用いる．
$\omega_E = \omega_X|_E \otimes \mathcal{O}_E(-1)$
であることに注意する．これは Duality for Schemes, 補題
\ref{duality-lemma-sheaf-with-exact-support-effective-Cartier} および
\ref{duality-lemma-twisted-inverse-image-closed} による．したがって
$\omega_X|_E = \omega_E(1)$ である．短完全列
$$
0 \to \omega_X(n + 1) \to \omega_X(n) \to i_*\omega_E(n + 1) \to 0
$$
を考える．Algebraic Curves, 補題 \ref{curves-lemma-vanishing-twist} により，
$H^1(E, \omega_E(n + 1)) = 0$ が $n \geq 0$ に対して成り立つ．
したがって写像
$$
\ldots \to H^1(X, \omega_X(2)) \to H^1(X, \omega_X(1)) \to H^1(X, \omega_X)
$$
は全射である．$H^1(X, \omega_X(n))$ は $n \gg 0$ に対して零なので
（Cohomology of Schemes, 補題 \ref{coherent-lemma-kill-by-twisting}），
(2) が成り立つと結論する．

\medskip\noindent
Algebraic Curves, 補題
\ref{curves-lemma-tensor-omega-with-globally-generated-invertible} により，
$\omega_X|_E = \omega_E \otimes \mathcal{O}_E(1)$ は大域的に生成される．
% P07-ERR-0404: Japanese states the intended backward reference; frozen grammar remains recorded.
上で見たように $H^1(X, \omega_X(1)) = 0$ なので，写像
$H^0(X, \omega_X) \to H^0(E, \omega_X|_E)$ は全射である．したがって
$\omega_X$ は大域的に生成される．ゆえに (3) が成り立つ．実際，補題
\ref{resolve-lemma-dualizing-rational} では写像を定義するために
$\Gamma(X, \omega_X) = \omega_A$ を用いている．
\end{proof}

\begin{lemma}
\label{resolve-lemma-rational-to-gorenstein}
$(A, \mathfrak m, \kappa)$ を，有理特異点を定める次元 $2$ の Nagata
正規局所整域とする．$A$ は双対化複体をもつと仮定する．このとき，
特異閉点におけるブローアップの有限列
$$
X = X_n \to X_{n - 1} \to \ldots \to X_1 \to X_0 = \Spec(A)
$$
で，$X_i$ が各 $i$ に対して正規であり，かつ双対化層 $\omega_X$
（$X$ の双対化層）が可逆 $\mathcal{O}_X$-加群となるものが存在する．
\end{lemma}

\begin{proof}
双対化加群 $\omega_A$ は有限 $A$-加群であり，一般点における茎は可逆である．
実際，$\omega_A \otimes_A K$ は分数体 $K$，すなわち $A$ の分数体の双対化加群なので，
階数 $1$ をもつ．したがってブローアップ $b : Y \to \Spec(A)$ で，
$\omega_A$ の $b$ に関する狭義変換が可逆 $\mathcal{O}_Y$-加群となるものが存在する．
Divisors, 補題 \ref{divisors-lemma-blowup-fitting-ideal} を参照せよ．
補題 \ref{resolve-lemma-dominate-by-normalized-blowing-up} により，正規化ブローアップの列
$$
X_n \to X_{n - 1} \to \ldots \to X_1 \to \Spec(A)
$$
を，$X_n$ が $Y$ を支配するように選べる．補題
\ref{resolve-lemma-blow-up-normal-rational} と帰納法により，各
$X_i \to X_{i - 1}$ は単にブローアップである．

\medskip\noindent
$\omega_{X_n}$ は可逆であると主張する．$\omega_{X_n}$ は連接
$\mathcal{O}_{X_n}$-加群なので，その茎が可逆加群であることを見れば十分である．
$x \in X_n$ が正則点ならば，正則スキームが Gorenstein であることから明らかである
（Dualizing Complexes, 補題 \ref{dualizing-lemma-regular-gorenstein}）．
$x$ が $X_n$ の特異点ならば，各像 $x_i \in X_i$（$x$ の像）は
それぞれ特異点である（正則点のブローアップは補題
\ref{resolve-lemma-blowup-regular} により正則だからである）．
補題 \ref{resolve-lemma-dualizing-rational} の標準写像
$f_n^*\omega_A \to \omega_{X_n}$ を考える．各 $i$ に対し，射
$X_{i + 1} \to X_i$ は $x_i$ のブローアップであるか，$x_i$ で同型である．
$x_i$ は常に特異点なので，補題 \ref{resolve-lemma-dualizing-blow-up-rational} と帰納法により，
% P07-ERR-0405: Japanese gives the family of maps a plural-compatible predicate.
写像 $f_i^*\omega_A \to \omega_{X_i}$ は $x_i$ における茎上で常に全射である．したがって
$$
(f_n^*\omega_A)_x \longrightarrow \omega_{X_n, x}
$$
は全射である．一方，$b$ の選び方により，$f_n^*\omega_A$ をその捩れ部分加群で
割ったものは可逆加群 $\mathcal{L}$ である．さらに双対化加群は捩れなしである
（Duality for Schemes, 補題 \ref{duality-lemma-dualizing-module}）．
したがって $\mathcal{L}_x \cong  \omega_{X_n, x}$ であり，証明は完了する．
\end{proof}






\section{形式弧}
\label{resolve-section-arcs}

\noindent
局所 Noether スキーム $X$ をとる．本節では，$X$ における {\it 形式弧}とは，
射 $a : T \to X$ であって，$T$ が完備離散付値環 $R$ のスペクトルであり，
その剰余体 $\kappa$ が，像 $p$（$\Spec(R)$ の閉点の像）の剰余体と
同一視されているものをいう．この場合，形式弧 $a$ は
{\it $p$ を中心とする}という．
形式弧 $T \to X$ が {\it 非特異}であるとは，誘導される写像
$\mathfrak m_p/\mathfrak m_p^2 \to \mathfrak m_R/\mathfrak m_R^2$
が全射であることをいう．

\medskip\noindent
$a : T \to X$，$T = \Spec(R)$ を，閉点 $p$（$X$ の閉点）を中心とする非特異形式弧とする．
$X$ は局所 Noether であると仮定する．
% P07-ERR-0407: the frozen source uses undefined x for the centre p.
$b : X_1 \to X$ を $X$ の $x$ におけるブローアップとする．
$a$ は非特異なので，元 $f \in \mathfrak m_p$ で $R$ の一意化元へ写るものが存在する．
特に，$T$ の一般点は $X$ の $p$ とは異なる点へ写る．言い換えると，$K$ を
$R$ の分数体とすれば，$a$ の制限は射
$\Spec(K) \to X \setminus \{p\}$ を定める．射 $b$ は固有であり，
$X \setminus \{x\}$ 上の同型なので，固有性の付値判定法を適用すると，
次の図式を可換にする一意な射 $a_1$ を得る．
$$
\xymatrix{
T \ar[r]_{a_1} \ar[rd]_a & X_1 \ar[d]^{b} \\
& X
}
$$
$p_1 \in X_1$ を $T$ の閉点の像とする．$p_1$ は
$\kappa = \kappa(p)$-有理点であり，$X_1 \to X$ の $x$ 上のファイバーに
属するので閉点である．
次の分解がある．
$$
\mathcal{O}_{X, x} \to \mathcal{O}_{X_1, p_1} \to R
$$
したがって，$a_1$ も非特異形式弧である．

\medskip\noindent
この操作を反復して，ブローアップの列
$$
\xymatrix{
T \ar[d]_a \ar[rd]_{a_1} \ar[rrd]_{a_2} \ar[rrrd]^{a_3} \\
(X, p) & (X_1, p_1) \ar[l] & (X_2, p_2) \ar[l] &
(X_3, p_3) \ar[l] & \ldots \ar[l]
}
$$
を得る．この種のブローアップ列は次のように特徴づけられる．

\begin{lemma}
\label{resolve-lemma-sequence-blowups}
局所 Noether スキーム $X$ をとる．ブローアップの列
$$
(X, p) = (X_0, p_0) \leftarrow (X_1, p_1) \leftarrow (X_2, p_2) \leftarrow
(X_3, p_3) \leftarrow \ldots
$$
で次の条件を満たすものをとる．
\begin{enumerate}
\item $p_i$ は閉点で，$p_{i - 1}$ へ写り，
$\kappa(p_i) = \kappa(p_{i - 1})$ である．
\item ある $x_1 \in \mathfrak m_p$ が存在し，$\mathfrak m_{p_i}$ における
その像は，$i > 0$ のとき例外因子 $E_i \subset X_i$ を定める．
\end{enumerate}
このとき，この列は上のような非特異形式弧 $a : T \to X$ から得られる．
\end{lemma}

\begin{proof}
$\mathcal{O}_n = \mathcal{O}_{X_n, p_n}$ および
$\mathcal{O} = \mathcal{O}_{X, p}$ と書く．極大イデアルをそれぞれ
$\mathfrak m \subset \mathcal{O}$ および $\mathfrak m_n \subset \mathcal{O}_n$
で表す．

\medskip\noindent
$x_1^t \not \in \mathfrak m_n^{t + 1}$ と主張する．実際，そうでなければ，
局所環 $\mathcal{O}_{n + 1}$ において元 $x_1^t$ は
$(t + 1)E_{n + 1}$ のイデアルに属する．これは $x_1$ が
$E_{n + 1}$ を定めるという仮定に矛盾する．

\medskip\noindent
各 $n$ に対して生成元 $y_{n, 1}, \ldots, y_{n, t_n}$ を
$\mathfrak m_n$ の生成元として選ぶ．仮定 (2) により
$\mathfrak m_n \mathcal{O}_{n + 1} = x_1\mathcal{O}_{n + 1}$
なので，$y_{n, i} = a_{n, i} x_1$ と書ける．ここで
$a_{n, i} \in \mathcal{O}_{n + 1}$ である．写像
$\mathcal{O}_n \to \mathcal{O}_{n + 1}$ は (1) により剰余体上の同型を
定めるので，$c_{n, i} \in \mathcal{O}_n$ を，$a_{n, i}$ と同じ
剰余類をもつように選べる．すると
$$
\mathfrak m_n = (x_1, z_{n, 1}, \ldots, z_{n, t_n}),
\quad z_{n, i} = y_{n, i} - c_{n, i} x_1
$$
となり，元 $z_{n, i}$ は $\mathfrak m_{n + 1}^2$ の元へ写る
（$\mathcal{O}_{n + 1}$ において）．

\medskip\noindent
次を考える．
$$
J_n = \Ker(\mathcal{O} \to \mathcal{O}_n/\mathfrak m_n^{n + 1})
$$
$\mathcal{O}/J_n$ の長さは $n + 1$ であり，さらに
% P07-ERR-0408: the frozen quotient leaves +J_n outside the denominator.
$\mathcal{O}/(x_1) + J_n$ は剰余体に等しいと主張する．$n = 0$ なら明らかである．
$n$ に対して主張が成り立つと仮定する．$f \in J_n$ とする．
このとき $\mathcal{O}_n$ において
$$
f = a x_1^{n + 1} + x_1^n A_1(z_{n, i}) +
x_1^{n - 1} A_2(z_{n, i}) + \ldots + A_{n + 1}(z_{n, i})
$$
と書ける．ここで $a \in \mathcal{O}_n$ であり，$A_i$ は次数 $i$ の斉次式で，
係数は $\mathcal{O}_n$ に属する．$\mathcal{O} \to \mathcal{O}_n$ は剰余体を
同一視するので，$a \in \mathcal{O}$ と選んでよい
（上の $z_{n, i}$ の構成と同様に論じればよい）．
$\mathcal{O}_{n + 1}$ における像をとると，$f$ と $a x_1^{n + 1}$ は
$\mathfrak m_{n + 1}^{n + 2}$ を法として同じ像をもつ．
% P07-ERR-0409: the frozen source uses undefined x_n instead of x_1.
$x_n^{n + 1} \not \in \mathfrak m_{n + 1}^{n + 2}$ なので，
$J_n/J_{n + 1}$ の長さは $1$ であり，主張が従う．

\medskip\noindent
$R = \lim \mathcal{O}/J_n$ を考える．これは $\mathfrak m$-進完備化，すなわち
$\mathcal{O}$ のその進完備化の商なので，完備 Noether 局所環である．一方，その長さは
有限でなく，$x_1$ が極大イデアルを生成する．したがって $R$ は完備離散付値環である．
写像 $\mathcal{O} \to R$ は局所準同型 $\mathcal{O}_n \to R$ へ，
各 $n$ に対して持ち上がる．これを示す方法は二つある．
(1) 各 $n$ に対して同様の手順で $\mathcal{O}_n \to R_n$ を構成し，
$\mathcal{O} \to \mathcal{O}_n \to R_n$ が同型 $R \to R_n$ を経由して
分解することを示す．または (2) Divisors, 補題
\ref{divisors-lemma-characterize-affine-blowup} を用い，$\mathcal{O}_n$ が
反復アフィン・ブローアップ代数の局所化であることから，写像
$\mathcal{O}_n \to R$ を明示的に構成する．以上から，このブローアップ列が
非特異弧 $a : T = \Spec(R) \to X$ に由来することは明らかである．
\end{proof}

\noindent
次の補題は，一種の N\'eron 非特異化補題である．

\begin{lemma}
\label{resolve-lemma-sequence-blowups-along-arc-becomes-nonsingular}
$(A, \mathfrak m, \kappa)$ を次元 $2$ の Noether 局所整域とする．
$A \to R$ を完備離散付値環への全射とする．これは非特異弧
$a : T = \Spec(R) \to \Spec(A)$ を定める．ブローアップの列を
$$
\Spec(A) = X_0 \leftarrow X_1 \leftarrow X_2 \leftarrow X_3 \leftarrow \ldots
$$
とし，これが $a$ から構成されるとする．$A_\mathfrak p$ が正則局所環であると
仮定する．ここで
$\mathfrak p = \Ker(A \to R)$ である．このとき，ある $i$ に対して
% P07-ERR-0410: the frozen statement uses undefined x_i instead of p_i.
スキーム $X_i$ は $x_i$ において正則である．
\end{lemma}

\begin{proof}
$x_1 \in \mathfrak m$ を $R$ の一意化元へ写る元とする．
$\kappa(\mathfrak p) = K$ は $R$ の分数体であることに注意する．
$\mathfrak p = (x_2, \ldots, x_r)$ と，$r$ が最小になるように書く．
$r = 2$ ならば $\mathfrak m = (x_1, x_2)$ であり，$A$ は正則なので補題は成り立つ．
$r > 2$ と仮定する．必要なら番号を付け替えることにより，$x_2$ は
$A_\mathfrak p$ の一意化元へ写ると仮定してよい．
% P07-ERR-0411: the frozen module quotient leaves +(x_2) outside the denominator.
$\mathfrak p/\mathfrak p^2 + (x_2)$ は $x_1$ のある冪で零化される．
$i > 2$ に対し，ある $n_i \geq 0$ と $a_i \in A$ が存在して
$$
x_1^{n_i} x_i - a_i x_2 = \sum\nolimits_{2 \leq j \leq k} a_{jk} x_jx_k
$$
となる．ここで $a_{jk} \in A$ である．
% P07-ERR-0412: every frozen quadratic sum omits the upper bound k <= r.
$n_i = 0$ となる $i$ があれば，生成元のリストから $x_i$ を除ける．
ここでいうリストは $\mathfrak p$ の生成元のリストである．そこで $r$ に関する
帰納法で結論を得る．ある $i$ に対して $a_i$ が単元ならば，生成元のリストから
$x_2$ を除ける．このリストも $\mathfrak p$ の生成元のリストであり，
同じ方法で結論を得る．したがって，次のいずれかが成り立つ．
$a_i \in \mathfrak p$，または
% P07-ERR-0413: the frozen prose uses m_1 although the indexed parameter is m_i.
$a_i = u_i x_1^{m_1} \bmod \mathfrak p$ であり，ここで $m_1 > 0$，
$u_i \in A$ は単元である．ゆえに次のいずれかを得る．
$$
x_1^{n_i} x_i = \sum\nolimits_{2 \leq j \leq k} a_{jk} x_jx_k
\quad\text{or}\quad
x_1^{n_i} x_i - u_i x_1^{m_i} x_2 =
\sum\nolimits_{2 \leq j \leq k} a_{jk} x_jx_k
$$
ブローアップ後に整数 $n_i$，$m_i$ が減少することを示せば，証明は完了する．

\medskip\noindent
これらの方程式がアフィン・ブローアップ代数
$A' = A[\mathfrak m/x_1]$ 上でどうなるかを見る．
$\mathfrak m = (x_1, \ldots, x_r)$ なので，$A'$ は $R$ 上，
% P07-ERR-0414: the frozen source names R although A is the algebra base.
$y_i = x_i/x_1$（$i \geq 2$）で生成される．明らかに $A \to R$ は
写像 $A' \to R$ へ延長され，その核は $(y_2, \ldots, y_r)$ である．
このとき，次のいずれかが成り立つ．
% P07-ERR-0415: the frozen transformed exponent is m_1 - 1 rather than m_i - 1.
$$
x_1^{n_i - 1} y_i = \sum\nolimits_{2 \leq j \leq k} a_{jk} y_jy_k
\quad\text{or}\quad
x_1^{n_i - 1} y_i - u_i x_1^{m_1 - 1} y_2 =
\sum\nolimits_{2 \leq j \leq k} a_{jk} y_jy_k
$$
これで証明は完了する．
\end{proof}




\section{完備化への基底変換}
\label{resolve-section-aux}

\noindent
次の簡単な補題は，以下で有用な道具となる．

\begin{lemma}
\label{resolve-lemma-iso-completions}
$(A, \mathfrak m, \kappa)$ を，有限生成な最大イデアル
$\mathfrak m$ をもつ局所環とする．$X$ を $A$ 上の概形とする．
$Y = X \times_{\Spec(A)} \Spec(A^\wedge)$ とおく．ここで
$A^\wedge$ は $\mathfrak m$-進の意味での $A$ の完備化である．
$q \in Y$ を，その像 $p \in X$ が $\Spec(A)$ の閉点上にある点とすると，
局所環の準同型 $\mathcal{O}_{X, p} \to \mathcal{O}_{Y, q}$ は
完備化の間の同型を誘導する．
\end{lemma}

\begin{proof}
$X$ はアフィンであると仮定してよい．このとき $X = \Spec(B)$ と書ける．
$\mathfrak q \subset B' = B \otimes_A A^\wedge$ を $q$ に対応する素イデアルとし，
$\mathfrak p \subset B$ を $p$ に対応する素イデアルとする．
Algebra, 補題 \ref{algebra-lemma-hathat-finitely-generated} により
$$
B'/(\mathfrak m^\wedge)^n B' =
A^\wedge/(\mathfrak m^\wedge)^n \otimes_A B =
A/\mathfrak m^n \otimes_A B = B/\mathfrak m^n B
$$
がすべての $n$ に対して成り立つ．$\mathfrak m B \subset \mathfrak p$ かつ
$\mathfrak m^\wedge B' \subset \mathfrak q$ であるから，
$B/\mathfrak p^n$ と $B'/\mathfrak q^n$ はともに，上で表示した環を
同じ素イデアルの $n$ 乗で割った商である．これで補題が従う．
\end{proof}

\begin{lemma}
\label{resolve-lemma-port-regularity-to-completion}
$(A, \mathfrak m, \kappa)$ を Noether 局所環とする．
$X \to \Spec(A)$ を局所有限型の射とする．
$Y = X \times_{\Spec(A)} \Spec(A^\wedge)$ とおく．$y \in Y$ の像を
$x \in X$ とする．このとき
\begin{enumerate}
\item $\mathcal{O}_{Y, y}$ が正則ならば，$\mathcal{O}_{X, x}$ は正則である，
\item $y$ が閉ファイバーに属するならば，$\mathcal{O}_{Y, y}$ が正則
$\Leftrightarrow \mathcal{O}_{X, x}$ が正則である，また
\item $X$ が $A$ 上固有ならば，$X$ が正則であることと
$Y$ が正則であることは同値である．
\end{enumerate}
\end{lemma}

\begin{proof}
$A \to A^\wedge$ は忠実平坦であるから
（Algebra, 補題 \ref{algebra-lemma-completion-faithfully-flat}），
$Y \to X$ は平坦である．したがって (1) は
Algebra, 補題 \ref{algebra-lemma-descent-regular} による．
補題 \ref{resolve-lemma-iso-completions} によれば，射 $Y \to X$ は
特殊ファイバーの点における完備局所環の間の同型を誘導する．
ゆえに (2) は More on Algebra, 補題
\ref{more-algebra-lemma-completion-regular} による．
$X$ が $A$ 上固有ならば，$Y$ は $A^\wedge$ 上固有であり
（Morphisms, 補題 \ref{morphisms-lemma-base-change-proper}），
$X$ と $Y$ のすべての閉点は閉ファイバーに属する．
したがって Properties, 補題 \ref{properties-lemma-characterize-regular} により，
$Y$ が正則概形であることと $X$ が正則概形であることは同値である．
\end{proof}

\begin{lemma}
\label{resolve-lemma-descend-admissible-blowup}
$(A, \mathfrak m)$ を，完備化 $A^\wedge$ をもつ Noether 局所環とする．
$U \subset \Spec(A)$ と $U^\wedge \subset \Spec(A^\wedge)$ を
それぞれの穴あきスペクトルとする．$Y \to \Spec(A^\wedge)$ が
$U^\wedge$-許容ブローアップならば，$U$-許容ブローアップ
$X \to \Spec(A)$ で
$Y = X \times_{\Spec(A)} \Spec(A^\wedge)$ を満たすものが存在する．
\end{lemma}

\begin{proof}
定義により，イデアル $J \subset A^\wedge$ で
$V(J) = \{\mathfrak m A^\wedge\}$ を満たし，かつ $Y$ が
$S^\wedge$ の，$J$ により定まる閉部分概形におけるブローアップとなるものが存在する．
Divisors, 定義 \ref{divisors-definition-admissible-blowup} を参照せよ．
$A^\wedge$ は Noether であるから，$\mathfrak m^n A^\wedge \subset J$ となる
$n$ が存在する．
$A^\wedge/\mathfrak m^n A^\wedge = A/\mathfrak m^n$ なので，
$\mathfrak m^n \subset I \subset A$ となるイデアルで，$J = I A^\wedge$ を満たすものが得られる．
% P07-ERR-0416: frozen authority uses the undefined base symbol S here; retained sidecar-only.
$X \to S$ を $I$ におけるブローアップとする．
$A \to A^\wedge$ は平坦であるから，Divisors, 補題
\ref{divisors-lemma-flat-base-change-blowing-up} により，
$X$ の基底変換は $Y$ であると結論できる．
\end{proof}

\begin{lemma}
\label{resolve-lemma-blowup-still-good}
$(A, \mathfrak m, \kappa)$ を次元 $2$ の Nagata 局所正規整域とする．
$A$ は有理特異点を定めると仮定する．また，$A^\wedge$ は $A$ の完備化であり，
正規であると仮定する．
このとき
\begin{enumerate}
\item $A^\wedge$ は有理特異点を定める，また
\item $X \to \Spec(A)$ が $\mathfrak m$ におけるブローアップならば，
閉点 $x \in X$ に対して完備化 $\mathcal{O}_{X, x}$ は正規である．
% P07-ERR-0417: frozen statement omits the completion superscript on O_{X,x}; retained sidecar-only.
\end{enumerate}
\end{lemma}

\begin{proof}
$Y \to \Spec(A^\wedge)$ を，$Y$ が正規である修正とする．
$H^1(Y, \mathcal{O}_Y) = 0$ を示さなければならない．Varieties, 補題
\ref{varieties-lemma-modification-normal-iso-over-codimension-1} により，
$Y \to \Spec(A^\wedge)$ は穴あきスペクトル
$U^\wedge = \Spec(A^\wedge) \setminus \{\mathfrak m^\wedge\}$ 上の同型である．
補題 \ref{resolve-lemma-dominate-by-scheme-modification} により，
$U^\wedge$-許容ブローアップ $Y' \to \Spec(A^\wedge)$ で $Y$ を支配するものが存在する．
補題 \ref{resolve-lemma-descend-admissible-blowup} により，
$U$-許容ブローアップ $X \to \Spec(A)$ で，その $A^\wedge$ への
基底変換が $Y$ を支配するものが存在する．
$A$ は Nagata なので，$X$ をその正規化で置き換えることができる．
この置換の後，$X \to \Spec(A)$ は正規修正である（ただし，もはや
$U$-許容ブローアップではない可能性がある）．
$H^1(X, \mathcal{O}_X) = 0$ である．これは $A$ が有理特異点を定めるからである．
したがって
$H^1(X \times_{\Spec(A)} \Spec(A^\wedge),
\mathcal{O}_{X \times_{\Spec(A)} \Spec(A^\wedge)}) = 0$
が平坦基底変換により従う（Cohomology of Schemes, 補題
\ref{coherent-lemma-flat-base-change-cohomology}，および Algebra, 補題
\ref{algebra-lemma-completion-flat} による $A \to A^\wedge$ の平坦性）．
補題 \ref{resolve-lemma-exact-sequence} により
$H^1(Y, \mathcal{O}_Y) = 0$ が得られる．

\medskip\noindent
最後に，$X \to \Spec(A)$ を $\Spec(A)$ の $\mathfrak m$ における
ブローアップとする．このとき $Y = X \times_{\Spec(A)} \Spec(A^\wedge)$ は
$\Spec(A^\wedge)$ の $\mathfrak m^\wedge$ におけるブローアップである．
補題 \ref{resolve-lemma-blow-up-normal-rational} により，$Y$ と $X$ はともに正規である．
一方，$A^\wedge$ は優秀である
（More on Algebra, 命題 \ref{more-algebra-proposition-ubiquity-excellent}）．
したがって，$Y$ の任意のアフィン開部分は優秀な正規整域のスペクトルである
（More on Algebra, 補題 \ref{more-algebra-lemma-finite-type-over-excellent}）．
ゆえに $y \in Y$ に対する環準同型
$\mathcal{O}_{Y, y} \to \mathcal{O}_{Y, y}^\wedge$ は正則であり，
More on Algebra, 補題 \ref{more-algebra-lemma-normal-goes-up} により
$\mathcal{O}_{Y, y}^\wedge$ は正規である．
$x \in X$ が特殊ファイバーの閉点ならば，一意な閉点
$y \in Y$ で $x$ の上にあるものが存在する．$\mathcal{O}_{X, x} \to \mathcal{O}_{Y, y}$ は
完備化の間の同型を誘導するので（補題 \ref{resolve-lemma-iso-completions}），
結論が従う．
\end{proof}

\begin{lemma}
\label{resolve-lemma-formally-unramified}
$(A, \mathfrak m)$ を局所 Noether 環とする．
$X$ を $A$ 上の概形とする．次を仮定する．
\begin{enumerate}
\item $A$ は解析的に非分岐である
（Algebra, 定義 \ref{algebra-definition-analytically-unramified}），
\item $X$ は $A$ 上局所有限型である，また
\item $X \to \Spec(A)$ は $X$ の既約成分の一般点において \'etale である．
\end{enumerate}
このとき $X$ の正規化は $X$ 上有限である．
\end{lemma}

\begin{proof}
$A$ は解析的に非分岐なので，Algebra, 補題
\ref{algebra-lemma-analytically-unramified-easy} により被約である．
$X$ の正規化は $X$ の被約化のみに依存するから，$X$ をその被約化
$X_{red}$ で置き換えてよい．$X_{red} \to X$ は開部分 $U$ 上の同型であることに注意する．
この開部分では $X \to \Spec(A)$ が \'etale である．実際，$U$ は被約である（Descent, 補題
\ref{descent-lemma-reduced-local-smooth}）．したがって，この置換の後にも
条件 (3) は成り立つ．さらに，$X = \Spec(B)$ はアフィンであると仮定してよいし，
そう仮定する．

\medskip\noindent
準同型
$$
K = \prod\nolimits_{\mathfrak p \subset A\text{ 極小}} \kappa(\mathfrak p)
\longrightarrow
K^\wedge = \prod\nolimits_{\mathfrak p^\wedge \subset A^\wedge\text{ 極小}}
\kappa(\mathfrak p^\wedge)
$$
は単射である．実際，$A \to A^\wedge$ は忠実平坦であり
（Algebra, 補題 \ref{algebra-lemma-completion-faithfully-flat}），
したがって極小素イデアルの集合の間に全射を誘導する
（平坦環準同型に対する下降定理による．Algebra, 節
\ref{algebra-section-going-up} を参照せよ）．環は Noether なので，
両辺は体の有限積である．
$L = \prod_{\mathfrak q \subset B\text{ 極小}} \kappa(\mathfrak q)$ とおく．
仮定 (3) により $L = B \otimes_A K$ であり，
$K \to L$ は有限 \'etale 環準同型である（これは
$A \to B$ が一般点上有限だからである．たとえば Algebra, 補題
\ref{algebra-lemma-generically-finite}，または Morphisms, 節
\ref{morphisms-section-generically-finite} のより詳しい結果を用いればよい）．
$B$ は被約なので $B \subset L$ である．したがって
$$
C = B \otimes_A A^\wedge \subset
L \otimes_A A^\wedge = L \otimes_K K^\wedge = M
$$
となる．このとき $M$ は $C$ の全商環であり，$K^\wedge$ 上の
有限分離代数として体の有限積である．したがって $C$ は被約であり，
その正規化 $C'$ は $C$ の $M$ における整閉包である．
$B'$ は $B$ の正規化であり，$B$ の $L$ における整閉包である．
$A \to A^\wedge$ の平坦性により，単射
$B' \otimes_A A^\wedge \to M$ が得られ，その像は $C'$ に含まれる．図式
$$
B' \otimes_A A^\wedge \longrightarrow C'
$$
を考える．$A^\wedge$ は Nagata なので（Algebra, 補題
\ref{algebra-lemma-Noetherian-complete-local-Nagata}），
$C'$ は $C = B \otimes_A A^\wedge$ 上有限である（Algebra, 補題
\ref{algebra-lemma-Noetherian-complete-local-Nagata} および
\ref{algebra-lemma-nagata-in-reduced-finite-type-finite-integral-closure} を参照せよ）．
$C$ は Noether だから，$B' \otimes_A A^\wedge$ は
$C = B \otimes_A A^\wedge$ 上有限であると結論できる．
したがって忠実平坦降下（Algebra, 補題
\ref{algebra-lemma-descend-properties-modules}）により，
$B'$ は $B$ 上有限である．これが示すべきことであった．
\end{proof}

\begin{lemma}
\label{resolve-lemma-normalization-completion}
$(A, \mathfrak m, \kappa)$ を Noether 局所環とする．
$X \to \Spec(A)$ を局所有限型の射とする．
$Y = X \times_{\Spec(A)} \Spec(A^\wedge)$ とおく．
$Y$ における特殊ファイバーの補集合が正規ならば，正規化
$X^\nu \to X$ は有限であり，$X^\nu$ の $\Spec(A^\wedge)$ への
基底変換は $Y$ の正規化を復元する．
\end{lemma}

\begin{proof}
$X = \Spec(B)$ がアフィンで，$B$ が有限型 $A$-代数である場合へ
直ちに帰着できる．$C = B \otimes_A A^\wedge$ とおけば
$Y = \Spec(C)$ となる．$A \to A^\wedge$ は忠実平坦なので，
任意の素イデアル $\mathfrak q \subset B$ に対して，素イデアル
$\mathfrak r \subset C$ で $\mathfrak q$ の上にあるものが存在する．このとき
$B_\mathfrak q \to C_\mathfrak r$ は忠実平坦である．したがって
$\mathfrak q$ が $\mathfrak m$ の上にないならば，$C_\mathfrak r$ は
$Y$ に関する仮定により正規であり，Algebra, 補題
\ref{algebra-lemma-descent-normal} により $B_\mathfrak q$ も正規である．
このようにして，$X$ は特殊ファイバーの外で正規であることがわかる．

\medskip\noindent
完備 Noether 局所環 $A^\wedge$ は Nagata であることを思い出そう
（Algebra, 補題 \ref{algebra-lemma-Noetherian-complete-local-Nagata}）．
したがって正規化 $Y^\nu \to Y$ は有限であり
（Morphisms, 補題 \ref{morphisms-lemma-nagata-normalization}），
特殊ファイバーの外では同型である．$Y^\nu = \Spec(C')$ と書こう．
このとき $C \to C'$ は有限であり，$V(\mathfrak m C)$ の外では同型である．
$B \to C$ は平坦で，同型 $B/\mathfrak m B \to C/\mathfrak m C$ を誘導するから，
有限環準同型 $B \to B'$ で，その $C$ への基底変換が
$C \to C'$ を復元するものが存在する．More on Algebra, 補題
\ref{more-algebra-lemma-application-formal-glueing} および注意
\ref{more-algebra-remark-formal-glueing-algebras} を参照せよ．
ゆえに有限射 $X' \to X$ で，特殊ファイバーの外では同型であり，
その基底変換が $Y^\nu \to Y$ を復元するものが得られる．
第1段落の議論により，$X'$ は特殊ファイバー上にない点で正規である．
特殊ファイバー上の点 $x \in X'$ に対して，対応する点
$y \in Y^\nu$ と平坦準同型
$\mathcal{O}_{X', x} \to \mathcal{O}_{Y^\nu, y}$ がある．
$\mathcal{O}_{Y^\nu, y}$ は正規なので，$\mathcal{O}_{X', x}$ も正規である．
Algebra, 補題 \ref{algebra-lemma-descent-normal} を参照せよ．
したがって $X'$ は正規であり，$X$ の正規化であることが従う．
\end{proof}

\begin{lemma}
\label{resolve-lemma-normalized-blowup-completion}
$(A, \mathfrak m, \kappa)$ を，その完備化 $A^\wedge$ が正規である
Noether 局所整域とする．このとき，正規化ブローアップの任意の列
$$
Y_n \to Y_{n - 1} \to \ldots \to Y_1 \to \Spec(A^\wedge)
$$
が与えられれば，（固有な）正規化ブローアップの列
$$
X_n \to X_{n - 1} \to \ldots \to X_1 \to \Spec(A)
$$
で，その $A^\wedge$ への基底変換が与えられた列を復元するものが存在する．
\end{lemma}

\begin{proof}
列 $Y_n \to \ldots \to Y_1 \to Y_0 = \Spec(A^\wedge)$ が与えられたとする．
$X_n \to \ldots \to X_1 \to X_0 = \Spec(A)$ を帰納的に構成する．
基底の場合は $i = 0$ である．$X_i$ で，その基底変換が $Y_i$ であるものが
与えられたとする．$Y'_i \to Y_i$ を閉点 $y_i \in Y_i$ における
ブローアップで，$Y_{i + 1}$ が $Y_i$ の正規化となるものとする．
% P07-ERR-0418: frozen authority names Y_i rather than the intermediate blowup Y'_i; retained sidecar-only.
$Y_i$ と $X_i$ の閉ファイバーは同型なので，点 $y_i$ は閉点
$x_i$ に対応し，後者は $X_i$ の特殊ファイバー上にある．
$X'_i \to X_i$ を $X_i$ の $x_i$ におけるブローアップとする．
このとき $X'_i$ の $\Spec(A^\wedge)$ への基底変換は $Y'_i$ と同型である．
補題 \ref{resolve-lemma-normalization-completion} により，正規化
$X_{i + 1} \to X'_i$ は有限であり，その $\Spec(A^\wedge)$ への
基底変換は $Y_{i + 1}$ と同型である．
\end{proof}














\section{有理二重点}
\label{resolve-section-rational-double-points}

\noindent
節 \ref{resolve-section-rational-singularities} では，$2$ 次元有理特異点の
特異点解消が Gorenstein の場合に帰着されることを論じた．
Gorenstein 有理曲面特異点は有理二重点である．
以下では，明示的な計算によってこれらを解消する．

\medskip\noindent
Examples, 節 \ref{examples-section-bad} の議論によれば，正規 Noether 局所整域
$A$ で，その完備化が $\mathbf{C}[[x, y, z]]/(z^2)$ と同型であるものが
存在する．この場合，$A$ は有理二重点特異点をもつと言えるが，
他方で $\Spec(A)$ は特異点解消をもたない．$A$ が Nagata 環ならば，
このような挙動は生じ得ない．Algebra, 補題
\ref{algebra-lemma-local-nagata-domain-analytically-unramified} を参照せよ．

\medskip\noindent
しかし，事態はさらに悪い．局所正規 Nagata 整域 $A$ で，その完備化が
$\mathbf{C}[[x, y, z]]/(yz)$ であるものと，その完備化が
$\mathbf{C}[[x, y, z]]/(y^2 - z^3)$ である別のものが存在する．
これは \cite{Nishimura-few} の例 2.5 である．このため，本節では
考える環の完備化が正規であると仮定する必要がある．

\begin{situation}
\label{resolve-situation-rational-double-point}
% P07-ERR-0419: the frozen English situation has a missing finite verb; Japanese states the hypotheses grammatically.
ここで $(A, \mathfrak m, \kappa)$ は，有理特異点を定め，完備化が正規であり，
Gorenstein である次元 $2$ の Nagata 局所正規整域とする．
$A$ は正則でないと仮定する．
\end{situation}

\noindent
本節の議論により，この状況では特異点を繰り返しブローアップすることで
$\Spec(A)$ が解消されることを示す．証明の途中で次の補題を用いる．

\begin{lemma}
\label{resolve-lemma-issquare}
$\kappa$ を体とし，$I \subset \kappa[x, y]$ をイデアルとする．
次が成り立つとする．
$$
a + b x + c y + d x^2 + exy + f y^2 \in I^2
$$
% P07-ERR-0420: frozen authority uses the undefined field k rather than kappa; retained sidecar-only.
ここで $a, b, c, d, e, f \in k$ はすべてが零ではないとする．
$I$ の $\kappa[x, y]$ における余長が $> 1$ ならば，
$a + b x + c y + d x^2 + exy + f y^2 = j(g + hx + iy)^2$
となる $j, g, h, i \in \kappa$ が存在する．
\end{lemma}

\begin{proof}
偏導関数 $b + 2dx + ey$ と $c + ex + 2fy$ を考える．
Leibniz 則により，これらは $I$ に含まれる．いずれかが零でなければ，
線形座標変換，すなわち $x \mapsto \alpha + \beta x + \gamma y$ および
$y \mapsto \delta + \epsilon x + \zeta y$ という形の変換を施すことで，
$x \in I$ と仮定してよい．このとき $I = (x)$ または $I = (x, F)$ であり，
$F$ は次数 $\geq 2$ の $y$ に関するモニック多項式である．
第1の場合，主張は明らかである．第2の場合，$I^2$ の任意の元は
次の形に書けることに注意する．
$$
A(x, y) x^2 + B(y) x F + C(y) F^2
$$
ここで $A(x, y) \in \kappa[x, y]$ かつ $B, C \in \kappa[y]$ である．
したがって
$$
a + b x + c y + d x^2 + exy + f y^2 = A(x, y) x^2 + B(y) x F + C(y) F^2
$$
であり，次数を考えれば $B = C = 0$ かつ $A$ は定数であることがわかる．

\medskip\noindent
証明を終えるには，両方の偏導関数が零である場合を扱う必要がある．
これは標数 $2$ の場合にしか起こらず，このとき
$$
a + d x^2 + f y^2 \in I^2
$$
$f$ は零でないと仮定してよい（そうでなければ $x$ と $y$ の役割を入れ替える）．
$f$ で割ることにより，$\kappa$ の標数が $2$ であって
$$
a + d x^2 + y^2 \in I^2
$$
となる場合を得る．$a$ と $d$ が $\kappa$ における平方ならば，証明は終わる．
そうでなければ，導分 $\theta : \kappa \to \kappa$ で，
$\theta(a) \not = 0$ または $\theta(d) \not = 0$ を満たすものが存在する．Algebra, 補題
\ref{algebra-lemma-derivative-zero-pth-power} を参照せよ．
これを $\kappa[x, y]$ の導分へ，$\theta(x) = \theta(y) = 0$ とおくことで延長できる．
このとき
$$
\theta(a) + \theta(d) x^2 \in I
$$
$\theta(d) = 0$ の場合は矛盾である．したがって，
$\alpha + x^2 \in I$ を満たすある $\alpha \in \kappa$ があると仮定してよい．
上の結果と合わせると $a + \alpha d + y^2 \in I$ が得られる．ゆえに
$$
J = (\alpha + x^2, a + \alpha d + y^2) \subset I
$$
であり，その余次元は高々 $2$ である．$J/J^2$ は
$\kappa[x, y]/J$ 上自由で，基底は $\alpha + x^2$ と
$a + \alpha d + y^2$ であることに注意する．したがって
$a + d x^2 + y^2 =
1 \cdot (a + \alpha d + y^2) + d \cdot (\alpha + x^2) \in I^2$ から，
包含 $J \subset I$ は真であることが従う．ゆえに，
$g + hx + iy + jxy$ という形の零でない元が $I$ の中に見つかる．
$j = 0$ ならば $I$ は一次形式を含むので，第1段落と同様に結論できる．
したがって $j \not = 0$ であり，$\dim_\kappa(I/J) = 1$ である
（そうでなければ，$I$ の中に上の形の元で $j = 0$ となるものを見つけられる）．
したがって $I$ は
$(\alpha + x^2, \beta + y^2, g + hx + iy + jxy)$
という形で，$j \not = 0$ かつ余長が $3$ である．
この場合，$a + dx^2 + y^2 \in I^2$ は不可能である．
これは直接計算で示せるが，次のように論じる方を選ぶ．すなわち，
この主張を証明するには $\kappa$ が代数閉であると仮定してよい．
そこで座標変換 $x \mapsto \sqrt{\alpha} + x$ および
$y \mapsto \sqrt{\beta} + y$ を行い，
$I = (x^2, y^2, g' + h'x + i'y + jxy)$ と仮定できる．ここで $j$ は同じものである．
このとき $g' = h' = i' = 0$ でなければ，$I$ の余長は $3$ ではない．
したがって $I = (x^2, y^2, xy)$ が得られ，結果は明らかである．
\end{proof}

\noindent
$(A, \mathfrak m, \kappa)$ を状況
\ref{resolve-situation-rational-double-point} のものとする．
$X \to \Spec(A)$ を $\mathfrak m$ における $\Spec(A)$ のブローアップとする．
補題 \ref{resolve-lemma-blow-up-normal-rational} により $X$ は正規である．
$X$ のすべての特異点は，補題 \ref{resolve-lemma-rational-propagates} により
有理特異点である．$\omega_A = A$ なので，補題
\ref{resolve-lemma-dualizing-blow-up-rational} から
$\omega_X \cong \mathcal{O}_X$ がわかる（規約については注意
\ref{resolve-remark-dualizing-setup} の議論を参照せよ）．
したがって $X$ のすべての特異点は Gorenstein である．
さらに，$X$ の閉点における局所環の完備化は，補題
\ref{resolve-lemma-blowup-still-good} により正規である．
言い換えれば，$\Spec(A)$ をブローアップすると正規曲面 $X$ が得られ，
その特異点は状況 \ref{resolve-situation-rational-double-point} のものとなる．
以下では，これを断りなく用いる．
（なお，以下の議論の途中で，これらの特異点が有限個であることもわかる．）

\medskip\noindent
$E \subset X$ を例外因子とする．補題 \ref{resolve-lemma-dualizing-blow-up-rational} により
$\omega_E = \mathcal{O}_E(-1)$ である．補題
\ref{resolve-lemma-cohomology-blow-up-rational} により
$\kappa = H^0(E, \mathcal{O}_E)$ である．したがって $E$ は Gorenstein 曲線であり，
Algebraic Curves, 節 \ref{curves-section-Riemann-Roch} で論じた Riemann--Roch により
$$
\chi(E, \mathcal{O}_E) = 1 - g = -(1/2) \deg(\omega_E) =
(1/2)\deg(\mathcal{O}_E(1))
$$
となる．ここで $g = \dim_\kappa H^1(E, \mathcal{O}_E) \geq 0$ である．
$\deg(\mathcal{O}_E(1))$ は Varieties, 補題
\ref{varieties-lemma-ampleness-in-terms-of-degrees-components} により正なので，
$g = 0$ かつ $\deg(\mathcal{O}_E(1)) = 2$ が得られる．したがって
$$
\dim_\kappa (\mathfrak m^n/\mathfrak m^{n + 1}) = 2n + 1
$$
である．これは補題 \ref{resolve-lemma-cohomology-blow-up-rational} と
$E$ 上の Riemann--Roch による．

\medskip\noindent
$x_1, x_2, x_3 \in \mathfrak m$ を，その像が
$\mathfrak m/\mathfrak m^2$ の基底となるように選ぶ．
$\dim_\kappa(\mathfrak m^2/\mathfrak m^3) = 5$ なので，
$x_i x_j$（$i \geq j$）の像は，この $\kappa$-ベクトル空間において
ある関係式を満たす．言い換えれば，$a_{ij} \in A$（$i \geq j$）で，
すべてが $\mathfrak m$ に含まれるわけではなく，次を満たすものを選べる．
$$
a_{11} x_1^2 + a_{12} x_1x_2 + a_{13}x_1x_3 + a_{22} x_2^2 +
a_{23} x_2x_3 + a_{33} x_3^2 =
\sum a_{ijk} x_ix_jx_k
$$
ここで，ある $a_{ijk} \in A$ に対して $i \leq j \leq k$ とする．
$a \mapsto \overline{a}$ により写像 $A \to \kappa$ を表す．二次形式
$q = \sum \overline{a}_{ij} t_i t_j \in \kappa[t_1, t_2, t_3]$
は，以上の選択のもとで $\kappa^*$ の元による乗法を除いてよく定まる．
議論の途中で $\overline{a}_{ij} = 0$ が $\kappa$ において成り立つことが
わかったなら，項 $a_{ij} x_i x_j$ を右辺に吸収して
$a_{ij} = 0$ と仮定できる．この操作は $a_{ijk}$ を変えるが，
他の $a_{i'j'}$ は変えない．

\medskip\noindent
このブローアップは $3$ 個の，「変数」$x_1, x_2, x_3$ に対応する
アフィンチャートで被覆される．対称性により，そのうち一つを調べれば十分である．
そのため
$$
A' = A[\mathfrak m/x_1]
$$
をアフィン・ブローアップ代数とする（Algebra, 節
\ref{algebra-section-blow-up}）．$x_1, x_2, x_3$ は $\mathfrak m$ を生成するので，
$A'$ は $y_2 = x_2/x_1$ と $y_3 = x_3/x_1$ により $A$ 上生成される．
式を簡単にするため，ときどき $y_1 = 1$ とおく．さらに，上の関係式を見ると
$$
a_{11} + a_{12} y_2 + a_{13} y_3 + a_{22} y_2^2 +
a_{23} y_2y_3 + a_{33} y_3^2 =
x_1 (\sum a_{ijk} y_iy_jy_k)
$$
が $A'$ において成り立つ．$x_1 \in A'$ は例外因子 $E$ を
$X$ のこのアフィン開部分上で定めることを思い出そう．したがって，
これは概形論的には
$$
\kappa[y_2, y_3]/
(\overline{a}_{11} + \overline{a}_{12} y_2 + \overline{a}_{13} y_3 +
\overline{a}_{22} y_2^2 + \overline{a}_{23} y_2y_3 + \overline{a}_{33} y_3^2)
$$
で与えられる．言い換えれば，
$E \subset \mathbf{P}^2_\kappa = \text{Proj}(\kappa[t_1, t_2, t_3])$ は
上で導入した二次形式 $q$ の零点概形である．

\medskip\noindent
二次形式 $q$ は，$A$ により定まる特異点の重要な不変量である．
$q$ が一次形式の平方に $\kappa^*$ の元を掛けたものならば
{\bf 場合 II}，そうでなければ {\bf 場合 I} と呼ぶ．
$x_1, x_2, x_3$ の選択を変えた後で
$$
x_3^2 = \sum a_{ijk}x_ix_jx_k
$$
が局所環 $A$ において成り立つことと，場合 II にあることは同値である．

\medskip\noindent
$\mathfrak m' \subset A'$ を $\mathfrak m$ の上にある最大イデアルとし，
その剰余体を $\kappa'$ とする．言い換えれば，$\mathfrak m'$ は
例外因子の閉点 $p \in E$ に対応する．全射
$$
\kappa[y_2, y_3] \to \kappa'
$$
の核は二つの元 $f_2, f_3 \in \kappa[y_2, y_3]$ で生成されることを思い出そう
（たとえば Algebra, 例 \ref{algebra-example-spec-kxy}，または
Algebra, 補題 \ref{algebra-lemma-dim-affine-space} の証明を参照せよ）．
$z_2, z_3 \in A'$ を，$f_2, f_3$ に $\kappa[y_2, y_3]$ において写るものとする．
$\mathfrak m' = (x_1, z_2, z_3)$ である．実際，$x_2$ と $x_3$ は
$x_1$ で $A'$ において割り切れるようになる．

\medskip\noindent
{\bf 主張．}$X$ が $p$ で特異ならば，$\kappa' = \kappa$ であるか，
場合 II にある．実際，$A'_{\mathfrak m'}$ が特異ならば
$\dim_{\kappa'} \mathfrak m'/(\mathfrak m')^2 = 3$ であり，したがって
$\dim_{\kappa'} \overline{\mathfrak m}'/(\overline{\mathfrak m}')^2 = 2$
となる．% P07-ERR-0421: frozen authority drops mathfrak on the following maximal-ideal symbol; retained sidecar-only.
ここで $\overline{m}'$ は
$\mathcal{O}_{E, p} = \mathcal{O}_{X, p}/x_1\mathcal{O}_{X, p}$ の最大イデアルである．
このことから
$$
q(1, y_2, y_3) =
\overline{a}_{11} + \overline{a}_{12} y_2 + \overline{a}_{13} y_3 +
\overline{a}_{22} y_2^2 + \overline{a}_{23} y_2y_3 + \overline{a}_{33} y_3^2
\in (f_2, f_3)^2
$$
が従う．そうでなければ，$z_2$ と $z_3$ の類の間に
$\overline{\mathfrak m}'/(\overline{\mathfrak m}')^2$ における関係式が存在することになる．
これで主張は補題 \ref{resolve-lemma-issquare} から従う．

\medskip\noindent
場合 I における解消．上の主張により，$X$ の任意の特異点は
$\kappa$-有理的である．そのような特異点 $p$ を選ぶ．
$x_1, x_2, x_3 \in \mathfrak m$ を，$p$ が上で述べたチャートに属し，
座標が $y_2 = y_3 = 0$ となるように選べる．特異点であるから，
主張の証明と同様に論じて
$q(1, y_2, y_3) \in (y_2, y_3)^2$ が得られる．
したがって $a_{11} = a_{12} = a_{13} = 0$ かつ
$q(t_1, t_2, t_3) = q(t_2, t_3)$ となるように選べる．よって
$$
E = V(q) \subset \mathbf{P}^1_\kappa
$$
% P07-ERR-0422: frozen authority prints P^1 although the described plane quadratic lies in P^2; retained sidecar-only.
は，$p$ で交わる二つの相異なる直線の和であるか，または次数 $2$ で一意な
$\kappa$-有理点をもつ曲線である
（細部は省略する．$q$ がスカラー倍を除いて一次形式の平方でないことを用いよ）．
いずれの場合にも，$X$ は一意な特異点 $p$ をもち，それは
$\kappa$-有理的であると結論できる．この場合には，もう少し情報が必要である．
まず，上の表示の高次項を見ると，$\overline{a}_{111} = 0$ であることがわかる．
これは $p$ が特異だからである．したがって
$a_{111} = b_{111} x_1 \bmod (x_2, x_3)$ と書ける．ここで
$b_{111} \in A$ である．このとき，$p$ における
$x_1, y_2, y_3$（$\mathfrak m'$ の生成元）に関する二次形式は
$$
q' =
\overline{b}_{111} t_1^2 +
\overline{a}_{112} t_1 t_2 + \overline{a}_{113} t_1 t_3 +
\overline{a}_{22} t_2^2 +
\overline{a}_{23} t_2 t_3 +
\overline{a}_{33} t_3^2
$$
$E' = V(q')$ は直線 $t_1 = 0$ と二点，または次数 $2$ の一点で交わる．
したがって $p$ は場合 I に属する．

\medskip\noindent
ブローアップ $X' \to X$ を $X$ の $p$ において行うと，再び特異点 $p'$ をもつとする．
このとき $p'$ は $\kappa$-有理点であり，さらにブローアップして
$X'' \to X'$ を得られる．この過程が停止しなければ，ブローアップの列
$$
\Spec(A) \leftarrow X \leftarrow X' \leftarrow X'' \leftarrow \ldots
$$
を得る．この状況に補題 \ref{resolve-lemma-sequence-blowups} が適用できることを示したい．
そのためには，元 $x_1$ の $\mathfrak m$ における選択について述べる必要がある．
$A$ が場合 I にあり，$X$ が特異点をもつと仮定する．このとき，
$x_1 \in \mathfrak m$ が {\it 良い座標} であるとは，$x_2, x_3$ の任意の
（同値なことに，ある）選択に対して，二次形式 $q(t_1, t_2, t_3)$ が，
$q(0, t_2, t_3)$ はスカラーと平方との積ではないという性質をもつこととする．
良い座標が存在することは上で見た．$x_1$ が良い座標ならば，特異点
$p \in E$ は $X$ の超曲面 $t_1 = 0$ 上にない．実際，この超曲面は有理点をもたないか，
もったとしてもその点で $X$ は特異でない．これは，$x_1$ の
$\mathcal{O}_{X, p}$ における像が例外因子 $E$ を切り出すという主張と同値であることに注意する．
上の計算から，$x_1$ が $A$ の良い座標ならば，
$x_1 \in \mathfrak m'\mathcal{O}_{X, p}$ は $p$ の良い座標であることがわかる．
もちろん，ここでは良い座標の概念が，計算に用いた $x_2$, $x_3$ の選択に
依存しないことを用いている．したがって $x_1$ は $p'$, $p''$ などにおいて
良い座標に写る．ゆえに補題 \ref{resolve-lemma-sequence-blowups} が適用でき，
このブローアップ列は非特異弧 $A \to R$ から生じる．
このとき写像 $A^\wedge \to R$ は全射である．$A$ の完備化は正規なので，
補題 \ref{resolve-lemma-sequence-blowups-along-arc-becomes-nonsingular} により，
有限回のブローアップの後，
$$
\Spec(A^\wedge) \leftarrow X^\wedge \leftarrow (X')^\wedge \leftarrow \ldots
$$
得られる概形 $(X^{(n)})^\wedge$ は正則である．
$(X^{(n)})^\wedge \to X^{(n)}$ は完備局所環の間の同型を誘導するので
（補題 \ref{resolve-lemma-iso-completions}），$X^{(n)}$ についても同じことが成り立つ．

\medskip\noindent
場合 II における解消．この場合
$$
x_3^2 = \sum a_{ijk}x_ix_jx_k
$$
が $A$ において成り立つような生成元 $x_1, x_2, x_3$ を
$\mathfrak m$ に対して選べる．このとき $q = t_3^2$ かつ
$E = 2C$ であり，$C$ は直線である．$A'$ において
$$
y_3^2 = x_1(\sum a_{ijk} y_iy_jy_k)
$$
を得ることを思い出そう．$X$ は正規であることがわかっているので，
離散付値環 $\mathcal{O}_{X, \xi}$ が得られる．ここで $\xi$ は $C$ の一般点である．
元 $y_3 \in A'$ は $\mathcal{O}_{X, \xi}$ の一意化元に写る．
$x_1$ は概形論的に $E$ を切り出し，これは $C$ に重複度 $2$ を付けたものなので，
$x_1$ は $y_3^2$ と単元との積として $\mathcal{O}_{X, \xi}$ に属する．
上の等式を見ると
$$
h(y_2) = \overline{a}_{111} + \overline{a}_{112} y_2 +
\overline{a}_{122} y_2^2 +
\overline{a}_{222} y_2^3
$$
は $\xi$ の剰余体で零でないと結論できる．
ここで $p \in C$ が特異点を定めると仮定する．このとき $y_3$ は
$p$ で零であり，上の主張の証明で用いた議論により，$p$ は $h$ の零点に対応する．
$h$ が $p$ に二重零点をもたないならば，二次形式 $q'$ は $p$ において
平方ではなく，$p$ は上で扱った場合 I に属すると結論できる\footnote{
$p$ における $A'$ の最大イデアルは $y_3, x_1$ と第三の元 $g$ により生成され，
その $\kappa[y_2]$ における像は，$h$ の $p$ に対応する素因子である．
この素因子が $h$ を二回割らないならば，$p$ における二次形式は
$$
y_3^2 - x_1((something)x_1 + (something)y_3 + (unit)g)
$$
という形となり，$\kappa[y_3, x_1, g]$ における平方には決してならない．}．
$h$ の次数は $3$ なので，場合 II に属する特異点 $p \in C$ は高々一つであり，
しかも $\kappa$-有理的である．$x_1, x_2, x_3$ の選択を変えれば，
これが点 $y_2 = y_3 = 0$ であると仮定してよい．このとき
$h = \overline{a}_{122} y_2^2 + \overline{a}_{222} y_2^3$ である．
さらに，$\overline{a}_{113} = 0$ でなければ，二次形式 $q'$ は
正しい形をもたない．
したがって局所環 $\mathcal{O}_{X, p}$ は次の段落で述べる特異点を定める．

\medskip\noindent
最後に扱うのは，生成元 $x_1, x_2, x_3$ を $\mathfrak m$ に対して
次が成り立つように選べる場合である．
$$
x_3^2 + x_1(a x_2^2 + b x_2x_3 + c x_3^2) \in \mathfrak m^4
$$
ここで $a, b, c \in A$ である．これは場合 II の部分クラスである．
$\overline{a} = 0$ ならば
$a = a_1 x_1 + a_2 x_2 + a_3 x_3$ と書け，ブローアップ後に
$$
y_3^2 + x_1(a_1 x_1 y_2^2 + a_2 x_1 y_2^3 + a_3 x_1 y_2^2 y_3 +
b y_2 y_3 + c y_3^2) = x_1^2 (\sum a_{ijkl}y_iy_jy_ky_l)
$$
を得る．これは $X$ が正規でないことを意味する\footnote{実際，この等式から
$1$ 次元の軌跡 $x_1 = y_3 = 0$ に沿って特異なものが生じるが，
正規曲面ではこれは起こり得ない．}．これは矛盾である．前段落の結果により，
ブローアップ $X$ が場合 II に属する特異点 $p$ をもつならば，
そのような点は一つだけで，$\kappa$-有理的である．
アフィン・ブローアップ代数 $A[\frac{\mathfrak m}{x_2}]$ と
$A[\frac{\mathfrak m}{x_3}]$ を計算すれば，$p$ は $X$ の対応する
開部分に含まれ得ないことが容易にわかる．
% P07-ERR-0423: frozen English omits the preposition in "contained in"; Japanese states the intended relation.
したがって $p$ は $A[\frac{\mathfrak m}{x_1}]$ のスペクトルに属する．
前と同様にブローアップすると，$p$ は座標 $y_2 = y_3 = 0$ をもつ点でなければならず，
新しい等式は
$$
y_3^2 + x_1(a y_2^2 + b y_2 y_3 + c y_3^2) \in (\mathfrak m')^4
$$
という形になる．これは前と同じ形をもち，$x_1$ が例外因子を定めるという
性質をもつ．したがって，この過程が停止しなければ無限ブローアップ列が得られ，
その各段階で $x_1$ は特異点の局所環における例外因子を定める．
ゆえに補題 \ref{resolve-lemma-sequence-blowups} と
\ref{resolve-lemma-sequence-blowups-along-arc-becomes-nonsingular}，および
前と同じ議論を用いて証明を終えることができる．

\begin{lemma}
\label{resolve-lemma-resolve-rational-double-points}
$(A, \mathfrak m, \kappa)$ を，有理特異点を定め，完備化が正規であり，
Gorenstein である次元 $2$ の局所正規 Nagata 整域とする．
このとき，特異閉点におけるブローアップの有限列
$$
X_n \to X_{n - 1} \to \ldots \to X_1 \to X_0 = \Spec(A)
$$
で，$X_n$ が正則であり，途中の各概形 $X_i$ が正規かつ
同じ型の特異点を有限個しかもたないものが存在する．
% P07-ERR-0424: frozen English has singular/plural disagreement in "each intervening schemes"; Japanese is grammatical.
\end{lemma}

\begin{proof}
これはまさに上の議論で証明したことである．
\end{proof}




\section{導かれる性質}
\label{resolve-section-existence-gives}

\noindent
本節では，Noether 整概形に正則オルタレーションが存在すると，
かなり多くの帰結が得られることを証明する．
「悪い」Noether 環に関心のない読者は，本節を飛ばしてよい．

\begin{lemma}
\label{resolve-lemma-regular-alteration-implies}
$Y$ を Noether 整概形とする．オルタレーション
$f : X \to Y$ で，$X$ が正則となるものが存在すると仮定する．
このとき正規化 $Y^\nu \to Y$ は有限であり，$Y$ は正則な稠密開部分をもつ．
\end{lemma}

\begin{proof}
$Y = \Spec(A)$ であり，$A$ が Noether 整域である場合を証明すれば十分である．
$B$ を $A$ の分数体における整閉包とする．
$C = \Gamma(X, \mathcal{O}_X)$ とおく．Cohomology of Schemes, 補題
\ref{coherent-lemma-proper-over-affine-cohomology-finite} により，
$C$ は有限 $A$-加群である．$X$ は正規なので
（Properties, 補題 \ref{properties-lemma-regular-normal}），
% P07-ERR-0425: frozen English twice omits the article in “C is a normal domain”; Japanese states the intended predicate grammatically.
$C$ は正規整域である
（Properties, 補題 \ref{properties-lemma-normal-integral-sections}）．
したがって $B \subset C$ であり，$B$ は $A$ 上有限である．
これは $A$ が Noether だからである．

\medskip\noindent
空でない開部分 $V \subset Y$ で，$f^{-1}V \to V$ が有限となるものが存在する．
Morphisms, 定義 \ref{morphisms-definition-alteration} を参照せよ．
$V$ を縮小した後，$f^{-1}V \to V$ は平坦であると仮定してよい
（Morphisms, 命題 \ref{morphisms-proposition-generic-flatness}）．
したがって $f^{-1}V \to V$ は忠実平坦である．このとき
Algebra, 補題 \ref{algebra-lemma-descent-regular} により $V$ は正則である．
\end{proof}

\begin{lemma}
\label{resolve-lemma-algebra-helper}
$(A, \mathfrak m)$ を Noether 局所環とする．$B \subset C$ を
有限 $A$-代数とする．(a) $B$ が正規環であり，かつ
(b) $\mathfrak m$-進完備化 $C^\wedge$ が正規環であると仮定する．
このとき $B^\wedge$ は正規環である．
\end{lemma}

\begin{proof}
可換図式
$$
\xymatrix{
B \ar[r] \ar[d] & C \ar[d] \\
B^\wedge \ar[r] & C^\wedge
}
$$
を考える．$\mathfrak m$-進完備化は有限 $A$-加群の圏において，
平坦 $A$-代数 $A^\wedge$ とのテンソル積で与えられるため完全であることを
思い出そう（Algebra, 補題 \ref{algebra-lemma-completion-flat}）．
Serre の判定法（Algebra, 補題 \ref{algebra-lemma-criterion-normal}）を用いて，
Noether 環 $B^\wedge$ が正規であることを証明する．
$\mathfrak q \subset B^\wedge$ を，$\mathfrak p \subset B$ の上にある素イデアルとする．
$\dim(B_\mathfrak p) \geq 2$ ならば
$\text{depth}(B_\mathfrak p) \geq 2$ であり，
$B_\mathfrak p \to B^\wedge_\mathfrak q$ は平坦なので
$\text{depth}(B^\wedge_\mathfrak q) \geq 2$ が得られる
（Algebra, 補題 \ref{algebra-lemma-apply-grothendieck}）．
$\dim(B_\mathfrak p) \leq 1$ ならば，$B_\mathfrak p$ は
離散付値環または体である．この場合 $C_\mathfrak p$ は
$B_\mathfrak p$ 上忠実平坦である（有限かつ捩れなしだからである）．
したがって $B^\wedge_\mathfrak p \to C^\wedge_\mathfrak p$ は忠実平坦であり，
$\mathfrak q$ における局所化の後も同じことが成り立つ．
$C^\wedge$ は，したがってその任意の局所化も $(S_2)$ なので，
Algebra, 補題 \ref{algebra-lemma-descent-Sk} により
$B^\wedge_\mathfrak p$ は $(S_2)$ である．以上から
$(S_2)$ は $B^\wedge$ に対して成り立つ．$B^\wedge$ が
$(R_1)$ であることを証明するには，素イデアル $\mathfrak q \subset B^\wedge$ で
$\dim(B^\wedge_\mathfrak q) \leq 1$ を満たすものだけを考えればよい．
Algebra, 補題 \ref{algebra-lemma-dimension-base-fibre-total} により
$\dim(B^\wedge_\mathfrak q) = \dim(B_\mathfrak p) +
\dim(B^\wedge_\mathfrak q/\mathfrak p B^\wedge_\mathfrak q)$
なので，$\dim(B_\mathfrak p) \leq 1$ が得られ，前と同様に
$B^\wedge_\mathfrak q \to C^\wedge_\mathfrak q$ は忠実平坦であることがわかる．
Algebra, 補題 \ref{algebra-lemma-descent-Rk} を用いて結論する．
\end{proof}

\begin{lemma}
\label{resolve-lemma-regular-alteration-implies-local}
$(A, \mathfrak m, \kappa)$ を Noether 局所整域とする．
オルタレーション $f : X \to \Spec(A)$ で，$X$ が正則となるものが
存在すると仮定する．このとき
\begin{enumerate}
% P07-ERR-0426: frozen authority reuses f for the alteration and the following ring element; retained sidecar-only.
\item 零でない $f \in A$ で，$A_f$ が正則となるものが存在する，
\item $B$ を $A$ の分数体における整閉包とすると，これは $A$ 上有限である，
\item $\mathfrak m$-進完備化によって $B$ から得られる環は正規環である，すなわち
$B$ の各極大イデアルにおける完備化は正規整域である，かつ
\item $A$ の一般形式的ファイバーは正則である．
\end{enumerate}
\end{lemma}

\begin{proof}
(1) と (2) は補題 \ref{resolve-lemma-regular-alteration-implies} から従う．
(3) の証明で記号を設定するため，同補題の証明の一部を繰り返す必要がある．
$C = \Gamma(X, \mathcal{O}_X)$ とおく．Cohomology of Schemes, 補題
\ref{coherent-lemma-proper-over-affine-cohomology-finite} により，
$C$ は有限 $A$-加群である．$X$ は正規なので
（Properties, 補題 \ref{properties-lemma-regular-normal}），
% P07-ERR-0425: second frozen missing-article occurrence; Japanese states the intended predicate grammatically.
$C$ は正規整域である
（Properties, 補題 \ref{properties-lemma-normal-integral-sections}）．
したがって $B \subset C$ であり，$B$ は $A$ 上有限である．
これは $A$ が Noether だからである．補題 \ref{resolve-lemma-algebra-helper} により，
(3) を証明するには $\mathfrak m$-進完備化 $C^\wedge$ が正規であることを
示せば十分である．

\medskip\noindent
Algebra, 補題 \ref{algebra-lemma-completion-finite-extension} により，
完備化 $C^\wedge$ は，$C$ の完備化たちの積である．ここで，各完備化は
$C$ の $\mathfrak m$ の上にある素イデアルで取る．それらは有限個で，
$\mathfrak m_1, \ldots, \mathfrak m_r$ という $C$ の極大イデアルである．
（$B$ に対する対応する結果が，補題の最後の主張を説明する．）
したがって，$A$ を $C_{\mathfrak m_i}$ で，$X$ を
$X_i = X \times_{\Spec(C)} \Spec(C_{\mathfrak m_i})$ で置き換えると，
次の段落で論じる場合に帰着する．
（Cohomology of Schemes, 補題
\ref{coherent-lemma-flat-base-change-cohomology} により
$\Gamma(X_i, \mathcal{O}) = C_{\mathfrak m_i}$ であることに注意せよ．）

\medskip\noindent
ここで $A$ は Noether 局所正規整域であり，$f : X \to \Spec(A)$ は
$\Gamma(X, \mathcal{O}_X) = A$ を満たす正則オルタレーションである．
$A^\wedge$ が，$A$ の完備化として正規整域であることを示さなければならない．
補題 \ref{resolve-lemma-port-regularity-to-completion} により
$Y = X \times_{\Spec(A)} \Spec(A^\wedge)$ は正則である．
Cohomology of Schemes, 補題
\ref{coherent-lemma-flat-base-change-cohomology} により
$\Gamma(Y, \mathcal{O}_Y) = A^\wedge$ なので，前と同様に
$A^\wedge$ は正規であると結論できる．実際，$Y$ は Properties, 補題
\ref{properties-lemma-regular-normal} により正規である．
$\Gamma(Y, \mathcal{O}_Y) = A^\wedge$ は局所環なので，これは連結である．
したがって $Y$ は正規かつ整である
（Noether 概形では連結かつ正規ならば整だからである）．ゆえに
$\Gamma(Y, \mathcal{O}_Y) = A^\wedge$ は Properties, 補題
\ref{properties-lemma-normal-integral-sections} により正規整域である．
これで (3) が証明された．

\medskip\noindent
(4) の証明．$\eta \in \Spec(A)$ で一般点を表し，下付き添字
$\eta$ で $\eta$ への基底変換を表す．$f$ はオルタレーションなので，
概形 $X_\eta$ は $\eta$ 上有限かつ忠実平坦である．
$Y = X \times_{\Spec(A)} \Spec(A^\wedge)$ は補題
\ref{resolve-lemma-port-regularity-to-completion} により正則なので，
$Y_\eta$ は正則である（$Y$ の開部分の極限だからである）．このとき
$Y_\eta \to \Spec(A^\wedge \otimes_A \kappa(\eta))$ は，一般形式的ファイバーへの
有限忠実平坦射である．Algebra, 補題 \ref{algebra-lemma-descent-regular}
により結論する．
\end{proof}











\section{特異点解消}
\label{resolve-section-resolution}



\noindent
まず定義を与える．

\begin{definition}
\label{resolve-definition-resolution}
$Y$ を Noether 整概形とする．$Y$ の {\it 特異点解消} とは，
修正 $f : X \to Y$ で $X$ が正則となるものをいう．
\end{definition}

\noindent
曲面の場合には，もう少し詳しい情報が必要になることがある．

\begin{definition}
\label{resolve-definition-resolution-surface}
$Y$ を $2$ 次元 Noether 整概形とする．次の列が存在するとき，
$Y$ は {\it 正規化ブローアップによる特異点解消} をもつという．
$$
Y_n \to Y_{n - 1} \to \ldots \to Y_1 \to Y_0 \to Y
$$
ここで
\begin{enumerate}
\item $Y_i$ は $Y$ 上固有である（$i = 0, \ldots, n$），
\item $Y_0 \to Y$ は正規化である，
\item $Y_i \to Y_{i - 1}$ は正規化ブローアップである
（$i = 1, \ldots, n$），かつ
\item $Y_n$ は正則である．
\end{enumerate}
\end{definition}

\noindent
条件 (1) から，正規化 $Y_0$ から $Y$ への射は有限，すなわち $Y$ 上有限であり，
正規化ブローアップに用いる正規化も有限であることに注意する．

\begin{lemma}
\label{resolve-lemma-existence-implies-existence-by-normalized-blowing-ups}
$(A, \mathfrak m, \kappa)$ を Noether 局所環とする．
$A$ は正規で次元 $2$ であると仮定する．
$\Spec(A)$ が特異点解消をもつならば，
$\Spec(A)$ は正規化ブローアップによる解消をもつ．
\end{lemma}

\begin{proof}
補題 \ref{resolve-lemma-regular-alteration-implies-local} により，
完備化 $A^\wedge$ は $A$ に対して正規である．
補題 \ref{resolve-lemma-port-regularity-to-completion} により，
$\Spec(A^\wedge)$ は特異点解消をもつ．
補題 \ref{resolve-lemma-normalized-blowup-completion} によれば，
% P07-ERR-0427: frozen English has a stray “of” before “comes”; Japanese states the intended provenance.
正規化ブローアップの任意の列
$Y_n \to Y_{n - 1} \to \ldots \to \Spec(A^\wedge)$ は，
正規化ブローアップの列 $X_n \to \ldots \to \Spec(A)$ から生じる．
さらに $Y_n$ が正則ならば，補題
\ref{resolve-lemma-port-regularity-to-completion} により $X_n$ は正則である．
したがって $A$ が完備である場合に補題を証明すれば十分である．

\medskip\noindent
% P07-ERR-0428: frozen English says “A is a complete”; Japanese states the predicate grammatically.
さらに $A$ は完備であると仮定する．$A$ が Nagata
（Algebra, 命題 \ref{algebra-proposition-ubiquity-nagata}），
優秀（More on Algebra, 命題
\ref{more-algebra-proposition-ubiquity-excellent}）であり，
双対化複体をもつこと
（Dualizing Complexes, 補題 \ref{dualizing-lemma-ubiquity-dualizing}）を用いる．
さらに，$A$ 上本質的有限型な任意の環についても同じことが成り立つ．
% P07-ERR-0429: frozen English twice uses “a excellent”; Japanese uses the grammatical predicate.
$B$ が優秀局所正規整域ならば，完備化 $B^\wedge$ は正規である
（$B \to B^\wedge$ は正則であり，More on Algebra, 補題
\ref{more-algebra-lemma-normal-goes-up} が適用できるからである）．
以下の証明では，これを断りなく用いる．

\medskip\noindent
$X \to \Spec(A)$ を特異点解消とする．
正規化ブローアップの列
$$
Y_n \to Y_{n - 1} \to \ldots \to Y_1 \to \Spec(A)
$$
で $X$ を支配するものを選ぶ
（補題 \ref{resolve-lemma-dominate-by-normalized-blowing-up}）．
射 $Y_n \to X$ は $X$ の有限個の点を除いて同型である．
したがって補題 \ref{resolve-lemma-dominate-by-blowing-up-in-points} を適用すると，
閉点におけるブローアップの列
$$
X_m \to X_{m - 1} \to \ldots \to X
$$
で，$X_m$ が $Y_n$ を支配するものが得られる．図式は
$$
\xymatrix{
& Y_n \ar[rd] \ar[rr] & & \Spec(A) \\
X_m \ar[rr] \ar[ru] & & X \ar[ru]
}
$$
である．補題を証明するには，$Y_n$ に有限回の正規化ブローアップを
施すと正則概形が得られることを示せば十分である．上の図式から，
$Y_n$ は特異点解消（すなわち $X_m$）をもつ．$Y_n$ は正規曲面なので，
$Y_n$ の特異点は高々有限個
$y_1, \ldots, y_t$ である
（$X_m \to Y_n$ は次元 $1$ のファイバーを除いて同型だからである．
Varieties, 補題
\ref{varieties-lemma-modification-normal-iso-over-codimension-1} を参照せよ）．

\medskip\noindent
$x_a \in X$ を $y_a$ の像とする．このとき $\mathcal{O}_{X, x_a}$ は
正則であり，したがって有理特異点を定める
（補題 \ref{resolve-lemma-regular-rational}）．
補題 \ref{resolve-lemma-rational-propagates} を
$\mathcal{O}_{X, x_a} \to \mathcal{O}_{Y_n, y_a}$ に適用すると，
$\mathcal{O}_{Y_n, y_a}$ が有理特異点を定めることがわかる．
補題 \ref{resolve-lemma-rational-to-gorenstein} により，
特異閉点におけるブローアップの有限列
$$
Y_{a, n_a} \to Y_{a, n_a - 1} \to \ldots \to \Spec(\mathcal{O}_{Y_n, y_a})
$$
で，$Y_{a, n_a}$ が Gorenstein，すなわち可逆な双対化加群をもつものが
存在する．（本質的には自明な）補題
\ref{resolve-lemma-equivalence-sequence-blowups} により，
$n' = \sum n_a$ とおくと，これらの列はブローアップの列
$$
Y_{n + n'} \to Y_{n + n' - 1} \to \ldots \to Y_n
$$
に対応し，$Y_{n + n'}$ は正規で，$Y_{n + n'}$ の局所環は
Gorenstein である．上で挙げた結果を用いると，$Y_{n + n'}$ を，
ブローアップの列 $X_{m + m'} \to \ldots \to X_m$ で支配でき，
この列は $Y_{n + n'}$ を次の図式のように支配する．
$$
\xymatrix{
 & Y_{n + n'} \ar[rr] & & Y_n \ar[rd] \ar[rr] & & \Spec(A) \\
X_{m + m'} \ar[ru] \ar[rr] & & X_m \ar[rr] \ar[ru] & & X \ar[ru]
}
$$
したがって再び $Y_{n + n'}$ は有限個の特異点
$y'_1, \ldots, y'_s$ をもつが，今回はそれらは有理二重点である．
より正確には，局所環 $\mathcal{O}_{Y_{n + n'}, y'_b}$ は
補題 \ref{resolve-lemma-resolve-rational-double-points} のものとなる．
上とまったく同様に論じれば，補題が成り立つと結論できる．
\end{proof}

\begin{lemma}
\label{resolve-lemma-resolve-complete}
$(A, \mathfrak m, \kappa)$ を Noether 完備局所環とする．
$A$ は次元 $2$ の正規整域であると仮定する．このとき
$\Spec(A)$ は特異点解消をもつ．
\end{lemma}

\begin{proof}
Noether 完備局所環は J-2
（More on Algebra, 命題 \ref{more-algebra-proposition-ubiquity-J-2}），
Nagata（Algebra, 命題 \ref{algebra-proposition-ubiquity-nagata}），
優秀（More on Algebra, 命題
\ref{more-algebra-proposition-ubiquity-excellent}）であり，
双対化複体をもつ
（Dualizing Complexes, 補題 \ref{dualizing-lemma-ubiquity-dualizing}）．
さらに，$A$ 上本質的有限型な任意の環についても同じことが成り立つ．
% P07-ERR-0429: second frozen “a excellent” occurrence; Japanese states the predicate grammatically.
$B$ が優秀局所正規整域ならば，完備化 $B^\wedge$ は正規である
（$B \to B^\wedge$ は正則であり，More on Algebra, 補題
\ref{more-algebra-lemma-normal-goes-up} が適用できるからである）．
% P07-ERR-0430: frozen English repeats “required”; Japanese states the requirement once.
言い換えれば，以下の証明で現れる局所環は，必要な「優秀性」の性質をもつ．

\medskip\noindent
$A_0 \subset A$ を，$A_0$ が正則完備局所環であり，
$A_0 \to A$ が有限となるように選ぶ．Algebra, 補題
\ref{algebra-lemma-complete-local-Noetherian-domain-finite-over-regular}
を参照せよ．これは分数体の有限拡大 $K/K_0$ を誘導する．
$[K : K_0]$ に関する帰納法で論じる．次数が $1$ の場合が基底であり，
このとき $A_0 = A$ なので結果は成り立つ．

\medskip\noindent
中間体 $K_0 \subset L \subset K$ で
$K_0 \not = L \not = K$ を満たすものがあるとする．$B \subset A$ を
$A_0$ の $L$ における整閉包とする．帰納法により特異点解消
$Y \to \Spec(B)$ を選ぶ．$X$ を
$Y \times_{\Spec(B)} \Spec(A)$ の正規化とする．図式は
$$
\xymatrix{
X \ar[r] \ar[d] & \Spec(A) \ar[d] \\
Y \ar[r] & \Spec(B)
}
$$
である．$A$ は J-2 なので，$X$ の正則軌跡は開である．
$X$ は正規曲面なので，$X$ は高々有限個の特異点
$x_1, \ldots, x_n$ しかもたず，それらは
$\dim(\mathcal{O}_{X, x_i}) = 2$ を満たす閉点である．
各 $i$ に対し，その像を $y_i \in Y$ とする．
$\mathcal{O}_{Y, y_i}^\wedge \to \mathcal{O}_{X, x_i}^\wedge$
は以前より小さい次数の有限射なので，帰納法の仮定により
$\mathcal{O}_{X, x_i}^\wedge$ は特異点解消をもつ．補題
\ref{resolve-lemma-existence-implies-existence-by-normalized-blowing-ups} により，
正規化ブローアップの列
$$
Z^\wedge_{i, n_i} \to \ldots \to Z^\wedge_{i, 1} \to
\Spec(\mathcal{O}_{X, x_i}^\wedge)
$$
で，$Z^\wedge_{i, n_i}$ が正則となるものが存在する．
補題 \ref{resolve-lemma-normalized-blowup-completion} により，
対応する正規化ブローアップの列
$$
Z_{i, n_i} \to \ldots \to Z_{i, 1} \to \Spec(\mathcal{O}_{X, x_i})
$$
が存在する．このとき $Z_{i, n_i}$ は補題
\ref{resolve-lemma-port-regularity-to-completion} により正則概形である．
補題 \ref{resolve-lemma-equivalence-sequence-normalized-blowups} により，
これらの正規化ブローアップを対応する列
$$
Z_n \to Z_{n - 1} \to \ldots \to Z_1 \to X
$$
に組み込める．もちろん，局所環を見れば $Z_n$ も正則である．
これで帰納段階が証明された．

\medskip\noindent
中間体 $K_0 \subset L \subset K$ で
$K_0 \not = L \not = K$ を満たすものがないと仮定する．
このとき $K/K_0$ は分離的であるか，または
% P07-ERR-0431: frozen English says “characteristic to K”; Japanese states that K has characteristic p.
$K$ の標数が $p$ で $[K : K_0] = p$ である．
すると補題 \ref{resolve-lemma-go-up-separable} または
\ref{resolve-lemma-go-up-degree-p} により，有理特異点への帰着が可能である．
補題 \ref{resolve-lemma-reduce-to-rational} により，正規修正
$X \to \Spec(A)$ で，任意の特異点 $x$ に対して $X$ の局所環
$\mathcal{O}_{X, x}$ が有理特異点を定めるものが存在すると結論できる．
$A$ は J-2 なので，$X$ は有限個の特異点
$x_1, \ldots, x_n$ をもつ．補題 \ref{resolve-lemma-rational-to-gorenstein} により，
特異閉点におけるブローアップの有限列
$$
X_{i, n_i} \to X_{i, n_i - 1} \to \ldots \to \Spec(\mathcal{O}_{X, x_i})
$$
で，$X_{i, n_i}$ が Gorenstein，すなわち可逆な双対化加群をもつものが存在する．
（本質的には自明な）補題 \ref{resolve-lemma-equivalence-sequence-blowups} により，
% P07-ERR-0432: frozen authority uses undefined summation index a instead of i; formula retained unchanged.
$n = \sum n_a$ とおくと，これらの列はブローアップの列
$$
X_n \to X_{n - 1} \to \ldots \to X
$$
に対応し，$X_n$ は正規で，$X_n$ の局所環は Gorenstein である．
再び $X_n$ は有限個の特異点 $x'_1, \ldots, x'_s$ をもつが，
今回はそれらは有理二重点である．より正確には，局所環
$\mathcal{O}_{X_n, x'_i}$ は補題
\ref{resolve-lemma-resolve-rational-double-points} のものとなる．
上とまったく同様に論じれば，補題が成り立つと結論できる．
\end{proof}

\noindent
最後に本章の主定理に到達する．

\begin{theorem}[Lipman]
\label{resolve-theorem-resolve}
\begin{reference}
\cite[151ページの定理]{Lipman}
\end{reference}
$Y$ を 2 次元 Noether 整概形とする．次の条件は同値である．
\begin{enumerate}
\item オルタレーション $X \to Y$ で $X$ が正則となるものが存在する，
\item $Y$ の特異点解消が存在する，
\item $Y$ は正規化ブローアップによる特異点解消をもつ，
\item 正規化 $Y^\nu \to Y$ は有限であり，$Y^\nu$ は有限個の特異点
$y_1, \ldots, y_m$ をもち，各 $y_i$ に対して
$\mathcal{O}_{Y^\nu, y_i}$ の完備化は正規である．
\end{enumerate}
\end{theorem}

\begin{proof}
含意 (3) $\Rightarrow$ (2) $\Rightarrow$ (1) は直ちに従う．

\medskip\noindent
オルタレーション $X \to Y$ で $X$ が正則となるものをとる．補題
\ref{resolve-lemma-regular-alteration-implies} により $Y^\nu \to Y$ は有限である．
Morphisms, 補題 \ref{morphisms-lemma-normalization-normal} による分解
$f : X \to Y^\nu$ を考える．射 $f$ は開集合 $V \subset Y^\nu$ 上有限であり，
Varieties, 補題 \ref{varieties-lemma-finite-in-codim-1} により，この開集合は
余次元が $\leq 1$ である $Y^\nu$ のすべての点を含む．次に $f$ は $V$ 上平坦である．
これは Algebra, 補題 \ref{algebra-lemma-CM-over-regular-flat} と，
Serre の判定法により次元 $\leq 2$ の正規局所環が Cohen--Macaulay
であること（Algebra, 補題 \ref{algebra-lemma-criterion-normal}）から従う．
したがって Algebra, 補題
\ref{algebra-lemma-descent-regular} により $V$ は正則である．
$Y^\nu$ は Noether なので，
$Y^\nu \setminus V = \{y_1, \ldots, y_m\}$ は有限である．補題
\ref{resolve-lemma-regular-alteration-implies-local} により
$\mathcal{O}_{Y^\nu, y_i}$ の完備化は正規である．
以上から (1) $\Rightarrow$ (4) が従う．

\medskip\noindent
(4) を仮定する．(3) を証明すればよい．ただちに $Y$ をその正規化で
置き換えてよい．特異点を $y_1, \ldots, y_m \in Y$ とする．補題
\ref{resolve-lemma-resolve-complete} および
\ref{resolve-lemma-existence-implies-existence-by-normalized-blowing-ups} を適用する．
% P07-ERR-0433: frozen English has the malformed phrase “we find there exists”; Japanese states the intended existence claim.
すると正規化ブローアップの有限列
$$
Y_{i, n_i} \to Y_{i, n_i - 1} \to \ldots \to \Spec(\mathcal{O}^\wedge_{Y, y_i})
$$
で $Y_{i, n_i}$ が正則となるものが存在する．補題
\ref{resolve-lemma-normalized-blowup-completion} により，対応する正規化ブローアップの列
$$
X_{i, n_i} \to \ldots \to X_{i, 1} \to \Spec(\mathcal{O}_{Y, y_i})
$$
が存在する．このとき補題 \ref{resolve-lemma-port-regularity-to-completion} により
$X_{i, n_i}$ は正則概形である．補題
\ref{resolve-lemma-equivalence-sequence-normalized-blowups} により，
これらの正規化ブローアップを対応する列
$$
X_n \to X_{n - 1} \to \ldots \to X_1 \to Y
$$
に組み込める．もちろん，局所環を見れば $X_n$ も正則である．
これで証明が完了した．
\end{proof}






\section{埋め込み特異点解消}
\label{resolve-section-embedded}

\noindent
曲面上の曲線が与えられたとき，その曲線を単純正規交差因子にする
ブローアップが存在する．本節では，1 次元局所 Noether 概形が正規であることと
正則であることが同値であるという事実
（Algebra, 補題 \ref{algebra-lemma-characterize-dvr}）を用いる．
また，局所 Noether 概形の任意の点がある閉点に特殊化するという事実
（Properties, 補題 \ref{properties-lemma-locally-Noetherian-closed-point}）
も用いる．

\begin{lemma}
\label{resolve-lemma-resolve-curve}
$Y$ を 1 次元 Noether 整概形とする．次の条件は同値である．
\begin{enumerate}
\item オルタレーション $X \to Y$ で $X$ が正則となるものが存在する，
\item $Y$ の特異点解消が存在する，
\item 閉点におけるブローアップの有限列
$Y_n \to Y_{n - 1} \to \ldots \to Y_1 \to Y$ で
$Y_n$ が正則となるものが存在する，かつ
\item 正規化 $Y^\nu \to Y$ は有限である．
\end{enumerate}
\end{lemma}

\begin{proof}
含意 (3) $\Rightarrow$ (2) $\Rightarrow$ (1) は直ちに従う．
含意 (1) $\Rightarrow$ (4) は補題
\ref{resolve-lemma-regular-alteration-implies} から従う．
正規な 1 次元概形は正則であるから，含意 (4) $\Rightarrow$ (2) も明らかである．
したがって，同値な条件 (1)，(2)，(4) が (3) を導くことを示せばよい．

\medskip\noindent
$f : X \to Y$ を特異点解消とする．$Y$ の次元は 1 なので，
Varieties, 補題 \ref{varieties-lemma-finite-in-codim-1} により
$f$ は有限である．次の分解を構成する．
$$
X \to \ldots \to Y_2 \to Y_1 \to Y
$$
ここで，$Y_i \to Y_{i - 1}$ は，$Y_{i - 1}$ が正則でないかぎり，
閉点におけるブローアップであり，同型ではない．
これらの射は（上と同じ理由により）すべて有限となり，対応する系
$$
f_*\mathcal{O}_X \supset \ldots \supset
f_{2, *}\mathcal{O}_{Y_2} \supset
f_{1, *}\mathcal{O}_{Y_1} \supset \mathcal{O}_Y
$$
を得る．ここで $f_i : Y_i \to Y$ は構造射である．
$Y$ は Noether なので，この連接部分加群の増大列は安定する
（Cohomology of Schemes, 補題 \ref{coherent-lemma-acc-coherent}）．
したがって，ある $n$ に対して $Y_n$ は所望のとおり正則である．
$Y_i$ を $Y_{i - 1}$ から構成するには，特異閉点
$y_{i - 1} \in Y_{i - 1}$ を選び，$Y_i \to Y_{i - 1}$ を
対応するブローアップとする．$X$ は正則で次元 $1$ である
（したがって閉点における局所環は離散付値環，特に単項イデアル整域である）ので，
イデアル層 $\mathfrak m_{y_{i - 1}} \cdot \mathcal{O}_X$ は可逆である．
ブローアップの普遍性（Divisors, 補題
\ref{divisors-lemma-universal-property-blowing-up}）により，これは分解
$X \to Y_i$ を与える．最後に，$Y_i \to Y_{i - 1}$ は同型ではない．
実際，$\mathfrak m_{y_{i - 1}}$ は可逆イデアルではない．
\end{proof}

\begin{lemma}
\label{resolve-lemma-blowup-curve}
$X$ を Noether 概形とする．$Y \subset X$ を，補題
\ref{resolve-lemma-resolve-curve} の同値な条件を満たす次元 $1$ の整閉部分概形とする．
このとき閉点におけるブローアップの有限列
$$
X_n \to X_{n - 1} \to \ldots \to X_1 \to X
$$
で，$Y$ の狭義変換が $X_n$ において正則曲線となるものが存在する．
\end{lemma}

\begin{proof}
補題 \ref{resolve-lemma-resolve-curve} により与えられるブローアップの列を
$Y_n \to Y_{n - 1} \to \ldots \to Y_1 \to Y$ とする．対応する列を
$X_n \to X_{n - 1} \to \ldots \to X_1 \to X$ とする．これは
$X$ のブローアップの列である．
Divisors, 補題 \ref{divisors-lemma-strict-transform} により，
狭義変換はブローアップだから，これでよい．
\end{proof}

\noindent
$X$ を局所 Noether 概形とし，$Y, Z \subset X$ を閉部分概形とする．
$p \in Y \cap Z$ を閉点とする．$Y$ は次元 $1$ の整概形であり，
$Y$ の生成点は $Z$ に含まれないと仮定する．この状況では不変量
\begin{equation}
\label{resolve-equation-multiplicity}
m_p(Y \cap Z) =
\text{length}_{\mathcal{O}_{X, p}}(\mathcal{O}_{Y \cap Z, p})
\end{equation}
を考えることができる．これは整数 $\geq 1$ である．実際，
$I, J \subset \mathcal{O}_{X, p}$ を $Y, Z$ に対応するイデアルとすると，
% P07-ERR-0434: the frozen quotient lacks parentheses around I + J; the formula is retained unchanged.
$\mathcal{O}_{Y \cap Z, p} = \mathcal{O}_{X, p}/I + J$
の台は $\{\mathfrak m_p\}$ に等しい．これは，$Y \cap Z$ が，
$Y$ の点で $p$ に特殊化する唯一のものを含まないと仮定したからである．
したがって Algebra, 補題 \ref{algebra-lemma-support-point} により
長さは有限である．

\begin{lemma}
\label{resolve-lemma-blowup-nonsingular-curves-meeting-at-point}
上の状況で，$X' \to X$ を $X$ の $p$ におけるブローアップとする．
$Y', Z' \subset X'$ を $Y, Z$ の狭義変換とする．
$\mathcal{O}_{Y, p}$ が正則ならば，次が成り立つ．
\begin{enumerate}
\item $Y' \to Y$ は同型である，
\item $Y'$ は例外ファイバー $E \subset X'$ とただ一つの点 $q$ で交わり，
% P07-ERR-0435: the frozen statement uses Y rather than its strict transform Y'; formula retained unchanged.
$m_q(Y \cap E) = 1$ である，
\item $q \in Z'$ でもあるならば，
% P07-ERR-0435: the second frozen invariant again uses Y rather than Y'; formula retained unchanged.
$m_q(Y \cap Z') < m_p(Y \cap Z)$ である．
\end{enumerate}
\end{lemma}

\begin{proof}
$\mathcal{O}_{X, p} \to \mathcal{O}_{Y, p}$ は全射であり，
$\mathcal{O}_{Y, p}$ は離散付値環なので，その極大イデアルの元
$x_1 \in \mathfrak m_p$ で，$\mathcal{O}_{Y, p}$ の一意化元に写るものを
選べる．アフィン開集合 $U = \Spec(A)$ で $p$ を含み，$x_1 \in A$ と
なるものを選ぶ．$\mathfrak m \subset A$ を $p$ に対応する極大イデアルとし，
$I, J \subset A$ を $Y, Z$ を $\Spec(A)$ において定めるイデアルとする．
$U$ を縮小すれば，$\mathfrak m = I + (x_1)$，言い換えれば概形論的に
$V(x_1) \cap U \cap Y = \{p\}$ と仮定できる．
$p$ は $Y$ 上の有効 Cartier 因子である．さらに $Y'$ は $Y$ の $p$ における
ブローアップなので（Divisors, 補題 \ref{divisors-lemma-strict-transform}），
Divisors, 補題 \ref{divisors-lemma-blow-up-effective-Cartier-divisor} により
$Y' \to Y$ は同型である．関係 $\mathfrak m = I + (x_1)$ から
$\mathfrak m^n \subset I + (x_1^n)$ が従うので，写像
$$
\psi : A[\textstyle{\frac{\mathfrak m}{x_1}}] \longrightarrow A/I
$$
を次のように定義できる．$y/x_1^n \in A[\frac{\mathfrak m}{x_1}]$ に対し，
$a$ の $A/I$ における類を対応させる．ここで $a$ は
$y \equiv ax_1^n \bmod I$ を満たすように選ぶ．
すると $\psi$ は，$Y \cap U$ から $X'$ への $U$ 上の射で，
$Y' \cong Y$ が与えるものに対応する．$x_1$ の
$A[\frac{\mathfrak m}{x_1}]$ における像は例外因子を切り出すので，
% P07-ERR-0436: the frozen proof writes m_q(Y', E) instead of an intersection; notation retained unchanged.
$m_q(Y', E) = 1$ と結論できる．最後に，$J \subset \mathfrak m$ だから，
イデアル $J' \subset A[\frac{\mathfrak m}{x_1}]$ は確かに，
元 $f/x_1$（$f \in J$）を含む．そこで $f \in J$ を，その像
$\overline{f}$ の $A/I$ における最小付値が $m_p(Y \cap Z)$ に
等しくなるように選ぶと，$\psi(f/x_1) = \overline{f}/x_1$ の
$A/I$ における付値は一つ小さい．これで補題の最後の主張が証明された．
\end{proof}

\begin{lemma}
\label{resolve-lemma-blowup-curves}
$X$ を Noether 概形とする．
% P07-ERR-0437: frozen English combines a singular article with the plural family; Japanese states the family grammatically.
$Y_i \subset X$（$i = 1, \ldots, n$）を，それぞれ次元 $1$ の整閉部分概形で，
補題 \ref{resolve-lemma-resolve-curve} の同値な条件を満たすものとする．
このとき閉点におけるブローアップの有限列
$$
X_n \to X_{n - 1} \to \ldots \to X_1 \to X
$$
% P07-ERR-0438: frozen English uses a singular subject with a plural predicate; Japanese states the indexed family.
で，狭義変換 $Y'_i \subset X_n$，すなわち $Y_i$ の $X_n$ における
狭義変換が互いに素な正則曲線となるものが存在する．
\end{lemma}

\begin{proof}
補題 \ref{resolve-lemma-blowup-curve} により，$Y_i$ は
$i = 1, \ldots, n$ に対して正則曲線であると仮定してよい．
各 $i \not = j$ と $p \in Y_i \cap Y_j$ に対して，不変量
$m_p(Y_i \cap Y_j)$（\ref{resolve-equation-multiplicity}）がある．
これらの数の最大値が $> 1$ ならば，最大値をとるすべての点 $p$ で
ブローアップすることにより，その値を減少させられる
（補題 \ref{resolve-lemma-blowup-nonsingular-curves-meeting-at-point}）．
最大値が $1$ ならば，同じ補題を用い，これらすべての点 $p$ で
ブローアップすることにより曲線を分離できる．
\end{proof}

\noindent
% P07-ERR-0439: frozen English says “contained on” rather than “contained in”; Japanese uses ordinary inclusion phrasing.
曲線が正則曲面に含まれるとき，しばしばこれを正規交差因子にしたい．

\begin{lemma}
\label{resolve-lemma-turn-into-effective-Cartier}
$X$ を次元 $2$ の正則概形とし，$Z \subset X$ を真の閉部分概形とする．
閉点におけるブローアップの列
$$
X_n \to \ldots \to X_1 \to X
$$
で，逆像 $Z_n$，すなわち $Z$ の $X_n$ における逆像が
有効 Cartier 因子となるものが存在する．
\end{lemma}

\begin{proof}
$D \subset Z$ を，$Z$ に含まれる最大の有効 Cartier 因子とする．
すると $\mathcal{I}_Z \subset \mathcal{I}_D$ であり，Divisors, 補題
\ref{divisors-lemma-codim-1-part} によりその商の台は閉点からなる．
したがって $\mathcal{I}_Z = \mathcal{I}_{Z'} \mathcal{I}_D$ と書ける．
ここで $Z' \subset X$ は，集合論的には有限個の閉点からなる閉部分概形である．
補題 \ref{resolve-lemma-make-ideal-principal} を適用すると，本補題の主張にある
ブローアップの列で，$\mathcal{I}_{Z'}\mathcal{O}_{X_n}$ が可逆となるものが得られる．
これで補題が証明された．
\end{proof}

\begin{lemma}
\label{resolve-lemma-embedded-resolution}
$X$ を次元 $2$ の正則概形とし，$Z \subset X$ を真の閉部分概形とする．
各既約成分 $Y \subset Z$ で次元 $1$ のものは補題
\ref{resolve-lemma-resolve-curve} の同値な条件を満たすと仮定する．
このとき閉点におけるブローアップの列
$$
X_n \to \ldots \to X_1 \to X
$$
で，逆像 $Z_n$，すなわち $Z$ の $X_n$ における逆像が，
単純正規交差因子に台をもつ
有効 Cartier 因子となるものが存在する．
\end{lemma}

\begin{proof}
$X' \to X$ を閉点 $p$ におけるブローアップとする．すると
$Z' \subset X'$，すなわち $Z$ の逆像の台は，$Z$ の狭義変換と
例外因子との和に含まれる．例外因子は正則曲線であり
（補題 \ref{resolve-lemma-blowup}），狭義変換 $Y'$ は，各既約成分 $Y$ に対して，
$Y$ 自身または $Y$ の $p$ におけるブローアップである．
したがってこの過程では，次元 $1$ の特異成分を新たに生じない．
ゆえに補題 \ref{resolve-lemma-turn-into-effective-Cartier} および
\ref{resolve-lemma-blowup-curves} により，$Z$ は有効 Cartier 因子であり，
既約成分 $Y$ は $Z$ のものすべてについて正則であると仮定してよい．
（もちろん，既約成分が互いに素であるとは仮定できない．
$Z$ の点をブローアップするたびに，例外因子という新しい既約成分を
$Z$ に加えるからである．）

\medskip\noindent
$Z$ は有効 Cartier 因子で，その既約成分 $Y_i$ は正則であると仮定する．
各 $i \not = j$ と $p \in Y_i \cap Y_j$ に対して，不変量
$m_p(Y_i \cap Y_j)$（\ref{resolve-equation-multiplicity}）がある．
これらの数の最大値が $> 1$ ならば，最大値をとるすべての点 $p$ で
ブローアップすることにより，その値を減少させられる
（補題 \ref{resolve-lemma-blowup-nonsingular-curves-meeting-at-point}．
「新しい」不変量 $m_{q_i}(Y'_i \cap E)$ は常に $1$ であることに注意せよ）．
最大値が $1$ であるとする．例えば，
$p \in Y_1 \cap \ldots \cap Y_r$ で，ある $r > 2$ に対して
他の成分には属さないならば，
$p$ をブローアップした後，$Y'_1, \ldots, Y'_r$ は $p$ の上の点では交わらず，
$m_{q_i}(Y'_i, E) = 1$ である．ここで $Y'_i \cap E = \{q_i\}$ である．
% P07-ERR-0440: frozen authority says more than 3 components, leaving triple intersections; numeral retained unchanged.
したがって，$3$ 個より多くの成分が交わる $Z$ の点を引き続き
ブローアップすると，任意の閉点 $p \in X$ において次のいずれかが
成り立つ状況に到達する：
(a) 曲線 $Y_i$ で $p$ を通るものはない；
(b) 曲線 $Y_i$ がちょうど一つ $p$ を通り，$\mathcal{O}_{Y_i, p}$ は正則である；
または (c) 曲線 $Y_i$, $Y_j$ がちょうど二つ $p$ を通り，
局所環 $\mathcal{O}_{Y_i, p}$, $\mathcal{O}_{Y_j, p}$ は正則で，
$m_p(Y_i \cap Y_j) = 1$ である．
これは $\sum Y_i$ が正則曲面 $X$ 上の単純正規交差因子であることを意味する．
\'Etale Morphisms, 補題 \ref{etale-lemma-strict-normal-crossings} を参照せよ．
\end{proof}






\section{例外曲線の収縮}
\label{resolve-section-minus-one}

\noindent
$X$ を Noether 概形とする．$E \subset X$ を，次の性質をもつ閉部分概形とする．
\begin{enumerate}
\item $E$ は $X$ 上の有効 Cartier 因子である，
\item 体 $k$ と概形の同型 $\mathbf{P}^1_k \to E$ が存在する，
\item 法層 $\mathcal{N}_{E/X}$ の引き戻しは
$\mathcal{O}_{\mathbf{P}^1}(-1)$ である．
\end{enumerate}
このような閉部分概形を {\it 第一種例外曲線} と呼ぶ．

\medskip\noindent
$X'$ を Noether 概形とし，$x \in X'$ を，$\mathcal{O}_{X', x}$ が
次元 $2$ の正則局所環となる閉点とする．$b : X \to X'$ を
$X'$ の $x$ におけるブローアップとする．この場合，例外ファイバー
$E \subset X$ は第一種例外曲線である．これは補題 \ref{resolve-lemma-blowup} から従う．

\medskip\noindent
問：任意の第一種例外曲線は，上のようなブローアップのファイバーとして
得られるだろうか．言い換えれば，概形の固有射 $X \to X'$ で，
$E$ が閉点 $x \in X'$ に写り，$\mathcal{O}_{X', x}$ が次元 $2$ の
正則局所環であり，$X$ が $X'$ の $x$ におけるブローアップとなるものは
常に存在するだろうか．これが成り立つとき，
{\it $E$ の収縮が存在する} という．

\begin{lemma}
\label{resolve-lemma-factor-through-contraction}
$X$ を Noether 概形とし，$E \subset X$ を第一種例外曲線とする．
$X \to X'$ が $E$ の収縮として存在するならば，それは次の普遍性をもつ：
$\varphi : X \to Y$ を，$\varphi(E)$ が一点となる任意の射とすると，
一意な分解 $X \to X' \to Y$ により $\varphi$ が分解する．
\end{lemma}

\begin{proof}
$b : X \to X'$ を $E$ の収縮とする．位相空間としての $X'$ は，
$X$ において $E$ のすべての点を一点に同一視する関係による商である．
実際，$b$ は固有
（Divisors, 補題 \ref{divisors-lemma-blowing-up-projective} および
Morphisms, 補題 \ref{morphisms-lemma-locally-projective-proper}）かつ全射なので，
位相空間の商写像を定める
（Topology, 補題 \ref{topology-lemma-closed-morphism-quotient-topology}）．
一方，標準写像 $\mathcal{O}_{X'} \to b_*\mathcal{O}_X$ は同型である．
実際，これは像点 $x \in X'$，すなわち $E$ の像の補集合上では明らかであり，
$x$ における茎では補題 \ref{resolve-lemma-cohomology-of-blowup} の (4) により同型である．
したがって対 $(X', \mathcal{O}_{X'})$ は，$X$ の位相空間としての商をとり，
その上に $b_*\mathcal{O}_X$ を構造層として与えることで構成される．

\medskip\noindent
$\varphi$ が与えられたとき，$\varphi' : X' \to Y$ を，
$\varphi = \varphi' \circ b$ を満たす一意な位相空間の写像とする．すると写像
$$
\varphi^\sharp : \varphi^{-1}\mathcal{O}_Y =
b^{-1}((\varphi')^{-1}\mathcal{O}_Y) \to \mathcal{O}_X
$$
は写像
$$
(\varphi')^\sharp :
(\varphi')^{-1}\mathcal{O}_Y \to b_*\mathcal{O}_X = \mathcal{O}_{X'}
$$
と随伴する．すると $(\varphi', (\varphi')^\sharp)$ は $X'$ から $Y$ への
環付き空間の射であり，所望の分解を与える．$\varphi$ は局所環付き空間の
射なので，$\varphi'$ もそうである．実際，確認すべきことは写像
$\mathcal{O}_{Y, y} \to \mathcal{O}_{X', x}$ が局所的であることだけである．
ここで $y \in Y$ は $E$ の $\varphi$ による像である．これは次の理由で成り立つ．
元 $f \in \mathfrak m_y$ の $X$ への引き戻しは $E$ のすべての点で
零となる関数である．したがって $f$ の $X'$ への引き戻しは，
$x$ の $X'$ における近傍上で定義された，同じ性質をもつ関数である．
ゆえにこの関数は所望のとおり $x$ で消える．
\end{proof}

\begin{lemma}
\label{resolve-lemma-contraction-unique}
$X$ を Noether 概形とし，$E \subset X$ を第一種例外曲線とする．
$E$ の収縮が存在するならば，一意な同型を除いて一意である．
\end{lemma}

\begin{proof}
これは補題 \ref{resolve-lemma-factor-through-contraction} の普遍性から直ちに従う．
\end{proof}

\begin{lemma}
\label{resolve-lemma-exceptional-first-kind-local}
$X$ を Noether 概形とし，$E \subset X$ を第一種例外曲線とする．
$E_n = nE$ とおき，その構造層を $\mathcal{O}_n$ と書く．このとき
$$
A = \lim H^0(E_n, \mathcal{O}_n)
$$
は次元 $2$ の完備 Noether 正則局所環であり，
$\Ker(A \to H^0(E_n, \mathcal{O}_n))$ はその極大イデアルの
$n$ 乗である．
\end{lemma}

\begin{proof}
同型 $\mathbf{P}^1_k \to E$ で，$E$ の $X$ における法層の引き戻しが
$\mathcal{O}(-1)$ となるものが存在することを思い出そう．すると
$H^0(E, \mathcal{O}_E) = k$ である．$\mathcal{O}_n(iE)$ で可逆層
$\mathcal{O}_X(iE)$ の $E_n$ への制限を表すことにする．ここでこれは
すべての $n \geq 1$ と $i \in \mathbf{Z}$ に対する記法である．
$\mathcal{O}_X(-nE)$ は $E_n$ のイデアル層であることを思い出そう．
したがって $d \geq 0$ に対し短完全列
$$
0 \to \mathcal{O}_E(-(d + n)E) \to
\mathcal{O}_{n + 1}(-dE) \to
\mathcal{O}_n(-dE) \to 0
$$
を得る．$\mathcal{O}_E(-(d + n)E) = \mathcal{O}_{\mathbf{P}^1_k}(d + n)$
なので，すべての $d \geq 0$ と $n \geq 1$ に対して第一コホモロジー群は
消える．したがって系 $H^0(E_n, \mathcal{O}_n(-dE))$ の遷移写像は全射である．
$d = 0$ とすれば，各遷移写像の核が冪零イデアルとなる環の全射の逆系を得る．
ゆえに $A = \lim H^0(E_n, \mathcal{O}_n)$ は剰余体 $k$ をもつ局所環であり，
その極大イデアルは
$$
\lim \Ker(H^0(E_n, \mathcal{O}_n) \to H^0(E, \mathcal{O}_E)) =
\lim H^0(E_n, \mathcal{O}_n(-E))
$$
である．この核の元 $x, y$ で，ある $k$-基底，すなわち
$H^0(E, \mathcal{O}_E(-E)) = H^0(\mathbf{P}^1_k, \mathcal{O}(1))$ の
基底に写るものを選ぶ．すると
$x^d, x^{d - 1}y, \ldots, y^d$ は
$\lim H^0(E_n, \mathcal{O}_n(-dE))$ の元であり，
$H^0(E, \mathcal{O}_E(-dE)) = H^0(\mathbf{P}^1_k, \mathcal{O}(d))$ の
基底に写る．このようにして，$A$ は核
$$
I_n = \Ker(A \longrightarrow H^0(E_n, \mathcal{O}_n))
$$
が定める線形位相に関して分離かつ完備であることがわかる．
$x, y \in I_1$，$I_d I_{d'} \subset I_{d + d'}$ であり，
$I_d/I_{d + 1}$ は自由 $k$-加群であり，その基底は
$x^d, x^{d - 1}y, \ldots, y^d$ である．$I_d = (x, y)^d$ を示そう．実際，
$z_e \in I_e$ かつ $e \geq d$ ならば
$$
z_e = a_{e, 0} x^d + a_{e, 1} x^{d - 1}y + \ldots + a_{e, d}y^d + z_{e + 1}
$$
と書ける．ここで $a_{e, j} \in (x, y)^{e - d}$ かつ
$z_{e + 1} \in I_{e + 1}$ である．これは上の
$I_d/I_{d + 1}$ の記述による．したがって
$z = z_d \in I_d$ から始め，帰納的に
$$
z = \sum\nolimits_{e \geq d} \sum\nolimits_j a_{e, j} x^{d - j} y^j
$$
を得る．ここである $a_{e, j} \in (x, y)^{e - d}$ を用いた．
すると $a_j = \sum_{e \geq d} a_{e, j}$ は存在する
（完備性と $a_{e, j} \in I_{e - d}$ による）．
% P07-ERR-0441: frozen final formula leaves e free instead of using the newly defined a_j; formula retained unchanged.
そして $z = \sum a_{e, j} x^{d - j} y^j$ を得る．
したがって $I_d = (x, y)^d$ である．ゆえに $A$ は
$(x, y)$-進完備である．Algebra, 補題
\ref{algebra-lemma-completion-Noetherian} により $A$ は Noether である．
また，次元が $2$ であることは，$(x, y)^d/(x, y)^{d + 1}$ の記述と
Algebra, 命題 \ref{algebra-proposition-dimension} により明らかである．
極大イデアルは二つの元で生成されるので正則である．
\end{proof}

\begin{lemma}
\label{resolve-lemma-contraction}
$X$ を Noether 概形とし，$E \subset X$ を第一種例外曲線とする．
次を満たす射 $f : X \to Y$ が存在すると仮定する．
\begin{enumerate}
\item $Y$ は Noether である，
\item $f$ は固有である，
\item $f$ は $E$ を一点 $y$（$Y$ の点）に写す，
\item $f$ は $E$ に属さない各点で準有限である．
\end{enumerate}
このとき $E$ の収縮が存在し，それは $f$ の Stein 分解である．
\end{lemma}

\begin{proof}
More on Morphisms, 定理
\ref{more-morphisms-theorem-stein-factorization-Noetherian}
を適用して Stein 分解 $X \to X' \to Y$ を得る．
このとき $X \to X'$ は補題のすべての仮定を満たす
（いくつかの詳細は省略する）．
したがって $Y$ を $X'$ で置き換えることにより，さらに
$f_*\mathcal{O}_X = \mathcal{O}_Y$ であり，$f$ のファイバーが
幾何学的に連結であると仮定してよい．

\medskip\noindent
$f_*\mathcal{O}_X = \mathcal{O}_Y$ であり，$f$ のファイバーが
幾何学的に連結であると仮定する．$y \in Y$ は閉点であることに注意する．
実際，$f$ は閉写像であり $E$ は閉である．制限
$f^{-1}(Y \setminus \{y\}) \to Y \setminus \{y\}$，すなわち $f$ の制限は有限射である
（More on Morphisms, 補題 \ref{more-morphisms-lemma-characterize-finite}）．
% P07-ERR-0442: frozen source repeats "since"; the two reasons are stated coordinately here.
$f_*\mathcal{O}_X = \mathcal{O}_Y$ であり，有限射はアフィンなので，
この制限は同型である．$\mathcal{O}_{Y, y}$ が次元 $2$ の正則環であることを
示すため，Cohomology of Schemes, 補題
\ref{coherent-lemma-formal-functions-stalk} の同型
$$
\mathcal{O}_{Y, y}^\wedge \longrightarrow
\lim H^0(X \times_Y \Spec(\mathcal{O}_{Y, y}/\mathfrak m_y^n), \mathcal{O})
$$
を考える．補題 \ref{resolve-lemma-exceptional-first-kind-local} と同様に
$E_n = nE$ とおく．このとき
$$
E_n \subset X \times_Y \Spec(\mathcal{O}_{Y, y}/\mathfrak m_y^n)
$$
である．なぜなら $E \subset X_y = X \times_Y \Spec(\kappa(y))$ だからである．
他方，集合論的に $E = f^{-1}(\{y\})$ である．これは $f$ のファイバーが
幾何学的に連結だからである．したがって概形論的ファイバー $X_y$ は，
概形論的に $E_n$ に，ある $n > 0$ に対して含まれる．実際，
連接 $\mathcal{O}_X$-加群 $\mathcal{F} = \mathcal{O}_{X_y}$ と
イデアル層 $\mathcal{I}$，すなわち $E$ のイデアル層に Cohomology of Schemes, 補題
\ref{coherent-lemma-power-ideal-kills-sheaf} を適用し，
$\mathcal{I}^n$ が $E_n$ のイデアル層であることを用いればよい．これにより
$$
X \times_Y \Spec(\mathcal{O}_{Y, y}/\mathfrak m_y^m) \subset E_{nm}
$$
が従う．したがって上に表示した逆極限は
$\lim H^0(E_n, \mathcal{O}_n)$ に等しく，これは補題
\ref{resolve-lemma-exceptional-first-kind-local} により二次元正則局所環である．
ゆえに $\mathcal{O}_{Y, y}$ も二次元正則局所環である．これはその完備化が
そうであることによる
（More on Algebra, 補題 \ref{more-algebra-lemma-completion-regular} および
\ref{more-algebra-lemma-completion-dimension}）．

\medskip\noindent
$f : X \to Y$ がブローアップ $b : Y' \to Y$，すなわち
$Y$ の $y$ におけるブローアップであることが残っている．読者には自身の証明を見つけることを勧める．
まず，補題 \ref{resolve-lemma-exceptional-first-kind-local} から，概形論的にも
$X_y = E$ であることが従う．$E$ のイデアル層は可逆なので，
$f^{-1}\mathfrak m_y \cdot \mathcal{O}_X$ は可逆である．
したがって分解
$$
X \to Y' \to Y
$$
を射 $f$ のものとして得る．これは Divisors, 補題
\ref{divisors-lemma-universal-property-blowing-up} のブローアップの普遍性による．
% P07-ERR-0443: E' itself is the exceptional fibre; the frozen malformed "of" is not rendered.
補題 \ref{resolve-lemma-blowup} により，例外ファイバー $E' \subset Y'$ は
第一種例外曲線であることを思い出そう．誘導される射を $g : E \to E'$ とする．
% P07-ERR-0444: frozen source invokes undefined a,b; symbols are retained without inventing definitions.
$E'$ と $E$ のいずれについても，余法層は $a$ と $b$（の引き戻し）で
生成されるので，標準写像
$g^*\mathcal{C}_{E'/Y'} \to \mathcal{C}_{E/X}$
(Morphisms, 補題 \ref{morphisms-lemma-conormal-functorial})
は全射である．両者は可逆なので，この写像は同型である．
$\mathcal{C}_{E/X}$ は正の次数をもつため，$g$ は定値射ではありえない．
したがって $g$ のファイバーは有限であり，ゆえに $g$ は有限射である
（上と同じ参照）．しかし $Y'$ は $E'$ のすべての点で正則，したがって正規であり，
$X \to Y'$ は双有理で $E'$ の外で同型なので，Varieties, 補題
\ref{varieties-lemma-modification-normal-iso-over-codimension-1}.
$X \to Y'$ は同型である．
\end{proof}

\begin{lemma}
\label{resolve-lemma-pic-blowup}
$b : X \to X'$ を第一種例外曲線 $E \subset X$ の収縮とする．
このとき短完全列
$$
0 \to \Pic(X') \to \Pic(X) \to \mathbf{Z} \to 0
$$
が存在する．ここで第一の写像は $b$ による引き戻しであり，第二の写像は
$\mathcal{L}$ を，$\mathcal{L}$ の例外曲線 $E$ 上での次数へ送る．
写像 $n \mapsto \mathcal{O}_X(-nE)$ はこの列の分裂を与える．
\end{lemma}

\begin{proof}
$E = \mathbf{P}^1_k$ なので，Divisors, 補題
\ref{divisors-lemma-Pic-projective-space-UFD}.
により $E$ の Picard 群は $\mathbf{Z}$ である．したがって最後の写像を
$\mathcal{L} \mapsto \mathcal{L}|_E$ とみなせる．第一種例外曲線の定義により，
$\mathcal{O}_X(E)$ の $E$ への制限の次数は $-1$ である．これらを合わせると，
$\Pic(X') \to \Pic(X)$ が単射であり，その像が $\mathcal{O}_E$ へと
$E$ 上で制限される可逆層全体であることを示せば十分である．

\medskip\noindent
可逆 $\mathcal{O}_{X'}$-加群 $\mathcal{L}'$ が与えられたとき，写像
$\mathcal{L}' \to b_*b^*\mathcal{L}'$ は同型であると主張する．
これは像点 $x \in X'$，すなわち $E$ の像を除けば明らかである．
$x$ における茎で同型であることを確かめるには，$X'$ を $x$ の開近傍で
置き換え，$\mathcal{L}'$ が $\mathcal{O}_{X'}$ であると仮定してよい．
すると写像 $\mathcal{O}_{X'} \to b_*\mathcal{O}_X$ が同型であることを
示せばよい．これは補題 \ref{resolve-lemma-cohomology-of-blowup} の (4) から従う．

\medskip\noindent
$\mathcal{L}$ を可逆 $\mathcal{O}_X$-加群で
$\mathcal{L}|_E = \mathcal{O}_E$ を満たすものとする．このとき
(1) $b_*\mathcal{L}$ は可逆であり，
(2) $b^*b_*\mathcal{L} \to \mathcal{L}$ は同型であると主張する．
(1) と (2) は $X' \setminus \{x\}$ 上では明らかである．したがって
$\Spec(\mathcal{O}_{X', x})$ への基底変換後に (1) と (2) を示せば十分である．
$b_*$ の計算は平坦基底変換と可換であり
（Cohomology of Schemes, 補題 \ref{coherent-lemma-flat-base-change-cohomology}），
$b^*$ および随伴写像の形成についても同様である．ところが $X'$ が
正則局所環のスペクトルならば，補題 \ref{resolve-lemma-blowup-pic} の Picard 群の
記述により $\mathcal{L}$ は自明である．これで主張が示された．

\medskip\noindent
前二段落で示した主張を合わせると，写像
$\mathcal{L} \mapsto b_*\mathcal{L}$ は写像
$$
\Pic(X') \longrightarrow \Ker(\Pic(X) \to \Pic(E))
$$
の逆写像であることがわかる．これで補題が示された．
\end{proof}

\begin{remark}
\label{resolve-remark-pic-blowup}
$b : X \to X'$ を第一種例外曲線 $E \subset X$ の収縮とする．
補題 \ref{resolve-lemma-pic-blowup} から同一視
$$
\Pic(X) = \Pic(X') \oplus \mathbf{Z}
$$
を得る．ここで $\mathcal{L}$ が対 $(\mathcal{L}', n)$ に対応することと，
$\mathcal{L} = (b^*\mathcal{L}')(-nE)$，すなわち
$\mathcal{L}(nE) = b^*\mathcal{L}'$ であることとは同値である．実際，補題
\ref{resolve-lemma-pic-blowup} の証明は $\mathcal{L}' = b_*\mathcal{L}(nE)$ を示す．
もちろん対応 $\mathcal{L} \mapsto \mathcal{L}'$ は群準同型である．
\end{remark}

\begin{lemma}
\label{resolve-lemma-lift-sections-and-h1}
$X$ を Noether 概形とし，$E \subset X$ を第一種例外曲線とする．
$\mathcal{L}$ を可逆 $\mathcal{O}_X$-加群とする．$n$ を，
$\mathcal{L}|_E$ が次数 $n$ をもつような整数とする．ここではこれを
$\mathbf{P}^1$ 上の可逆加群とみなしている．このとき
\begin{enumerate}
\item $H^1(X, \mathcal{L}) = 0$ かつ $n \geq 0$ ならば，
$H^1(X, \mathcal{L}(iE)) = 0$ が $0 \leq i \leq n + 1$ に対して成り立つ．
\item $n \leq 0$ ならば，
$H^1(X, \mathcal{L}) \subset H^1(X, \mathcal{L}(E))$ である．
\end{enumerate}
\end{lemma}

\begin{proof}
Divisors, 補題 \ref{divisors-lemma-Pic-projective-space-UFD} により
$\mathcal{L}|_E = \mathcal{O}(n)$ であることに注意する．短完全列
$$
0 \to \mathcal{L} \to \mathcal{L}(E) \to \mathcal{L}(E)|_E \to 0,
$$
に付随するコホモロジー長完全列と帰納法を用い，さらに
$H^1(\mathbf{P}^1, \mathcal{O}(d)) = 0$ が $d \geq -1$ に対して，
$H^0(\mathbf{P}^1, \mathcal{O}(d)) = 0$ が $d \leq -1$ に対して成り立つことを
用いればよい．いくつかの詳細は省略する．
\end{proof}

\begin{lemma}
\label{resolve-lemma-contract-ample}
$S = \Spec(R)$ をアフィン Noether 概形とする．$X \to S$ を固有射とし，
$\mathcal{L}$ を $X$ 上の豊富な可逆層とする．$E \subset X$ を
第一種例外曲線とする．このとき
\begin{enumerate}
\item 収縮 $b : X \to X'$ が $E$ に対して存在する，
\item $X'$ は $S$ 上固有である，
\item 可逆 $\mathcal{O}_{X'}$-加群 $\mathcal{L}'$ は豊富であり，ここで
$\mathcal{L}'$ は注意 \ref{resolve-remark-pic-blowup} のものである．
\end{enumerate}
\end{lemma}

\begin{proof}
補題 \ref{resolve-lemma-lift-sections-and-h1} と同様に，$n$ を
$\mathcal{L}|_E$ の次数とする．$n > 0$ であることに注意する．これは
$\mathcal{L}$ が $E$ 上で豊富だからである
（Varieties, 補題 \ref{varieties-lemma-ample-curve} および
Properties, 補題 \ref{properties-lemma-ample-on-closed}）．
$\mathcal{L}$ をある冪で置き換えることにより，
$H^i(X, \mathcal{L}^{\otimes e}) = 0$ がすべての $i > 0$ と $e > 0$ に対して
成り立つと仮定してよい．Cohomology of Schemes,
補題 \ref{coherent-lemma-vanshing-gives-ample} を参照せよ．
最後に $\mathcal{L}$ をさらに別の冪で置き換えることにより，大域切断
$t_0, \ldots, t_n$ が $\mathcal{L}$ に存在し，それらが閉埋め込み
$\psi : X \to \mathbf{P}^n_S$ を定めると仮定してよい．Morphisms, 補題
\ref{morphisms-lemma-finite-type-over-affine-ample-very-ample}.

\medskip\noindent
$\mathcal{M} = \mathcal{L}(nE)$ とおく．すると
$\mathcal{M}|_E \cong \mathcal{O}_E$ である．短完全列
$$
0 \to \mathcal{M}(-E) \to \mathcal{M} \to \mathcal{O}_E \to 0
$$
があり，$H^1(X, \mathcal{M}(-E))$ は零なので
（補題 \ref{resolve-lemma-lift-sections-and-h1} および $n > 0$ による），
切断 $s_{n + 1}$ を $\mathcal{M}$ から，$\mathcal{M}|_E$ を生成するよう選べる．
さらに，$s_0, \ldots, s_n$ で $\mathcal{M}$ の切断であって，
上で選んだ $t_0, \ldots, t_n$，すなわち $\mathcal{L}$ の切断から，
$\mathcal{L} \subset \mathcal{L}(nE) = \mathcal{M}$ を通じて得られるものを表す．
% P07-ERR-0445: frozen source says "Combined the sections"; intended collective sense is rendered.
切断 $s_0, \ldots, s_n, s_{n + 1}$ は合わせて $\mathcal{M}$ を
$X$ の各点で生成し，したがって射
$$
\varphi : X \longrightarrow \mathbf{P}^{n + 1}_S
$$
を $S$ 上に定める．Constructions, 補題
\ref{constructions-lemma-projective-space} を参照せよ．

\medskip\noindent
以下で補題 \ref{resolve-lemma-contraction} の条件を確認する．これが済めば，
Stein 分解 $X \to X' \to \mathbf{P}^{n + 1}_S$，すなわち $\varphi$ のものが
所望の収縮であることがわかり，(1) が示される．さらに射
$X' \to \mathbf{P}^{n + 1}_S$ は有限なので，$X'$ は $S$ 上固有である
（Morphisms, 補題 \ref{morphisms-lemma-finite-proper} および
\ref{morphisms-lemma-composition-proper}）．これで (2) が示された．
$X'$ は豊富な可逆層をもつことに注意する．実際，
$\mathcal{M}'$，すなわち $\mathcal{O}_{\mathbf{P}^{n + 1}_S}(1)$ の引き戻しは
Morphisms, 補題 \ref{morphisms-lemma-pullback-ample-tensor-relatively-ample}
により豊富である．また $\mathcal{M}'$ の引き戻しは $\mathcal{M}$ であり，
これは $X$ 上での引き戻しである
（Constructions, 補題 \ref{constructions-lemma-projective-space}）．
最後に $\mathcal{M} = \mathcal{L}(nE)$ である．上の議論では元の
$\mathcal{L}$ を正の冪で置き換えたので，Properties, 補題
\ref{properties-lemma-ample-power-ample} により，補題の (3) で述べた
可逆 $\mathcal{O}_{X'}$-加群 $\mathcal{L}'$ は $X'$ 上で豊富である．

\medskip\noindent
容易な観察として，$\mathbf{P}^{n + 1}_S$ は Noether であり，
$\varphi$ は固有である．詳細は省略する．

\medskip\noindent
次に，$U = X \setminus E$ の任意の点は，開部分概形 $W$，すなわち
$\mathbf{P}^{n + 1}_S$ において最初の $n + 1$ 個の斉次座標の
いずれかが零でない部分に写る．
他方，$E$ の任意の点は最初の $n + 1$ 個の斉次座標がすべて零である点に写り，
特に $W$ の補集合に写る．さらに，分解
$$
U = \varphi^{-1}(W) \xrightarrow{\varphi|_U}
W \xrightarrow{\text{pr}} \mathbf{P}^n_S
$$
が $\psi|_U$ に対してあることは明らかである．ここで $\text{pr}$ は最初の $n + 1$ 個の座標を
用いる射影であり，$\psi : X \to \mathbf{P}^n_S$ は上で選んだ埋め込みである．
したがって $\varphi|_U : U \to W$ は準有限である．

\medskip\noindent
最後に写像 $\varphi|_E : E \to \mathbf{P}^{n + 1}_S$ を考える．
任意の点 $x \in E$ に対し，像 $\varphi(x)$ の最初の $n + 1$ 個の座標は
零である．すなわち射 $\varphi|_E$ は閉部分概形
$\mathbf{P}^0_S \cong S$ を経由する．Schemes, 補題
\ref{schemes-lemma-morphism-into-affine} により，射
$E \to S = \Spec(R)$ は
$E \to \Spec(H^0(E, \mathcal{O}_E)) \to \Spec(R)$ と分解する．
仮定により $H^0(E, \mathcal{O}_E)$ は体なので，$E$ は
$S \subset \mathbf{P}^{n + 1}_S$ の一点に写る．これで証明が完了する．
\end{proof}

\begin{lemma}
\label{resolve-lemma-contract-when-quasi-projective}
$S$ を Noether 概形とし，$f : X \to S$ を有限型の射とする．
$E \subset X$ を $f$ の一つのファイバーに含まれる第一種例外曲線とする．
\begin{enumerate}
\item $X$ が $S$ 上射影的ならば，収縮 $X \to X'$ が $E$ に対して存在し，
$X'$ は $S$ 上射影的である．
\item $X$ が $S$ 上準射影的ならば，収縮 $X \to X'$ が $E$ に対して存在し，
$X'$ は $S$ 上準射影的である．
\end{enumerate}
\end{lemma}

\begin{proof}
いずれの場合も，豊富な可逆加群と（準）射影的射に関する標準的な結果を用いれば，
補題 \ref{resolve-lemma-contract-ample} から従う．

\medskip\noindent
(1) の証明．$f$ が射影的であるとは，$f$ が固有であり，$f$-豊富な
可逆加群 $\mathcal{L}$ が存在することを意味する．Morphisms, 補題
\ref{morphisms-lemma-projective-is-quasi-projective-proper} および定義
\ref{morphisms-definition-quasi-projective} を参照せよ．$U \subset S$ を
$E$ の像を含むアフィン開とする．補題 \ref{resolve-lemma-contract-ample} により，
収縮 $c : f^{-1}(U) \to V'$ が $E$ に対して存在し，豊富な可逆加群
$\mathcal{N}'$ が $V'$ 上に存在する．その $f^{-1}(U)$ への引き戻しは
$\mathcal{L}(nE)|_{f^{-1}(U)}$ に等しい．$v \in V'$ を，$c$ が
$v$ のブローアップとなるような閉点とする．このとき $V'$ と $X \setminus E$ を
$f^{-1}(U) \setminus E = V' \setminus \{v\}$
に沿って貼り合わせ，概形 $X'$ を $S$ 上に得る．射 $c$ と
$\text{id}_{X \setminus E}$ は射 $b : X \to X'$ に貼り合わさり，これは
$E$ の収縮である．$U$ の $X'$ における逆像は $U$ 上固有である．他方，
$X' \to S$ を $v$ の $S$ における像の補集合に制限すると，その開集合への
$X \to S$ の制限と同型である．したがって $X' \to S$ は固有である
（固有性は基底上局所的である；Morphisms, 補題
\ref{morphisms-lemma-proper-local-on-the-base}）．最後に，$\mathcal{N}'$ と
$\mathcal{L}|_{X \setminus E}$ の
$f^{-1}(U) \setminus E = V' \setminus \{v\}$ への制限は同型な可逆加群なので，
可逆加群 $\mathcal{L}'$ として $X'$ 上に貼り合わさる．$\mathcal{L}'$ の
$U$ の逆像への制限は，$X'$ において豊富である．これは $\mathcal{N}'$ についてそうだからである．
$S$ のアフィン開で $v$ の像を避けるものの上でも，$\mathcal{L}$ についてそうなので同じことが成り立つ．
ゆえに Morphisms, 補題 \ref{morphisms-lemma-characterize-relatively-ample}
により $\mathcal{L}'$ は $(X' \to S)$-相対的に豊富である．これで (1) が示された．

\medskip\noindent
(2) の証明．Morphisms, 補題
\ref{morphisms-lemma-quasi-projective-open-projective} により，$X$ を
$\overline{X}$，すなわち $S$ 上射影的な概形の開部分概形として書ける．(1) により収縮
$b : \overline{X} \to \overline{X}'$ が存在し，$\overline{X}'$ は $S$ 上射影的である．
% P07-ERR-0446: frozen source puts the image in the wrong ambient scheme; inclusion formula retained unchanged.
そこで $X' \subset \overline{X}$ を $X \to \overline{X}'$ の像とする．
$b$ は $E$ の外で同型なので，これは開である．すると $X \to X'$ は
所望の収縮である．$X'$ は $S$ 上準射影的であることに注意する．これは
(1) の証明の構成により $S$-相対的に豊富な可逆加群をもつからである．
\end{proof}

\begin{lemma}
\label{resolve-lemma-regular-dim-2-quasi-projective}
$S$ を Noether 概形とする．$f : X \to S$ を有限型の分離射とし，
$X$ は次元 $2$ の正則概形であるとする．このとき $X$ は $S$ 上準射影的である．
\end{lemma}

\begin{proof}
Chow の補題
（Cohomology of Schemes, 補題 \ref{coherent-lemma-chow-Noetherian}）により，
固有射 $\pi : X' \to X$ で，稠密開 $U \subset X$ 上で同型となり，かつ
$X' \to S$ が H-準射影的となるものが存在する．補題
\ref{resolve-lemma-extend-rational-map-blowing-up} により，閉点におけるブローアップの列
$$
X_n \to \ldots \to X_1 \to X_0 = X
$$
と，$S$-射 $X_n \to X'$ で有理写像 $U \to X'$ を延長するものが存在する．
$X_n \to X$ は射影的である．これは Divisors, 補題
\ref{divisors-lemma-blowing-up-projective} および Morphisms, 補題
\ref{morphisms-lemma-composition-projective} による．したがって $X_n \to X'$ は
Morphisms, 補題 \ref{morphisms-lemma-projective-permanence} により射影的である．
ゆえに $X_n \to S$ は Morphisms, 補題
\ref{morphisms-lemma-composition-quasi-projective} により準射影的である
（射影的射が準射影的であることについては Morphisms, 補題
\ref{morphisms-lemma-projective-quasi-projective} を参照せよ）．補題
\ref{resolve-lemma-contract-when-quasi-projective} と収縮の一意性
（補題 \ref{resolve-lemma-contraction-unique}）により，所望のとおり
$X_{n - 1}, \ldots, X_0 = X$ は $S$ 上準射影的である．
\end{proof}

\begin{lemma}
\label{resolve-lemma-regular-dim-2-projective}
$S$ を Noether 概形とする．$f : X \to S$ を固有射とし，
$X$ は次元 $2$ の正則概形であるとする．このとき $X$ は $S$ 上射影的である．
\end{lemma}

\begin{proof}
これは補題 \ref{resolve-lemma-regular-dim-2-quasi-projective} と Morphisms, 補題
\ref{morphisms-lemma-projective-is-quasi-projective-proper} から従う．
\end{proof}






% P07-ERR-0447: frozen English heading omits "of"; Japanese expresses the intended genitive relation.
\section{双有理写像の分解}
\label{resolve-section-factorizing}

\noindent
非特異曲面の間の固有双有理射は，二次変換の列として与えられる．

\begin{lemma}
\label{resolve-lemma-proper-birational-regular-surfaces}
$f : X \to Y$ を，次元 $2$ の正則な整 Noether 概形の間の固有双有理射とする．
このとき $f$ は閉点におけるブローアップの列である．
\end{lemma}

\begin{proof}
$V \subset Y$ を，$f$ が同型となる最大の開部分とする．すると $V$ は
余次元 $1$ の，$V$ の点をすべて含む
（Varieties, 補題 \ref{varieties-lemma-modification-normal-iso-over-codimension-1}）．
$y \in Y$ を $V$ に含まれない閉点とする．$f$ がブローアップ
$b : Y' \to Y$，すなわち $Y$ の $y$ におけるものを経由して分解することを
示したい．実際，これが成り立てば，ファイバー $f^{-1}(y)$ の少なくとも一つ
（実はちょうど一つ）の成分が $Y'$ の例外曲線に同型に写る．
% P07-ERR-0448: frozen source says "less that"; intended strict comparison is rendered.
さらに $X \to Y'$ のファイバー内の曲線数は $X \to Y$ のファイバー内の
曲線数より真に少なくなるので，帰納法で結論を得る．いくつかの詳細は省略する．

\medskip\noindent
補題 \ref{resolve-lemma-extend-rational-map-blowing-up} により，
% P07-ERR-0449: frozen source uses the malformed plural "blowing ups"; standard Japanese term is used.
閉点におけるブローアップの列
$$
X' = X_n \to X_{n - 1} \to \ldots \to X_1 \to X_0 = X
$$
で，ファイバー $f^{-1}(y)$ の上にあるものと，射 $X' \to Y'$ であって，図式
$$
\xymatrix{
X' \ar[d]_{f'} \ar[r] & X \ar[d]^f \\
Y' \ar[r] & Y
}
$$
を可換にするものが存在する．射 $X' \to Y'$ が $X$ を経由して分解することを
示したい．そうすれば $n$ に関する帰納法により，$X' \to X$ が $X$ の
閉点 $x \in X$，すなわち $y$ に写る点におけるブローアップである場合に帰着できる．

\medskip\noindent
$E \subset X'$ をブローアップ $X' \to X$ の例外ファイバーとする．
$E$ が $Y'$ の一点に写るならば，補題 \ref{resolve-lemma-factor-through-contraction}
により所望の分解を得る．そうでなければ矛盾が生じることを示そう．実際，
$f'(E)$ が一点でなければ，$E' = f'(E)$ は $Y'$ の例外曲線でなければならない．
図式は次のとおりである：
$$
\xymatrix{
E \ar[r] \ar[d]_g & X' \ar[d]_{f'} \ar[r] & X \ar[d]^f \\
E' \ar[r] & Y' \ar[r] & Y
}
$$
前と同様に論じると，$f'$ は $E'$ の生成点のある開近傍で同型である．
したがって $g : E \to E'$ は有限双有理射である．すると $g$ の逆写像
（有理写像）は Morphisms, 補題 \ref{morphisms-lemma-extend-across} により
至る所で定義され，$g$ は同型である．写像
$$
g^*\mathcal{C}_{E'/Y'} \longrightarrow \mathcal{C}_{E/X'}
$$
を考える．これは Morphisms, 補題 \ref{morphisms-lemma-conormal-functorial}
の写像である．その始域と終域は次数 $1$ の可逆加群であり，これは
$E = E' = \mathbf{P}^1_\kappa$ 上のものである．またこの写像は零でない
% P07-ERR-0450: frozen source says "in the generic point"; Japanese uses the standard pointwise phrasing.
（$f'$ は $E$ の生成点で同型である）ので，これは同型である．
Morphisms, 補題 \ref{morphisms-lemma-two-immersions} により
$\Omega_{X'/Y'}|_E = 0$ と結論する．これは $f'$ が $E$ の各点で不分岐であることを意味する
（Morphisms, 補題 \ref{morphisms-lemma-unramified-at-point}）．
したがって $f'$ は $E$ の各点で準有限である
（Morphisms, 補題 \ref{morphisms-lemma-unramified-quasi-finite}）．
ゆえに最大の開 $V' \subset Y'$，すなわち $f'$ が同型となるものは，Varieties, 補題
\ref{varieties-lemma-modification-normal-iso-over-codimension-1} により $E'$ を含む．
% P07-ERR-0451: frozen source names E' in X' although E' lies in Y'; token retained unchanged.
これはさらに，$y$ の $X'$ における逆像が $E'$ であることを意味する．
したがって $y$ の $X$ における逆像は $x$ である．ゆえに $x \in X$ は，
Varieties, 補題 \ref{varieties-lemma-modification-normal-iso-over-codimension-1}
により $f$ が同型となる最大の開部分に属する．これは $y$ がこの開部分に
属さないと仮定したことに矛盾する．
\end{proof}

\begin{lemma}
\label{resolve-lemma-birational-regular-surfaces}
$S$ を Noether 概形とする．$X$ と $Y$ を，$S$ 上固有な整概形で，
次元 $2$ の正則概形であるものとする．このとき $X$ と $Y$ が
$S$-双有理であることと，$S$-射の図式
$$
X = X_0 \leftarrow X_1 \leftarrow \ldots \leftarrow X_n = Y_m
\to \ldots \to Y_1 \to Y_0 = Y
$$
で，各射が閉点におけるブローアップであるものが存在することとは同値である．
\end{lemma}

\begin{proof}
$U \subset X$ を開とし，$f : U \to Y$ を与えられた $S$-有理写像とする
（これは $S$-有理写像として可逆である）．補題
\ref{resolve-lemma-extend-rational-map-blowing-up} により，$f$ を
$X_n \to \ldots \to X_1 \to X_0 = X$ と $f_n : X_n \to Y$ に分解できる．
$X_n$ は $S$ 上固有であり，$Y$ は $S$ 上分離的なので，射 $f_n$ は固有である．
明らかに $f_n$ は双有理である．したがって補題
\ref{resolve-lemma-proper-birational-regular-surfaces} により，$f_n$ は収縮の合成である．
逆向きの証明は省略する．
\end{proof}
